% Source: Lee, "Introduction to Riemannian Manifolds" (2nd ed.), Chapter 2 "Riemannian Metrics", PDF pages 23-68 of the cited source.
% Transcribed from the cited source and organized according to
% leanblueprint conventions (semantic \label, bracketed titles, \uses dependency edges, \dcref to Lee's own numbering) below.
\section{Definitions}

For $v=(v^1,\ldots,v^n)$ and $w=(w^1,\ldots,w^n)$ in $\mathbb R^n$, the Euclidean dot product is
\[
v\mathbin\cdot w=\sum_{i=1}^n v^i w^i.
\]
An inner product on a real vector space $V$ is a map $\langle\ ,\ \rangle:V\times V\to\mathbb R$ that is symmetric, bilinear, and satisfies $\langle v,v\rangle\ge0$ with equality only for $v=0$. A vector space with a specified inner product is an inner product space. Define $|v|=\langle v,v\rangle^{1/2}$.

\begin{lemma}[Polarization Identity]
\label{lem:polarization-identity}
\lean{LeeLib.Ch02.real_inner_eq_inner_add_add_self_sub_inner_sub_sub_self_div_four}
\leanok


\dcref{ch2:2.1}
For every inner product and $v,w\in V$,
\[
\langle v,w\rangle=\frac14\bigl(\langle v+w,v+w\rangle-\langle v-w,v-w\rangle\bigr).
\]
\end{lemma}
\begin{proof}
\leanok
\sketch Expanding both terms by bilinearity and symmetry gives $\langle v+w,v+w\rangle=\langle v,v\rangle+2\langle v,w\rangle+\langle w,w\rangle$ and $\langle v-w,v-w\rangle=\langle v,v\rangle-2\langle v,w\rangle+\langle w,w\rangle$; subtraction and division by $4$ yield the formula.
\end{proof}

For nonzero $v,w$, their angle is the unique $\theta\in[0,\pi]$ with
\[
\cos\theta=\frac{\langle v,w\rangle}{|v||w|}.
\]
Vectors are orthogonal when their inner product is zero. For a subspace $S\subseteq V$, its orthogonal complement $S^\perp$ is the subspace of vectors orthogonal to every element of $S$. A family $(v_i)$ is orthonormal when $\langle v_i,v_j\rangle=\delta_{ij}$.

\begin{proposition}[Gram-Schmidt Algorithm]
\label{prop:gram-schmidt}
\lean{LeeLib.Ch02.exists_orthonormalBasis_span_image_Iic_eq}
\leanok


\dcref{ch2:2.3}
If $(v_1,\ldots,v_n)$ is an ordered basis of an $n$-dimensional inner product space $V$, there is an orthonormal ordered basis $(b_1,\ldots,b_n)$ such that
\[
\operatorname{span}(b_1,\ldots,b_k)=\operatorname{span}(v_1,\ldots,v_k),\qquad 1\le k\le n.
\]
\end{proposition}
\begin{proof}
\leanok
\sketch Set $b_1=v_1/|v_1|$ and, for $2\le j\le n$,
\[
b_j=\frac{v_j-\sum_{i<j}\langle v_j,b_i\rangle b_i}{\left|v_j-\sum_{i<j}\langle v_j,b_i\rangle b_i\right|}.
\]
The numerator is nonzero because $v_j$ is outside the preceding span; the usual Gram--Schmidt computation gives orthonormality and the stated span flag.
\end{proof}

An invertible linear map $F:V\to W$ between inner product spaces is a linear isometry when $\langle F(v),F(w)\rangle_W=\langle v,w\rangle_V$.

\begin{proposition}[Classification of Inner Product Spaces]
\label{prop:inner-product-spaces-linearly-isometric}
\lean{LeeLib.Ch02.nonempty_linearIsometryEquiv_of_finrank_eq}
\uses{prop:gram-schmidt}
\leanok



Finite-dimensional inner product spaces of equal dimension are linearly isometric.
\end{proposition}
\begin{proof}
\leanok
\sketch Choose orthonormal bases $(v_i)$ and $(w_i)$ by Proposition \ref{prop:gram-schmidt}; the unique linear map $F(v_i)=w_i$ is an isomorphism and, by orthonormality, preserves inner products.
\end{proof}

\section{Riemannian Metrics}

\begin{definition}[Riemannian Metric]
\label{def:riemannian-metric}
\lean{LeeLib.Ch02.RiemannianMetric}
\leanok


A Riemannian metric on a smooth manifold $M$ is a smooth covariant $2$-tensor $g$ whose value $g_p$ is an inner product on $T_pM$ for every $p$. A Riemannian manifold is a pair $(M,g)$.
\end{definition}

\begin{lemma}[Convexity of the set of inner products]
\label{lem:convex-inner-products}
\lean{LeeLib.Ch02.convex_symm_posDef}
\leanok


For $p\in M$, if $q_1,q_2$ are inner products on $T_pM$ and $a,b\ge0$ with $a+b=1$, then $a q_1+b q_2$ is an inner product.
\end{lemma}
\begin{proof}
\leanok
\sketch Symmetry and bilinearity are preserved by linear combination. For $v\ne0$, at least one coefficient is positive and both $q_i(v,v)$ are positive, so $(a q_1+b q_2)(v,v)>0$.
\end{proof}

\begin{lemma}[Positive definite forms are coercive]
\label{lem:posdef-coercive}
\lean{LeeLib.Ch02.isVonNBounded_of_posDef}
\leanok


Let $V$ be finite-dimensional normed and let $q$ be bilinear with $q(v,v)>0$ for $v\ne0$. Then $\{v:q(v,v)<1\}$ is bounded; equivalently, some $c>0$ satisfies $q(v,v)\ge c|v|^2$ for all $v$.
\end{lemma}
\begin{proof}
\leanok
\sketch If $V\ne\{0\}$, compactness of the unit sphere gives a minimum $c>0$ of $q(v,v)$ there. Scaling $u=v/|v|$ yields $q(v,v)\ge c|v|^2$; hence $q(v,v)<1$ implies $|v|<c^{-1/2}$. The zero space is vacuous.
\end{proof}

\begin{proposition}[Every smooth manifold admits a Riemannian metric]
\label{prop:existence-riemannian-metric}
\lean{LeeLib.Ch02.exists_riemannianMetric}
\uses{def:riemannian-metric}
\leanok



\dcref{ch2:2.4}
Every smooth manifold admits a Riemannian metric.
\end{proposition}
\begin{proof}
\leanok
\uses{def:riemannian-metric,lem:convex-inner-products,lem:posdef-coercive}
\sketch On each chart domain $U_\alpha$ pull back the Euclidean metric of Definition \ref{ex:euclidean-metric} to $g_\alpha$. For a smooth partition of unity $(\psi_\alpha)$ subordinate to this cover, set
\[
g=\sum_\alpha\psi_\alpha g_\alpha,
\]
extending each summand by zero outside $\operatorname{supp}\psi_\alpha$. Local finiteness gives a smooth tensor. At $p$, the nonzero weights form a finite convex combination of inner products, hence give an inner product by Lemma \ref{lem:convex-inner-products}; Lemma \ref{lem:posdef-coercive} supplies the usual finite-dimensional topology.
\end{proof}

The same definition applies to manifolds with boundary. We use ``metric'' for a Riemannian metric. For $v,w\in T_pM$, write $\langle v,w\rangle_g=g_p(v,w)$ and $|v|_g=\langle v,v\rangle_g^{1/2}$, omitting the subscript when clear.

\begin{definition}[The Euclidean Metric]
\label{ex:euclidean-metric}
\lean{LeeLib.Ch02.euclideanMetric}
\uses{def:riemannian-metric}
\leanok



\dcref{ch2:2.6}
The Euclidean metric $\bar g$ on $\mathbb R^n$ is the pointwise dot product under $T_x\mathbb R^n\cong\mathbb R^n$:
\[
\left\langle\sum_i v^i\partial_i|_x,\sum_j w^j\partial_j|_x\right\rangle_{\bar g}=\sum_{i=1}^n v^i w^i.
\]
More generally, an inner product space carries the constant metric with that inner product.
\end{definition}

\section{Isometries}

\begin{definition}[Metric-Preserving Maps]
\label{def:metric-preserving}
\lean{LeeLib.Ch02.IsMetricPreserving}
\uses{def:riemannian-metric}
\leanok



For Riemannian manifolds with or without boundary, a map $\varphi:(M,g)\to(\widetilde M,\widetilde g)$ preserves the metric when $\varphi^*\widetilde g=g$, equivalently
\[
\widetilde g_{\varphi(p)}(d\varphi_pv,d\varphi_pw)=g_p(v,w)
\quad(p\in M, v,w\in T_pM).
\]
Each $d\varphi_p$ is then a linear isometry onto its image.
\end{definition}

\begin{definition}[Isometry]
\label{def:isometry}
\lean{LeeLib.Ch02.IsIsometry}
\uses{def:metric-preserving, def:riemannian-metric}
\leanok



An isometry is a diffeomorphism that preserves the metric.
\end{definition}

\begin{lemma}[Metric-Preserving Maps Are Immersions]
\label{lem:metric-preserving-immersion}
\lean{LeeLib.Ch02.IsMetricPreserving.injective_mfderiv}
\uses{def:metric-preserving, lem:pullback-immersion-metric}
\leanok



If $\varphi$ preserves the metric, each $d\varphi_p$ is injective.
\end{lemma}
\begin{proof}
\leanok
\uses{lem:pullback-immersion-metric}
\sketch Since $\varphi^*\widetilde g=g$ is positive definite, Lemma \ref{lem:pullback-immersion-metric} implies that every $d\varphi_p$ is injective.
\end{proof}

\begin{lemma}[The Inverse of an Isometry]
\label{lem:isometry-symm}
\lean{LeeLib.Ch02.IsIsometry.symm}
\uses{def:isometry}
\leanok



If $\varphi:(M,g)\to(\widetilde M,\widetilde g)$ is an isometry, then $\varphi^{-1}$ is an isometry in the opposite direction.
\end{lemma}
\begin{proof}
\leanok
\uses{def:isometry,def:metric-preserving}
\sketch For $q=\varphi(p)$, differentiate $\varphi\circ\varphi^{-1}=\operatorname{id}$ to obtain $d\varphi_p\circ d(\varphi^{-1})_q=\operatorname{id}$. Applying $\varphi^*\widetilde g=g$ to $d(\varphi^{-1})_qv,d(\varphi^{-1})_qw$ gives
\[
g_p(d(\varphi^{-1})_qv,d(\varphi^{-1})_qw)=\widetilde g_q(v,w),
\]
so $(\varphi^{-1})^*g=\widetilde g$.
\end{proof}

\begin{definition}[Local Isometry]
\label{def:local-isometry}
\lean{LeeLib.Ch02.IsLocalIsometry}
\uses{def:isometry, def:metric-preserving, def:riemannian-metric}
\leanok



A map between Riemannian manifolds is a local isometry if every point has a neighborhood on which it is an isometry onto an open subset. Equivalently, it is a local diffeomorphism preserving the metric pointwise.
\end{definition}

\begin{lemma}[Metric-Preserving Maps in Equal Dimensions Are Local Diffeomorphisms]
\label{lem:metric-preserving-equidim-local-diffeomorphism}
\lean{LeeLib.Ch02.IsMetricPreserving.isLocalDiffeomorph}
\uses{def:metric-preserving, lem:metric-preserving-immersion, thm:inverse-function-theorem-manifolds}
\leanok



If $M$ and $\widetilde M$ have equal dimension and $\varphi:M\to\widetilde M$ is smooth with $\varphi^*\widetilde g=g$, then $\varphi$ is a local diffeomorphism.
\end{lemma}
\begin{proof}
\leanok
\uses{lem:metric-preserving-immersion,thm:inverse-function-theorem-manifolds}
\sketch Lemma \ref{lem:metric-preserving-immersion} makes $d\varphi_p$ injective; equal finite dimensions make it invertible. The inverse function theorem, Theorem \ref{thm:inverse-function-theorem-manifolds}, then gives a local diffeomorphism at every $p$.
\end{proof}

\begin{proposition}
\label{prop:characterize-local-isometry}
\dcref{ch2:2.7}


\uses{def:local-isometry,def:metric-preserving,def:isometry,lem:metric-preserving-equidim-local-diffeomorphism}
For equal-dimensional Riemannian manifolds, a smooth map is a local isometry
if and only if it preserves the metric.
\end{proposition}
\begin{proof}
\leanok
\uses{def:local-isometry,lem:metric-preserving-equidim-local-diffeomorphism}
\sketch A local isometry preserves the metric on each neighborhood, hence everywhere. Conversely, metric preservation and Lemma \ref{lem:metric-preserving-equidim-local-diffeomorphism} give local diffeomorphisms; their restrictions preserve the metric and are local isometries.
\end{proof}

A Riemannian $n$-manifold is flat when every point has a neighborhood isometric to an open subset of Euclidean $n$-space. Isometries of $(M,g)$ form the group $\operatorname{Iso}(M,g)$ under composition.

\section{Local Representations for Metrics}

In coordinates $(x^1,\ldots,x^n)$ on $U\subseteq M$,
\[
g=g_{ij}\,dx^i\otimes dx^j=g_{ij}\,dx^i dx^j,\tag{2.7}
\]
where $g_{ij}=\langle\partial_i,\partial_j\rangle$ are smooth, symmetric, and nonsingular at every point. In Euclidean coordinates,
\[
\bar g=\sum_i dx^i dx^i=\sum_i(dx^i)^2=\delta_{ij}dx^i dx^j.\tag{2.8}
\]
For a smooth frame $(E_i)$ with dual coframe $(e^i)$,
\[
g=g_{ij}e^i\otimes e^j,\qquad g_{ij}=\langle E_i,E_j\rangle.\tag{2.9}
\]
For smooth vector fields $X=X^iE_i$ and $Y=Y^jE_j$, $\langle X,Y\rangle=g_{ij}X^iY^j$ is smooth; $|X|=\langle X,X\rangle^{1/2}$ is continuous and smooth where $X\ne0$.

\begin{lemma}[Smoothness of the pairing of vector fields]
\label{lem:inner-product-vector-fields-smooth}
\lean{LeeLib.Ch02.RiemannianMetric.contMDiff_innerAt}
\uses{def:riemannian-metric}
\leanok



\dcref{ch2:2.2}
For smooth vector fields $X,Y$ on a Riemannian manifold with or without boundary, $\langle X,Y\rangle:M\to\mathbb R$ is smooth.
\end{lemma}
\begin{proof}
\leanok
\uses{def:riemannian-metric}
\sketch In a smooth local frame, $\langle X,Y\rangle=g_{ij}X^iY^j$, a finite sum of products of smooth functions.
\end{proof}

\begin{lemma}[Continuity of the length of a vector field]
\label{lem:length-vector-field-continuous}
\lean{LeeLib.Ch02.RiemannianMetric.continuous_normAt}
\uses{lem:inner-product-vector-fields-smooth}
\leanok



\dcref{ch2:2.2}
For smooth $X$ on a Riemannian manifold with or without boundary, the function $|X|=\langle X,X\rangle^{1/2}$ is continuous.
\end{lemma}
\begin{proof}
\leanok
\uses{lem:inner-product-vector-fields-smooth}
\sketch The pairing is smooth and nonnegative by Lemma \ref{lem:inner-product-vector-fields-smooth}; compose it with the continuous square root on $[0,\infty)$.
\end{proof}

\begin{lemma}[Smoothness of the length away from the zero set]
\label{lem:length-vector-field-smooth}
\lean{LeeLib.Ch02.RiemannianMetric.contMDiffAt_normAt}
\uses{lem:inner-product-vector-fields-smooth}
\leanok



\dcref{ch2:2.2}
For smooth $X$ on a Riemannian manifold with or without boundary, $|X|$ is smooth at every point where $X$ is nonzero, hence on the open nonzero set.
\end{lemma}
\begin{proof}
\leanok
\uses{lem:inner-product-vector-fields-smooth}
\sketch At $p$ with $X_p\ne0$, positivity gives $\langle X_p,X_p\rangle>0$; the pairing is smooth by Lemma \ref{lem:inner-product-vector-fields-smooth}, and the square root is smooth on $(0,\infty)$.
\end{proof}

A local frame $(E_i)$ is orthonormal when $\langle E_i,E_j\rangle=\delta_{ij}$. With dual coframe $(e^i)$, this is equivalent to $g=\sum_i(e^i)^2$.

\begin{proposition}[Existence of Orthonormal Frames]
\label{prop:existence-orthonormal-frames}
\lean{LeeLib.Ch02.exists_orthonormalFrame}
\uses{prop:gram-schmidt, def:riemannian-metric}
\leanok


\dcref{ch2:2.8}

If $(M,g)$ is a Riemannian $n$-manifold with or without boundary and $(X_i)$ is a smooth local frame on $U\subseteq M$, there is a smooth orthonormal frame $(E_j)$ on $U$ with
\[
\operatorname{span}(E_1|_p,\ldots,E_k|_p)=\operatorname{span}(X_1|_p,\ldots,X_k|_p)
\]
for every $p$ and $k$. In particular, every point has a neighborhood with a smooth orthonormal frame.
\end{proposition}
\begin{proof}
\leanok
\uses{prop:gram-schmidt}
\sketch Apply Proposition \ref{prop:gram-schmidt} fibrewise. Writing
\[
u_j=X_j-\sum_{i<j}\frac{\langle u_i,X_j\rangle}{\langle u_i,u_i\rangle}u_i,
\qquad E_j=\frac{u_j}{|u_j|},
\]
linear independence of the original frame gives $u_j\ne0$ and hence nonvanishing denominators. Inductively the pairings and quotients are smooth, so the $E_j$ are smooth; fibrewise Gram--Schmidt gives orthonormality and the span flag.
\end{proof}

\begin{remark}
\label{rem:coordinate-frame-warning}
\uses{prop:existence-orthonormal-frames}

An orthonormal frame need not be a coordinate frame. A neighborhood admits orthonormal coordinate vector fields exactly when the metric is flat (locally Euclidean).
\end{remark}

For a Riemannian manifold, define
\[
UTM=\{(p,v)\in TM:|v|_g=1\}.\tag{2.10}
\]
\begin{proposition}[Properties of the Unit Tangent Bundle]
\label{prop:unit-tangent-bundle-properties}
\dcref{ch2:2.9}
If $(M,g)$ is a Riemannian manifold with or without boundary, then $UTM$ is a smooth properly embedded codimension-$1$ submanifold with boundary in $TM$, with $\partial(UTM)=\pi^{-1}(\partial M)$. It is connected exactly when $M$ is connected, and compact exactly when $M$ is compact.
\end{proposition}

\section{Methods for Constructing Riemannian Metrics}

\subsection*{Riemannian Submanifolds}

\begin{lemma}
\label{lem:pullback-immersion-metric}
\lean{LeeLib.Ch02.pullbackForm_posDef_iff_immersion}
\uses{def:riemannian-metric}
\leanok



\dcref{ch2:2.11}
Let $(\widetilde M,\widetilde g)$ be Riemannian with or without boundary, let $M$ be smooth with or without boundary, and let $F:M\to\widetilde M$ be smooth. Then $F^*\widetilde g$ is a Riemannian metric on $M$ if and only if $F$ is an immersion.
\end{lemma}
\begin{proof}
\leanok
\uses{def:riemannian-metric}
\sketch The pullback is smooth and symmetric, with
\[
(F^*\widetilde g)_p(v,w)=\widetilde g_{F(p)}(dF_pv,dF_pw).
\]
If $dF_p$ is injective, positivity of $\widetilde g$ gives positivity of the pullback. Conversely, a vector in $\ker dF_p$ would have positive pullback norm but also zero pullback norm, so the kernel is trivial.
\end{proof}

\begin{definition}[The metric induced by an immersion]
\label{def:induced-metric}
\lean{LeeLib.Ch02.pullbackMetric}
\uses{lem:pullback-immersion-metric, def:riemannian-metric}
\leanok



\dcref{ch2:2.11}
For a smooth immersion $F:M\to(\widetilde M,\widetilde g)$ between manifolds with or without boundary, the induced metric is the Riemannian metric $F^*\widetilde g$. An immersion or embedding with $F^*\widetilde g=g$ is respectively an isometric immersion or embedding. For a submanifold, the induced metric is the restriction of the ambient metric to tangent vectors.
\end{definition}

\begin{definition}[Spheres]
\label{ex:spheres}
\lean{LeeLib.Ch02.roundMetric}
\uses{def:induced-metric, ex:euclidean-metric}
\leanok



\dcref{ch2:2.13}
For $n>0$, the unit sphere $S^n\subseteq\mathbb R^{n+1}$ carries the round (standard) metric induced by its inclusion $\iota$. Its differential identifies $T_pS^n$ with $(\mathbb Rp)^\perp$, so $\iota^*\bar g$ is positive definite by Lemma \ref{lem:pullback-immersion-metric}.
\end{definition}

\begin{proposition}[Existence of Adapted Orthonormal Frames]
\label{prop:adapted-orthonormal-frames}
\uses{def:riemannian-metric, lem:partial-frame-extension, lem:pushforward-smooth, lem:ambient-tangent-bundle-riemannian, prop:existence-orthonormal-frames, lem:pullback-immersion-metric}
\dcref{ch2:2.14}

If $(\widetilde M,\widetilde g)$ has no boundary and $M\subseteq\widetilde M$ is an embedded smooth submanifold with or without boundary, then for every $p\in M$ some ambient neighborhood of $p$ has a smooth orthonormal frame whose first $\dim M$ members are tangent to $M$.
\end{proposition}

For a vector bundle $V\to B$, a local $k$-frame means $k$ smooth sections linearly independent in every fibre (not necessarily spanning). For an immersion, adapted frames live in the ambient pullback bundle of Lemma \ref{lem:ambient-tangent-bundle-riemannian}.

\begin{lemma}[Ordered extension of a linearly independent family]
\label{lem:ordered-linear-extension}
\lean{LeeLib.Ch02.exists_linearIndependent_castAdd}
\leanok


Let $\dim V=m$, let $k,d\ge0$ with $k+d\le m$, and let $(c_1,\ldots,c_k)$ be linearly independent. There is a linearly independent $(w_1,\ldots,w_{k+d})$ with $w_i=c_i$ for $i\le k$.
\end{lemma}
\begin{proof}
\leanok
\sketch Induct on $d$. For the step, choose $y$ outside the span of the given family; appending $y$ preserves linear independence, and apply the induction hypothesis to the enlarged family.
\end{proof}

\begin{lemma}[Linear independence of smooth sections is an open condition]
\label{lem:open-linear-independence}
\lean{LeeLib.Ch02.exists_isLocalIndepOn_nhds}
\leanok


If smooth sections $(s_1,\ldots,s_k)$ of a vector bundle are linearly independent at $x_0$, they are linearly independent on some neighborhood of $x_0$.
\end{lemma}
\begin{proof}
\leanok
\sketch In a fibre trivialization, an independent tuple defines an injective map $L:\mathbb R^k\to F$ and hence $\kappa>0$ with $\|La\|\ge\kappa\|a\|$. Perturbations with $\sum_i\|v_i'-v_i\|<\kappa$ remain independent since $\|L'a\|\ge(\kappa-\delta)\|a\|$. Continuity of the section coordinates then gives the required neighborhood.
\end{proof}

\begin{lemma}[Extension of a partial frame]
\label{lem:partial-frame-extension}
\lean{LeeLib.Ch02.exists_isLocalFrameOn_extend}
\uses{lem:ordered-linear-extension, lem:open-linear-independence}
\leanok



If $(s_1,\ldots,s_k)$ is a local $k$-frame of a rank-$m$ bundle over $u$, then $k\le m$ and each $x_0\in u$ has a neighborhood on which it extends to a smooth frame $(E_1,\ldots,E_m)$ with $E_i=s_i$ for $i\le k$.
\end{lemma}
\begin{proof}
\leanok
\sketch Trivialize near $x_0$, extend the fibre tuple by Lemma \ref{lem:ordered-linear-extension} to a basis, and define the extra sections to be constant in the trivialization. Lemma \ref{lem:open-linear-independence} shrinks the neighborhood so that the resulting $m$ sections remain a frame.
\end{proof}

\begin{lemma}[Smoothness of the pushforward of a vector field]
\label{lem:pushforward-smooth}
\lean{LeeLib.Ch02.contMDiffOn_pushforward}
\uses{lem:ambient-tangent-bundle-riemannian}
\leanok



If $F:M\to\widetilde M$ is smooth and $X$ is a smooth vector field on an open $u\subseteq M$, then $x\mapsto dF_x(X_x)$ is a smooth section of $F^*T\widetilde M$ over $u$.
\end{lemma}
\begin{proof}
\leanok
\sketch In local trivializations, the section has coordinate expression $D(x)\xi(x)$, where $\xi$ represents $X$ and $D(x)=\Theta_{F(x)}\circ dF_x\circ\Sigma_x^{-1}$ represents the differential of $F$. Both factors are smooth.
\end{proof}

\begin{lemma}[The Ambient Tangent Bundle Along a Submanifold]
\label{lem:ambient-tangent-bundle-riemannian}
\lean{Bundle.ContMDiffRiemannianMetric.pullback}
\uses{def:riemannian-metric, def:induced-metric}
\leanok



For a Riemannian $(\widetilde M,\widetilde g)$ and smooth $F:M\to\widetilde M$, the pullback bundle $F^*T\widetilde M$, with fibre $T_{F(p)}\widetilde M$, has the smoothly varying fibre metric $\widetilde g_{F(p)}$. For an inclusion this is $T\widetilde M|_M$, containing the orthogonal splitting into tangent and normal spaces and the projections of Proposition \ref{prop:normal-bundle}.
\end{lemma}
\begin{proof}
\leanok
\uses{def:riemannian-metric}
\sketch In a pulled-back trivialization, the fibre metric is the coordinate representation of $\widetilde g$ evaluated at $F(p)$, hence smooth; symmetry and positivity are inherited fibrewise.
\end{proof}

\begin{proposition}[Existence of Adapted Orthonormal Frames, immersed form]
\label{prop:adapted-orthonormal-frames-immersed}
\lean{LeeLib.Ch02.exists_adapted_orthonormalFrame}
\uses{def:riemannian-metric, lem:partial-frame-extension, lem:pushforward-smooth, lem:ambient-tangent-bundle-riemannian, prop:existence-orthonormal-frames, lem:pullback-immersion-metric}
\leanok



If $F:M^n\to\widetilde M^m$ is a smooth immersion, then $n\le m$ and every $p\in M$ has a neighborhood carrying an orthonormal frame $(E_1,\ldots,E_m)$ of the bundle $F^*T\widetilde M$ from Lemma \ref{lem:ambient-tangent-bundle-riemannian} such that
\[
\operatorname{span}(E_1|_x,\ldots,E_n|_x)=\operatorname{im}dF_x
\]
for every $x$ in the neighborhood.
\end{proposition}
\begin{proof}
\leanok
\uses{lem:partial-frame-extension,lem:pushforward-smooth,lem:ambient-tangent-bundle-riemannian,prop:existence-orthonormal-frames}
\sketch Push forward a local frame of $TM$ to a local $n$-frame of $F^*T\widetilde M$ using Lemma \ref{lem:pushforward-smooth}; independence follows from immersion. Extend it by Lemma \ref{lem:partial-frame-extension}. Lemma \ref{lem:ambient-tangent-bundle-riemannian} equips the bundle with the ambient metric, so Proposition \ref{prop:existence-orthonormal-frames} applies fibrewise. The Gram--Schmidt flag puts the first $n$ vectors in $\operatorname{im}dF_x$; both subspaces have dimension $n$, so they coincide.
\end{proof}

\begin{definition}[The Normal Space]
\label{def:normal-space}
\lean{LeeLib.Ch02.normalSpace}
\uses{lem:ambient-tangent-bundle-riemannian}



For an immersion $F:M^n\to\widetilde M^m$, identify $T_xM$ with $\operatorname{im}dF_x\subseteq T_{F(x)}\widetilde M$. The normal space is
\[
N_xM=\{v\in T_{F(x)}\widetilde M:\widetilde g_{F(x)}(v,w)=0\text{ for every }w\in T_xM\}.
\]
It is a linear subspace, and $T_{F(x)}\widetilde M=T_xM\oplus N_xM$ orthogonally.
\end{definition}

\begin{lemma}[Local Orthonormal Frames for the Normal Space]
\label{lem:normal-space-orthonormal-frame}
\lean{LeeLib.Ch02.exists_orthonormalFrame_normalSpace}
\uses{def:normal-space, prop:adapted-orthonormal-frames-immersed}
\leanok



For an immersion $F:M^n\to\widetilde M^m$, $n\le m$ and every point has a neighborhood with smooth sections $N_1,\ldots,N_{m-n}$ of $F^*T\widetilde M$ whose values form an orthonormal basis of $N_xM$ at each $x$.
\end{lemma}
\begin{proof}
\leanok
\uses{prop:adapted-orthonormal-frames-immersed}
\sketch Use Proposition \ref{prop:adapted-orthonormal-frames-immersed} and take $N_j=E_{n+j}$. Orthonormality makes these sections normal to the span of the first $n$ vectors; expanding any normal vector in the full orthonormal frame forces its first $n$ coefficients to vanish.
\end{proof}

\begin{proposition}[The Normal Bundle]
\label{prop:normal-bundle}
\lean{LeeLib.Ch02.contMDiffVectorBundle_normalSpace}
\uses{lem:ambient-tangent-bundle-riemannian, def:normal-space, lem:normal-space-orthonormal-frame, lem:local-frame-criterion-for-subbundles}
\leanok



\dcref{ch2:2.16}
If $M^n$ is an immersed or embedded submanifold with or without boundary of a Riemannian $m$-manifold $\widetilde M$, then $NM=\bigsqcup_{p\in M}N_pM$ is a smooth rank-$(m-n)$ vector subbundle of $T\widetilde M|_M$.
\end{proposition}
\begin{proof}
\leanok
\uses{lem:normal-space-orthonormal-frame,lem:local-frame-criterion-for-subbundles,lem:ambient-tangent-bundle-riemannian}
\sketch The ambient tangent bundle is Riemannian by Lemma \ref{lem:ambient-tangent-bundle-riemannian}, so Lemma \ref{lem:local-frame-criterion-for-subbundles} applies. Lemma \ref{lem:normal-space-orthonormal-frame} supplies local smooth frames of all fibres $N_pM$, of dimension $m-n$; Lemma \ref{lem:local-frame-criterion-for-subbundles} therefore makes their union a smooth subbundle.
\end{proof}

\begin{remark}
\label{rem:normal-bundle-projections}
\uses{prop:normal-bundle}

The orthogonal splitting defines smooth bundle maps
\[
\pi^\top:T\widetilde M|_M\to TM,\qquad \pi^\perp:T\widetilde M|_M\to NM,
\]
whose fibres are the projections onto tangent and normal spaces. In an adapted orthonormal frame,
\[
\pi^\top\!\left(\sum_{a=1}^mX^aE_a\right)=\sum_{a=1}^nX^aE_a,\qquad
\pi^\perp\!\left(\sum_{a=1}^mX^aE_a\right)=\sum_{a=n+1}^mX^aE_a.
\]
\end{remark}

\begin{proof}
\uses{thm:local-embedding-theorem,lem:local-frame-criterion-for-subbundles,prop:adapted-orthonormal-frames-immersed,lem:normal-space-orthonormal-frame}
\sketch Locally embed the submanifold and use an adapted orthonormal frame; the last $m-n$ restrictions form a smooth normal frame, so the local-frame criterion gives the normal subbundle. The displayed coordinate truncations define the unique orthogonal projections and show their smoothness.
\end{proof}

\begin{proposition}[Existence of Outward-Pointing Normal]
\label{prop:outward-pointing-normal}
\dcref{ch2:2.17}
If $(M,g)$ is a smooth Riemannian manifold with boundary, $N\partial M$ is a smooth rank-$1$ bundle over $\partial M$ and has a unique smooth outward-pointing unit normal field.
\end{proposition}

For a local parametrization $X:U\to\widetilde M$ of an immersed $n$-submanifold in $\mathbb R^m$, where $U\subseteq\mathbb R^n$ (or $U\subseteq\mathbb R^n_+$ at the boundary),
\[
X^*\bar g=\sum_{i=1}^m(dX^i)^2
=\sum_{i=1}^m\sum_{j,k=1}^n\frac{\partial X^i}{\partial u^j}\frac{\partial X^i}{\partial u^k}\,du^jdu^k.
\]

\begin{example}[Metrics in Graph Coordinates]
\label{ex:metrics-in-graph-coordinates}
\lean{LeeLib.Ch02.graphMap, LeeLib.Ch02.graphMetric}
\uses{def:induced-metric, ex:euclidean-metric}
\leanok
\dcref{ch2:2.19}



For smooth $f:U\subseteq\mathbb R^n\to\mathbb R$, its graph has parametrization $X(u)=(u,f(u))$ and induced metric
\[
X^*\bar g=(du^1)^2+\cdots+(du^n)^2+df^2.
\]
On the upper hemisphere, for $u^2+v^2<1$ and $X(u,v)=(u,v,\sqrt{1-u^2-v^2})$,
\[
\hat g=du^2+dv^2+\left(\frac{u\,du+v\,dv}{\sqrt{1-u^2-v^2}}\right)^2
 =\frac{(1-v^2)du^2+(1-u^2)dv^2+2uv\,du\,dv}{1-u^2-v^2}.
\]
\end{example}

\begin{example}[Surfaces of Revolution]
\label{ex:surfaces-of-revolution}
\lean{LeeLib.Ch02.surfaceOfRevolutionMap, LeeLib.Ch02.surfaceOfRevolutionMetric, LeeLib.Ch02.surfaceOfRevolutionMetric_unitSpeed, LeeLib.Ch02.surfaceOfRevolutionMetric_sphere, LeeLib.Ch02.surfaceOfRevolutionMetric_cylinder, LeeLib.Ch02.surfaceOfRevolutionMetric_torus}
\uses{def:induced-metric, ex:euclidean-metric}
\dcref{ch2:2.20}



For an embedded curve $C\subset\{(r,z):r>0\}$ with local parametrization $\gamma(t)=(a(t),b(t))$, the surface of revolution is parametrized by
\[
X(t,\theta)=(a(t)\cos\theta,a(t)\sin\theta,b(t)).\tag{2.11}
\]
Its induced metric is
\[
X^*\bar g=(a'(t)^2+b'(t)^2)dt^2+a(t)^2d\theta^2,
\]
and for unit-speed $\gamma$ this is $dt^2+a(t)^2d\theta^2$. For the semicircle $\gamma(\varphi)=(\sin\varphi,\cos\varphi)$ it is the sphere metric $d\varphi^2+\sin^2\varphi\,d\theta^2$; for $\gamma(t)=(2+\cos t,\sin t)$ it is the torus metric $dt^2+(2+\cos t)^2d\theta^2$; for $\gamma(t)=(1,t)$ it is the cylinder metric $dt^2+d\theta^2$, so the parametrization is an isometric immersion.
\end{example}

\begin{example}[The n-Torus as a Riemannian Submanifold]
\label{ex:n-torus-riemannian-submanifold}
\uses{ex:covering-n-torus}
\dcref{ch2:2.21}

The covering map $X:\mathbb R^n\to\mathbb T^n$ restricts locally to parametrizations whose induced metric is Euclidean in the coordinates $(u^i)$; hence the induced metric on $\mathbb T^n$ is flat.
\end{example}

\section{Riemannian Products}

For Riemannian manifolds $(M_1,g_1)$ and $(M_2,g_2)$, identify
$T_{(p_1,p_2)}(M_1\times M_2)$ with
$T_{p_1}M_1\oplus T_{p_2}M_2$. The product metric is
\begin{equation}
g_{(p_1,p_2)}((v_1,v_2),(w_1,w_2))=g_1|_{p_1}(v_1,w_1)+g_2|_{p_2}(v_2,w_2). \tag{2.12}
\end{equation}
In product coordinates its matrix is block diagonal,
\[
(g_{ij})=\begin{pmatrix}(g_1)_{ab}&0\\0&(g_2)_{cd}\end{pmatrix},
\]
and products of more factors are defined analogously. For the projections $\pi_i$,
\begin{equation}
g_1\oplus g_2=\pi_1^*g_1+\pi_2^*g_2. \tag{2.12'}
\end{equation}

\begin{lemma}[Sums of smooth sections]
\label{lem:smooth-tensor-field-add}
\lean{ContMDiff.add_section}

\mathlibok
Let $V$ be a smooth vector bundle over a smooth manifold $B$. If $s,t$ are smooth sections of $V$, then $s+t$ is smooth.
\end{lemma}

\begin{lemma}[Smooth rescalings of smooth sections]
\label{lem:smooth-tensor-field-smul}
\lean{ContMDiff.smul_section}

\mathlibok
Let $V$ be a smooth vector bundle over a smooth manifold $B$, let $s$ be a smooth section, and let $c:B\to\mathbb R$ be smooth. Then $b\mapsto c(b)s_b$ is smooth.
\end{lemma}

\begin{lemma}[Positive definiteness of the product form]
\label{lem:product-form-positive-definite}
\lean{LeeLib.Ch02.prodForm_pos}
\uses{def:riemannian-metric}
\leanok



For Riemannian manifolds $(M_1,g_1)$ and $(M_2,g_2)$, the form $g_1\oplus g_2$ defined by (2.12) is positive definite on $M_1\times M_2$.
\end{lemma}

\begin{proof}
\leanok
\uses{def:riemannian-metric}
\sketch For nonzero $v=(v_1,v_2)$, at least one component is nonzero, and
\[
(g_1\oplus g_2)(v,v)=g_1(v_1,v_1)+g_2(v_2,v_2)>0
\]
because both summands are nonnegative and a nonzero component gives a strictly positive summand.
\end{proof}

\begin{definition}[The Product Metric]
\label{def:product-metric}
\lean{LeeLib.Ch02.prodMetric}
\uses{def:riemannian-metric, lem:product-form-positive-definite, lem:smooth-tensor-field-add, lem:pullback-immersion-metric}
\leanok



For Riemannian manifolds $(M_1,g_1)$ and $(M_2,g_2)$, the \emph{product metric} on $M_1\times M_2$ is the metric (2.12). It is symmetric, positive definite by Lemma \ref{lem:product-form-positive-definite}, and smooth by (2.12') and Lemma \ref{lem:smooth-tensor-field-add}, since each pullback summand is smooth.
\end{definition}

\begin{proposition}[Induced Metric on $T^n$]
\label{prop:product-metric-torus}


\dcref{ch2:2.23}
\uses{def:product-metric,def:induced-metric}
The metric on $T^n=(S^1)^n$ induced by its standard product embedding equals
the product of the usual induced metrics on the factors
$S^1\subseteq\mathbb R^2$. In the formalized case $n=2$,
$T^2=S^1\times S^1$:  is the round metric on
$S^1\subseteq\mathbb C$,  is their product
metric, and  identifies it
with the metric induced by the $L^2$-embedding into $\mathbb R^4$.
\end{proposition}

For a strictly positive smooth $f:M_1\to\mathbb R^+$, define the warped form on $M_1\times M_2$ by
\begin{equation}
g_{(p_1,p_2)}((v_1,v_2),(w_1,w_2))=g_1|_{p_1}(v_1,w_1)+f(p_1)^2g_2|_{p_2}(v_2,w_2).
\end{equation}
Unless $f$ is constant, this is generally not a product metric.

\begin{lemma}[Positive definiteness of the warped product form]
\label{lem:warped-product-form-positive-definite}
\lean{LeeLib.Ch02.warpedProductForm_pos}
\uses{def:riemannian-metric}
\leanok



If $f:M_1\to\mathbb R$ is nowhere vanishing and $(M_i,g_i)$ are Riemannian, then $g_1\oplus f^2g_2$ is positive definite.
\end{lemma}

\begin{proof}
\leanok
\uses{def:riemannian-metric,lem:product-form-positive-definite}
\sketch For nonzero $(v_1,v_2)$,
\[
(g_1\oplus f^2g_2)(v,v)=g_1(v_1,v_1)+f(p_1)^2g_2(v_2,v_2).
\]
Here $f(p_1)^2>0$; both summands are nonnegative and one is strictly positive.
\end{proof}

\begin{definition}[Warped Products]
\label{def:warped-product}
\lean{LeeLib.Ch02.warpedProductMetric}
\uses{def:product-metric, def:riemannian-metric, lem:warped-product-form-positive-definite, lem:smooth-tensor-field-add, lem:smooth-tensor-field-smul}
\leanok



For Riemannian manifolds $(M_1,g_1)$ and $(M_2,g_2)$ and a strictly positive smooth $f:M_1\to\mathbb R^+$, the \emph{warped product} $M_1\times_fM_2$ is $M_1\times M_2$ with metric $g_1\oplus f^2g_2$. It is smooth because
\[
g_1\oplus f^2g_2=\pi_1^*g_1+(f\circ\pi_1)^2\pi_2^*g_2,
\]
using Lemmas \ref{lem:smooth-tensor-field-smul} and \ref{lem:smooth-tensor-field-add}; positivity is Lemma \ref{lem:warped-product-form-positive-definite}. Since only $f^2$ occurs, a nowhere-vanishing smooth real $f$ also suffices.
\end{definition}

\begin{proposition}[Constant warping gives the product metric]
\label{prop:warped-product-constant-one}
\lean{LeeLib.Ch02.warpedProductMetric_one}
\uses{def:warped-product, def:product-metric}
\leanok



If $f\equiv1$, the warped product metric is $g_1\oplus g_2$.
\end{proposition}

\begin{proof}
\leanok
\uses{def:warped-product,def:product-metric}
\sketch Substituting $f=1$ into the defining bilinear form gives (2.12) at every point, hence the metrics coincide.
\end{proof}

\begin{example}[Warped Products]
\label{ex:warped-products}
\uses{ex:surfaces-of-revolution, def:warped-product, prop:warped-product-constant-one, REF:ch2-3}
\dcref{ch2:2.24}

\begin{enumerate}
\item[(a)] $f\equiv1$ gives the product metric by Proposition \ref{prop:warped-product-constant-one}.
\item[(b)] If $H=\{(r,z):r>0\}$, $C\subset H$ is an embedded smooth curve
with its induced Euclidean metric, and $S_C$ is the corresponding surface of
revolution, then with the standard metric on $S^1$ and $f(r,z)=r$, the surface
$S_C$ is isometric to $C\times_fS^1$.
\item[(c)] For the standard coordinate $\rho$ on $\mathbb R^+$, the map
$(\rho,\omega)\mapsto\rho\omega$ is an isometry from
$\mathbb R^+\times_\rho S^{n-1}$ to $\mathbb R^n\setminus\{0\}$ with its
Euclidean metric.
\end{enumerate}
\end{example}

\section{Riemannian Submersions}

Let $\pi:\tilde M\to M$ be a smooth submersion and let $\tilde g$ be a metric on $\tilde M$. Set
\[
V_x=\ker d\pi_x=T_x(\pi^{-1}(\pi(x))),\qquad H_x=V_x^{\perp_{\tilde g_x}}.
\]
Thus $T_x\tilde M=H_x\oplus V_x$. A vector field is horizontal (respectively vertical) when its values lie in $H_x$ (respectively $V_x$); a horizontal lift of $X\in\mathfrak X(M)$ is a horizontal field $\tilde X$ with $d\pi_x(\tilde X_x)=X_{\pi(x)}$.

\begin{proposition}[The Smooth Horizontal--Vertical Splitting of a Vector Field]
\label{prop:horizontal-vector-field-properties}
\lean{LeeLib.Ch02.existsUnique_horizontal_add_vertical_field}
\uses{def:submersion, def:vertical-space, def:horizontal-space, def:riemannian-metric, lem:tangent-space-horizontal-vertical-splitting, lem:horizontal-projection-field-smooth}
\leanok
\dcref{ch2:2.25}



Every smooth vector field $W$ on $\tilde M$ has a unique decomposition $W=W^H+W^V$ into smooth horizontal and vertical fields.
\end{proposition}

\begin{proof}
\leanok
\sketch Pointwise set $W_x^H=L_x(d\pi_xW_x)$ and $W_x^V=W_x-W_x^H$. The tangent-space splitting gives uniqueness and the horizontal-projection lemma gives smoothness of $W^H$; $W^V$ is then smooth.
\end{proof}

\begin{lemma}[The Pushforward of a Vector Field is a Smooth Section Along the Map]
\label{lem:pushforward-along-map-smooth}
\lean{LeeLib.Ch02.contMDiff_mfderivAlong}
\uses{def:submersion}
\leanok



For a smooth $F:\tilde M\to M$ and smooth $W$ on $\tilde M$, the map $x\mapsto dF_x(W_x)$ is a smooth section of $F^*TM$ along $F$.
\end{lemma}

\begin{proof}
\leanok
\sketch The bundled differential $(x,v)\mapsto(F(x),dF_xv)$ is smooth, as is $x\mapsto(x,W_x)$; compose them.
\end{proof}

\begin{lemma}[The Horizontal Lift of a Smooth Section Along $\pi$ is Smooth]
\label{lem:horizontal-lift-along-smooth}
\lean{LeeLib.Ch02.contMDiffAt_horizontalLiftAlong}
\uses{def:horizontal-lift-tangent, lem:horizontal-lift-smooth-in-data, def:submersion}
\leanok



If $\pi:\tilde M\to M$ is a smooth submersion with metric $\tilde g$ and $X_x\in T_{\pi(x)}M$ depends smoothly on $x$, then $x\mapsto L_x(X_x)$ is a smooth vector field.
\end{lemma}

\begin{proof}
\leanok
\sketch In tangent-bundle trivializations the lift is $\mathcal L(B(x),A(x))u(x)$, where $B(x)$ and $A(x)$ represent $\tilde g_x$ and $d\pi_x$, and $u(x)$ represents $X_x$. Lemma \ref{lem:horizontal-lift-smooth-in-data} gives smoothness.
\end{proof}

\begin{lemma}[Smoothness of the Horizontal Projection of a Vector Field]
\label{lem:horizontal-projection-field-smooth}
\lean{LeeLib.Ch02.contMDiff_horizontalProjField}
\uses{lem:pushforward-along-map-smooth, lem:horizontal-lift-along-smooth, def:horizontal-space}
\leanok



For smooth $W$ on $\tilde M$, the field $W_x^H=L_x(d\pi_x(W_x))$ is smooth.
\end{lemma}

\begin{proof}
\leanok
\sketch Lemma \ref{lem:pushforward-along-map-smooth} makes $x\mapsto d\pi_x(W_x)$ a smooth section along $\pi$, and Lemma \ref{lem:horizontal-lift-along-smooth} lifts it smoothly.
\end{proof}

\begin{definition}[Smooth submersion]
\label{def:submersion}
\lean{LeeLib.Ch02.IsSubmersion}
\leanok


A smooth map $\pi:\tilde M\to M$ is a \emph{submersion} if $d\pi_x:T_x\tilde M\to T_{\pi(x)}M$ is surjective for every $x$.
\end{definition}

\begin{definition}[The horizontal lift of a surjection]
\label{def:horizontal-lift-linear}
\lean{LeeLib.Ch02.horizontalLift}
\leanok


Let $B$ be a positive definite symmetric bilinear form on a finite-dimensional real vector space $E$, let $E'$ be finite-dimensional, and let $A:E\to E'$ be surjective. Writing $A^t:(E')^*\to E^*$ for the transpose and identifying $B$ with $E\to E^*$, $v\mapsto B(v,\cdot)$, the $B$-weighted horizontal lift (right inverse) is
\[
L=B^{-1}A^t(AB^{-1}A^t)^{-1}:E'\to E.
\]
Lemma \ref{lem:horizontal-lift-linear-spec} gives its characterization and existence of the inverses.
\end{definition}

\begin{lemma}[The horizontal lift is the unique $B$-orthogonal right inverse]
\label{lem:horizontal-lift-linear-spec}
\lean{LeeLib.Ch02.horizontalLift_unique}
\uses{def:horizontal-lift-linear}
\leanok



With the preceding hypotheses, there is exactly one $L:E'\to E$ such that
\begin{enumerate}
\item[(i)] $A(Lu)=u$ for all $u\in E'$, and
\item[(ii)] $B(Lu,w)=0$ for all $u\in E'$ and $w\in\ker A$,
\end{enumerate}
and it is the displayed map.
\end{lemma}

\begin{proof}
\leanok
\uses{def:horizontal-lift-linear}
\sketch Let $S=B^{-1}A^t$. Then $B(S\xi,w)=\xi(Aw)$. Positive definiteness and surjectivity imply $AS$ is injective, hence invertible in finite dimensions. Setting $L=S(AS)^{-1}$ proves (i) and (ii). If $L,L'$ satisfy (i)--(ii), $d=Lu-L'u\in\ker A$ and $B(d,d)=0$, so $d=0$.
\end{proof}

\begin{lemma}[The horizontal lift depends smoothly on the data]
\label{lem:horizontal-lift-smooth-in-data}
\lean{LeeLib.Ch02.contMDiffAt_horizontalLift}
\uses{def:horizontal-lift-linear}
\leanok



Let $N$ be a smooth manifold, and let $x\mapsto B_x$ and $x\mapsto A_x\in\operatorname{Hom}(E,E')$ be smooth families of bilinear forms on $E$ and linear maps. If $B_{x_0}$ is positive definite and $A_{x_0}$ surjective, then the corresponding lifts $L_x$ are smooth at $x_0$.
\end{lemma}

\begin{proof}
\leanok
\uses{def:horizontal-lift-linear,lem:horizontal-lift-linear-spec}
\sketch Use $L_x=B_x^{-1}A_x^t(A_xB_x^{-1}A_x^t)^{-1}$. Transpose and composition are smooth, and inversion is smooth at the invertible operators supplied by Lemma \ref{lem:horizontal-lift-linear-spec}.
\end{proof}

\begin{definition}[The vertical space]
\label{def:vertical-space}
\lean{LeeLib.Ch02.verticalSpace}
\uses{def:submersion}
\leanok



For a submersion $\pi:\tilde M\to M$, the \emph{vertical space} at $x$ is $V_x=\ker d\pi_x\subseteq T_x\tilde M$.
\end{definition}

\begin{definition}[The horizontal space]
\label{def:horizontal-space}
\lean{LeeLib.Ch02.horizontalSpace}
\uses{def:vertical-space, def:riemannian-metric}
\leanok



For a submersion $\pi$ with metric $\tilde g$, the \emph{horizontal space} at $x$ is
\[
H_x=\{v\in T_x\tilde M:\tilde g_x(v,w)=0\text{ for every }w\in V_x\}.
\]
\end{definition}

\begin{definition}[The horizontal lift of a tangent vector]
\label{def:horizontal-lift-tangent}
\lean{LeeLib.Ch02.horizontalLiftAt}
\uses{def:horizontal-lift-linear, def:horizontal-space, def:submersion}
\leanok



For a smooth submersion $\pi:\tilde M\to M$ with metric $\tilde g$, the horizontal lift at $x$ is $L_x:T_{\pi(x)}M\to T_x\tilde M$ obtained from Definition \ref{def:horizontal-lift-linear} with $B=\tilde g_x$ and $A=d\pi_x$. It is the unique map with $d\pi_xL_xu=u$ and $L_xu\in H_x$.
\end{definition}

\begin{lemma}[The tangent space splits]
\label{lem:tangent-space-horizontal-vertical-splitting}
\lean{LeeLib.Ch02.isCompl_horizontalSpace_verticalSpace}
\uses{def:vertical-space, def:horizontal-space, def:horizontal-lift-tangent}
\leanok



For every $x$, $T_x\tilde M=H_x\oplus V_x$.
\end{lemma}

\begin{proof}
\leanok
\uses{def:vertical-space,def:horizontal-space,def:horizontal-lift-tangent,lem:horizontal-lift-linear-spec}
\sketch A vector in $H_x\cap V_x$ has zero $\tilde g_x$-norm, hence is zero. For $v\in T_x\tilde M$, set $v^H=L_x(d\pi_xv)$ and $v^V=v-v^H$; Lemma \ref{lem:horizontal-lift-linear-spec} gives $v^H\in H_x$ and $v^V\in V_x$.
\end{proof}

\begin{proposition}[Existence and uniqueness of horizontal lifts]
\label{prop:horizontal-lift-of-vector-field}
\lean{LeeLib.Ch02.exists_unique_horizontalLift}
\uses{def:horizontal-lift-tangent, def:horizontal-space, def:submersion, def:riemannian-metric, lem:horizontal-lift-linear-spec, lem:horizontal-lift-smooth-in-data}
\leanok



Every $X\in\mathfrak X(M)$ has a unique smooth horizontal lift $\tilde X$ satisfying $d\pi_x(\tilde X_x)=X_{\pi(x)}$.
\end{proposition}

\begin{proof}
\leanok
\uses{def:horizontal-lift-tangent,lem:horizontal-lift-linear-spec,lem:horizontal-lift-smooth-in-data,def:horizontal-space}
\sketch Define $\tilde X_x=L_x(X_{\pi(x)})$. The right-inverse and orthogonality properties give the lift conditions. In local tangent trivializations, $L_x$ is the smooth lift of the smooth families representing $\tilde g_x$ and $d\pi_x$, so $\tilde X$ is smooth. Pointwise uniqueness follows from Lemma \ref{lem:horizontal-lift-linear-spec}.
\end{proof}

\begin{proposition}[Every Horizontal Vector is the Value of a Lift]
\label{prop:horizontal-lift-realizes-horizontal-vector}
\lean{LeeLib.Ch02.exists_horizontalLift_eq}
\uses{def:horizontal-space, def:horizontal-lift-tangent, prop:horizontal-lift-of-vector-field, lem:vector-field-extension-point, lem:horizontal-lift-linear-spec}
\leanok
\dcref{ch2:2.25}



For $x\in\tilde M$ and $v\in H_x$, there is $X\in\mathfrak X(M)$ whose horizontal lift satisfies $\tilde X_x=v$.
\end{proposition}

\begin{proof}
\leanok
\sketch Extend $d\pi_xv$ to a smooth field $X$ on $M$ using Lemma \ref{lem:vector-field-extension-point}. Its lift satisfies $\tilde X_x=L_x(d\pi_xv)=v$ because $v$ is horizontal and $L_x$ is the inverse of $d\pi_x$ on $H_x$.
\end{proof}

\begin{proposition}
\label{prop:horizontal-vector-field-not-lift}


\dcref{ch2:2.26}
\uses{def:horizontal-space,def:horizontal-lift-tangent,ex:riemannian-submersions}
For $\pi:\mathbb R^2\to\mathbb R$, $\pi(x,y)=x$, the field
$\tilde W=y\partial_x$ is horizontal but is not the horizontal lift of any
vector field on $\mathbb R$.
\end{proposition}

\begin{example}[Riemannian Submersions]
\label{ex:riemannian-submersions}
\lean{LeeLib.Ch02.isRiemannianSubmersion_euclidean_fst, LeeLib.Ch02.isRiemannianSubmersion_prodMetric_fst, LeeLib.Ch02.isRiemannianSubmersion_prodMetric_snd, LeeLib.Ch02.isRiemannianSubmersion_warpedProductMetric_fst}
\uses{def:riemannian-submersion, def:product-metric, def:warped-product}
\leanok


\dcref{ch2:2.27}

\begin{enumerate}
\item[(a)] The projection $\mathbb R^{n+k}\to\mathbb R^n$ is a Riemannian submersion for Euclidean metrics.
\item[(b)] Both projections from a product metric $M\times N$ are Riemannian submersions.
\item[(c)] For a warped product $M\times_fN$, the projection to $M$ is a Riemannian submersion; the projection to $N$ generally is not. In (a), the Euclidean metric is the product metric (not the sup-norm on a plain product); in (c), only the first assertion is formalized, with horizontal spaces given by the factor subspaces ().
\end{enumerate}
\end{example}

\begin{definition}[Riemannian submersion]
\label{def:riemannian-submersion}
\lean{LeeLib.Ch02.IsRiemannianSubmersion}
\uses{def:submersion, def:horizontal-space, def:riemannian-metric}
\leanok



A smooth submersion $\pi:(\tilde M,\tilde g)\to(M,g)$ is a \emph{Riemannian submersion} if
\[
g_{\pi(x)}(d\pi_xv,d\pi_xw)=\tilde g_x(v,w)\qquad(v,w\in H_x).
\]
\end{definition}

\begin{definition}[The inner product induced on the base]
\label{def:submersion-inner}
\lean{LeeLib.Ch02.submersionInner}
\uses{def:horizontal-lift-tangent, def:riemannian-metric}
\leanok



For a submersion with metric $\tilde g$, define the bilinear form on $T_{\pi(x)}M$ by
\[
\beta_x(u,v)=\tilde g_x(L_xu,L_xv).
\]
\end{definition}

\begin{lemma}[A Riemannian submersion's base metric is the induced form]
\label{lem:riemannian-submersion-inner-eq}
\lean{LeeLib.Ch02.IsRiemannianSubmersion.inner_eq}
\uses{def:riemannian-submersion, def:submersion-inner}
\leanok



If $\pi:(\tilde M,\tilde g)\to(M,g)$ is a Riemannian submersion, then $g_{\pi(x)}=\beta_x$.
\end{lemma}

\begin{proof}
\leanok
\uses{def:riemannian-submersion,def:submersion-inner,lem:horizontal-lift-linear-spec}
\sketch Apply the submersion identity to the horizontal vectors $L_xu,L_xv$ and use $d\pi_xL_xu=u$ and $d\pi_xL_xv=v$.
\end{proof}

\begin{proposition}[Uniqueness of the submersion metric]
\label{prop:riemannian-submersion-metric-unique}
\lean{LeeLib.Ch02.isRiemannianSubmersion_unique}
\uses{lem:riemannian-submersion-inner-eq}
\leanok



For a Riemannian metric $\tilde g$ on $\tilde M$ and a surjective smooth submersion $\pi:\tilde M\to M$, at most one metric on $M$ makes $\pi$ a Riemannian submersion.
\end{proposition}

\begin{proof}
\leanok
\uses{lem:riemannian-submersion-inner-eq}
\sketch If $g_1,g_2$ both work and $y=\pi(x)$, Lemma \ref{lem:riemannian-submersion-inner-eq} gives $(g_1)_y=\beta_x=(g_2)_y$.
\end{proof}

For a smooth action of a group $G$ on $\tilde M$, call it \emph{vertical} when $\pi(a\cdot x)=\pi(x)$, \emph{transitive on fibers} when any two points in one fiber differ by an element of $G$, and \emph{isometric} when each action map preserves $\tilde g$.

\begin{lemma}[The action's differential is invertible]
\label{lem:mfderiv-smul-inv-smul}
\lean{LeeLib.Ch02.IsIsometricAction.mfderiv_smul_inv_smul}
\uses{def:riemannian-metric}
\leanok



For a smooth action, $d(a\cdot {-})_x$ and $d(a^{-1}\cdot {-})_{a\cdot x}$ are mutually inverse.
\end{lemma}

\begin{proof}
\leanok
\sketch Differentiate $(a\cdot {-})\circ(a^{-1}\cdot {-})=\operatorname{Id}$ and use the chain rule.
\end{proof}

\begin{lemma}[A vertical action preserves the splitting]
\label{lem:horizontal-space-smul}
\lean{LeeLib.Ch02.horizontalSpace_smul}
\uses{lem:mfderiv-smul-inv-smul, def:vertical-space, def:horizontal-space}
\leanok



If the action preserves $\tilde g$ and is vertical, then $d(a\cdot {-})_x$ maps $V_x$ onto $V_{a\cdot x}$ and maps $H_x$ into $H_{a\cdot x}$.
\end{lemma}

\begin{proof}
\leanok
\uses{lem:mfderiv-smul-inv-smul,def:vertical-space,def:horizontal-space}
\sketch Differentiate $\pi\circ(a\cdot {-})=\pi$ to obtain $d\pi_{a\cdot x}\,d(a\cdot {-})_x=d\pi_x$, proving the vertical assertion using the inverse differential. If $v\in H_x$ and $w=d(a\cdot {-})_xz$ with $z\in V_x$, metric invariance gives $\tilde g_{a\cdot x}(d(a\cdot {-})_xv,w)=\tilde g_x(v,z)=0$.
\end{proof}

\begin{lemma}[The induced form is constant on fibers]
\label{lem:submersion-inner-smul}
\lean{LeeLib.Ch02.submersionInner_smul}
\uses{lem:horizontal-space-smul, def:submersion-inner, lem:horizontal-lift-linear-spec}
\leanok



Under the hypotheses of Lemma \ref{lem:horizontal-space-smul}, $\beta_{a\cdot x}=\beta_x$ on $T_{\pi(x)}M$.
\end{lemma}

\begin{proof}
\leanok
\uses{lem:horizontal-space-smul,def:submersion-inner,lem:horizontal-lift-linear-spec}
\sketch With $\Phi=d(a\cdot {-})_x$, the map $\Phi L_x$ is a horizontal right inverse of $d\pi_{a\cdot x}$, hence equals $L_{a\cdot x}$ by Lemma \ref{lem:horizontal-lift-linear-spec}. Metric invariance then gives $\beta_{a\cdot x}(u,v)=\beta_x(u,v)$.
\end{proof}

\begin{lemma}[The induced form depends only on the fiber]
\label{lem:submersion-inner-congr}
\lean{LeeLib.Ch02.submersionInner_congr}
\uses{lem:submersion-inner-smul, def:submersion-inner, lem:horizontal-space-smul}
\leanok



If the hypotheses of Lemma \ref{lem:horizontal-space-smul} hold and $G$ is transitive on fibers, then $\pi(x)=\pi(y)$ implies $\beta_x=\beta_y$.
\end{lemma}

\begin{proof}
\leanok
\uses{lem:submersion-inner-smul,def:submersion-inner}
\sketch Choose $a$ with $y=a\cdot x$ and apply Lemma \ref{lem:submersion-inner-smul}.
\end{proof}

\begin{lemma}[The horizontal lift is smooth in tangent coordinates]
\label{lem:horizontal-lift-in-coordinates-smooth}
\lean{LeeLib.Ch02.contMDiffAt_horizontalLiftAt_inTangentCoordinates}
\uses{lem:horizontal-lift-smooth-in-data, def:horizontal-lift-tangent, lem:horizontal-lift-linear-spec}
\leanok



For a submersion of $(\tilde M,\tilde g)$ and $x_0\in\tilde M$, the family $L_x:T_{\pi(x)}M\to T_x\tilde M$ is smooth at $x_0$ when read in tangent coordinates at $x_0$ and $\pi(x_0)$.
\end{lemma}

\begin{proof}
\leanok
\uses{lem:horizontal-lift-smooth-in-data,def:horizontal-lift-tangent,lem:horizontal-lift-linear-spec}
\sketch In tangent trivializations, $L_x$ is the lift of the smooth families representing $\tilde g_x$ and $d\pi_x$. Their value at $x_0$ is positive definite and surjective, so Lemma \ref{lem:horizontal-lift-smooth-in-data} applies.
\end{proof}

\begin{theorem}
\label{thm:group-action-riemannian-submersion}
\lean{LeeLib.Ch02.existsUnique_isRiemannianSubmersion_metric}
\uses{prop:riemannian-submersion-metric-unique, lem:submersion-inner-congr, prop:local-section-manifold, def:riemannian-submersion, lem:bilinear-comp-of-family-smooth, lem:horizontal-lift-in-coordinates-smooth, REF:ch2-6}
\leanok
\dcref{ch2:2.28}



Let $(\tilde M,\tilde g)$ be Riemannian, let $\pi:\tilde M\to M$ be a surjective smooth submersion, and let $G$ act on $\tilde M$. If the action is isometric, vertical, and transitive on fibers, then a unique Riemannian metric on $M$ makes $\pi$ a Riemannian submersion.
\end{theorem}

\begin{proof}
\leanok
\uses{prop:riemannian-submersion-metric-unique,lem:submersion-inner-congr,prop:local-section-manifold,def:submersion-inner,def:riemannian-submersion,lem:bilinear-comp-of-family-smooth,lem:horizontal-lift-in-coordinates-smooth}
\sketch Define $g_y=\beta_x$ for any $x$ with $\pi(x)=y$; this is well defined by Lemma \ref{lem:submersion-inner-congr}, symmetric, and positive definite because $L_xu\ne0$ for $u\ne0$. A local section $\sigma$ of $\pi$ gives $g_y(u,v)=\tilde g_{\sigma(y)}(L_{\sigma(y)}u,L_{\sigma(y)}v)$. Lemma \ref{lem:horizontal-lift-in-coordinates-smooth} and Lemma \ref{lem:bilinear-comp-of-family-smooth} make this smooth. Finally the defining identity holds on horizontal vectors since $L_xd\pi_x$ is the identity on $H_x$. Uniqueness is Proposition \ref{prop:riemannian-submersion-metric-unique}.
\end{proof}

\begin{corollary}
\label{cor:quotient-manifold-lie-group-riemannian-metric}
\uses{thm:quotient-manifold-theorem, thm:group-action-riemannian-submersion}
\dcref{ch2:2.29}

If a Lie group $G$ acts smoothly, freely, properly, and isometrically on $(\tilde M,\tilde g)$, then $\tilde M/G$ has a unique smooth manifold structure and Riemannian metric for which the quotient map is a Riemannian submersion.
\end{corollary}

\begin{proof}
\sketch The quotient manifold theorem gives the smooth structure and a smooth submersion quotient map. The action is vertical and fiber-transitive by definition, so Theorem \ref{thm:group-action-riemannian-submersion} supplies the unique metric.
\end{proof}

\begin{example}[The Fubini--Study Metric]
\label{ex:fubini-study-metric}
\uses{ex:complex-projective-spaces, thm:local-section-theorem, thm:group-action-riemannian-submersion}
\dcref{ch2:2.30}

For $n>0$, let $p:S^{2n+1}\to\mathbb{CP}^n$ be the restriction of the span map $\pi:\mathbb C^{n+1}\setminus\{0\}\to\mathbb{CP}^n$. It is surjective. If $\sigma$ is a local section of $\pi$ and $\nu(z)=z/|z|$, then $\bar\sigma=\nu\circ\sigma$ is a local section of $p$ because $p\circ\nu=\pi$, so $p$ is a submersion. The $S^1$-action $\lambda\cdot(z^1,\ldots,z^{n+1})=(\lambda z^1,\ldots,\lambda z^{n+1})$ on the round sphere is isometric, vertical, and fiber-transitive. Hence Theorem \ref{thm:group-action-riemannian-submersion} gives a unique metric on $\mathbb{CP}^n$ making $p$ a Riemannian submersion; this is the Fubini--Study metric.
\end{example}

\section{Riemannian Coverings}

\begin{definition}
\label{def:smooth-covering-map}
\lean{LeeLib.Ch02.IsSmoothCoveringMap}
\uses{prop:elementary-covering-maps}
\leanok



A smooth map $\pi:\tilde M\to M$ is a \emph{smooth covering map} if it is surjective, a topological covering map, and a local diffeomorphism.
\end{definition}

\begin{definition}
\label{def:covering-automorphism-group}
\lean{LeeLib.Ch02.CoveringAut}
\uses{def:smooth-covering-map}
\leanok



For a smooth covering $\pi:\tilde M\to M$, a \emph{covering automorphism} is a diffeomorphism $\varphi:\tilde M\to\tilde M$ with $\pi\circ\varphi=\pi$. These form $\operatorname{Aut}_\pi(\tilde M)$ under composition.
\end{definition}

\begin{definition}
\label{def:normal-covering-fibre-transitive}
\lean{LeeLib.Ch02.IsNormalCovering}
\uses{def:covering-automorphism-group, prop:normal-covering-transitive-action}
\leanok



The covering $\pi$ is \emph{normal} if $\operatorname{Aut}_\pi(\tilde M)$ acts transitively on every fiber. Equivalently, the image of $\pi_*:\pi_1(\tilde M,\tilde x)\to\pi_1(M,\pi(\tilde x))$ is normal (Proposition C.21).
\end{definition}

\begin{definition}
\label{def:covering-automorphism-invariant-metric}
\lean{LeeLib.Ch02.IsCoveringAutInvariant}
\uses{def:covering-automorphism-group, def:metric-preserving}
\leanok



A metric $\tilde g$ on $\tilde M$ is \emph{invariant under all covering automorphisms} if $\varphi^*\tilde g=\tilde g$ for every $\varphi\in\operatorname{Aut}_\pi(\tilde M)$.
\end{definition}

\begin{definition}
\label{def:riemannian-covering}
\lean{LeeLib.Ch02.IsRiemannianCovering}
\uses{def:smooth-covering-map, def:metric-preserving}
\leanok



A smooth covering $\pi:(\tilde M,\tilde g)\to(M,g)$ is a \emph{Riemannian covering} if $\pi^*g=\tilde g$.
\end{definition}

\begin{lemma}[A Riemannian Covering Is a Local Isometry]
\label{lem:riemannian-covering-local-isometry}
\lean{LeeLib.Ch02.IsRiemannianCovering.isLocalIsometry}
\uses{def:riemannian-covering, def:local-isometry, def:smooth-covering-map}
\leanok



Every Riemannian covering is a local isometry.
\end{lemma}

\begin{proof}
\leanok
\uses{def:local-isometry,def:smooth-covering-map}
\sketch A smooth covering is a local diffeomorphism, and $\pi^*g=\tilde g$ means that each local inverse chart preserves the metric.
\end{proof}

\begin{proposition}
\label{prop:equidimensional-riemannian-submersion-local-isometry}
\lean{LeeLib.Ch02.isRiemannianSubmersion_iff_isMetricPreserving}
\uses{def:riemannian-submersion, def:metric-preserving, def:submersion}
\leanok
\dcref{ch2:2.31}



If $\pi:\tilde M\to M$ is a smooth submersion with injective differential everywhere, then $\pi$ is a Riemannian submersion iff $\pi^*g=\tilde g$.
\end{proposition}

\begin{proof}
\leanok
\sketch Injectivity gives $V_x=0$ and hence $H_x=T_x\tilde M$, so the submersion identity is exactly $\pi^*g=\tilde g$. Conversely that identity implies the submersion identity on $H_x$.
\end{proof}

\begin{proposition}
\label{prop:riemannian-covering-metric}
\lean{LeeLib.Ch02.existsUnique_isRiemannianCovering_metric}
\uses{def:smooth-covering-map, def:normal-covering-fibre-transitive, def:covering-automorphism-invariant-metric, def:riemannian-covering, prop:elementary-covering-maps, prop:normal-covering-transitive-action, thm:group-action-riemannian-submersion, prop:equidimensional-riemannian-submersion-local-isometry}
\leanok
\dcref{ch2:2.31}



If $\pi:\tilde M\to M$ is a smooth normal covering and $\tilde g$ is invariant under all covering automorphisms, there is a unique metric $g$ on $M$ for which $\pi$ is a Riemannian covering.
\end{proposition}

\begin{proof}
\leanok
\sketch Proposition A.49 makes $\pi$ a surjective smooth submersion; normality gives fiber transitivity of the vertical automorphism action. Theorem \ref{thm:group-action-riemannian-submersion} gives a unique submersion metric, and injectivity of $d\pi$ together with Proposition \ref{prop:equidimensional-riemannian-submersion-local-isometry} identifies it with the unique covering metric.
\end{proof}

\begin{proposition}
\label{prop:riemannian-quotient-metric}
\uses{prop:quotient-orbit-space-covering-map, prop:riemannian-covering-metric}
\dcref{ch2:2.32}

If a discrete Lie group $\Gamma$ acts smoothly, freely, properly, and isometrically on $(\widetilde M,\tilde g)$, then $\widetilde M/\Gamma$ has a unique metric making the quotient map a normal Riemannian covering.
\end{proposition}

\begin{proof}
\sketch Proposition C.23 gives a smooth normal covering quotient map; Proposition \ref{prop:riemannian-covering-metric} supplies the unique metric.
\end{proof}

\begin{corollary}
\label{cor:isometric-to-quotient}
\uses{prop:covering-automorphisms-discrete-lie-group, prop:quotient-orbit-space-covering-map, prop:uniqueness-smooth-quotients}
\dcref{ch2:2.33}

If connected Riemannian manifolds $(M,g)$ and $(\widetilde M,\tilde g)$ are related by a normal Riemannian covering $\pi:\widetilde M\to M$ and $\Gamma=\operatorname{Aut}_\pi(\widetilde M)$, then $M$ is isometric to $\widetilde M/\Gamma$.
\end{corollary}

\begin{proof}
\sketch With the discrete topology, $\Gamma$ acts smoothly, freely, and properly; Proposition C.23 gives the quotient covering $q:\widetilde M\to\widetilde M/\Gamma$. Normality makes $\pi$ and $q$ constant on each other's fibers, so Proposition A.19 gives a diffeomorphism $F:M\to\widetilde M/\Gamma$ with $F\circ\pi=q$. Since both maps are local isometries, $F$ is a global isometry.
\end{proof}

\begin{example}
\label{ex:real-projective-space-metric}
\uses{ex:universal-covering-projective-space, prop:riemannian-quotient-metric}
\dcref{ch2:2.34}

The action of $\Gamma=\{\pm1\}$ on $S^n$ is smooth, free, proper, and isometric. Its quotient is $\mathbb{RP}^n$ and the quotient map is a smooth normal covering, so Proposition \ref{prop:riemannian-quotient-metric} gives the unique Riemannian covering metric on $\mathbb{RP}^n$.
\end{example}

\begin{example}[The Open Möbius Band]
\label{ex:open-mobius-band}
\uses{prop:riemannian-quotient-metric, REF:ch2-8}
\dcref{ch2:2.35}

The open Möbius band is $M=\mathbb R^2/\mathbb Z$ for the action
$n\cdot(x,y)=(x+n,(-1)^ny)$. This action is smooth, free, proper, and
isometric, so $M$ has a flat metric for which the quotient map is a Riemannian
covering.
\end{example}

\section{Basic Constructions on Riemannian Manifolds}

Throughout this section, $M$ is a smooth manifold, possibly with boundary.

\section{Raising and Lowering Indices}

Let $g$ be a Riemannian metric on $M$. The bundle map
\[
\widetilde g\colon TM\longrightarrow T^*M,
\qquad \widetilde g(v)(w)=g_p(v,w)
\]
for $v,w\in T_pM$ is smooth and satisfies $\widetilde g(X)(Y)=g(X,Y)$ for smooth
vector fields $X,Y$. If $(E_i)$ is a local frame with dual coframe $(e^i)$ and
$g=g_{ij}e^i\otimes e^j$, then
\[
\widetilde g(X)=(g_{ij}X^i)e^j=X_je^j,
\qquad X_j=g_{ij}X^i.
\]
This operation is called \emph{lowering an index}, and $\widetilde g(X)$ is denoted
$X^\flat$.

The inverse matrix $(g^{ij})$ is symmetric and satisfies
\[
g^{ij}g_{jk}=g_{kj}g^{ji}=\delta_k^i.
\]
Consequently,
\[
\widetilde g^{-1}(\omega)=\omega^iE_i,
\]
where
\begin{equation}
\omega^i=g^{ij}\omega_j. \tag{2.13}
\end{equation}
The inverse operation is \emph{raising an index}; $\widetilde g^{-1}(\omega)$ is denoted
$\omega^\sharp$. The mutually inverse maps $\flat$ and $\sharp$ are the musical
isomorphisms.

\begin{definition}[Lowering an index]
\label{def:flat-map}
\lean{LeeLib.Ch02.flat}
\uses{def:riemannian-metric}
\leanok



For a Riemannian manifold $(M,g)$ and $p\in M$, the \emph{flat map} is
\[
\flat\colon T_pM\longrightarrow T_p^*M,
\qquad v^\flat=g_p(v,\cdot).
\]
\end{definition}

\begin{lemma}[The flat map is a bijection]
\label{lem:flat-bijective}
\lean{LeeLib.Ch02.flat_bijective}
\uses{def:flat-map, def:riemannian-metric}
\leanok



Let $(M,g)$ be a finite-dimensional Riemannian manifold and $p\in M$. Then
$\flat\colon T_pM\to T_p^*M$ is bijective.
\end{lemma}

\begin{proof}
\leanok
\uses{def:flat-map,def:riemannian-metric}
\sketch The Riesz representation theorem identifies the finite-dimensional inner product space
$T_pM$ isometrically with its dual by $v\mapsto g_p(v,\cdot)$, which is the flat map.
\end{proof}

\begin{definition}[Raising an index]
\label{def:sharp-map}
\lean{LeeLib.Ch02.sharp}
\uses{def:flat-map, lem:flat-bijective}
\leanok



For a finite-dimensional Riemannian manifold $(M,g)$ and $p\in M$, the
\emph{sharp map} is
\[
\sharp=\flat^{-1}\colon T_p^*M\longrightarrow T_pM.
\]
Thus $\alpha^\sharp$ is the unique vector satisfying
\[
g_p(\alpha^\sharp,w)=\alpha(w)
\qquad (w\in T_pM).
\]
\end{definition}

For $f\in C^\infty(M)$, the gradient has the coordinate expression
\[
\operatorname{grad}f=(g^{ij}E_if)E_j
\]
and is characterized by
\begin{equation}
df_p(w)=\langle\operatorname{grad}f|_p,w\rangle
\qquad (p\in M,\ w\in T_pM). \tag{2.14}
\end{equation}
In an orthonormal frame, the components of $\operatorname{grad}f$ agree with those
of $df$; this need not hold in a general frame.

\begin{definition}[The gradient]
\label{def:gradient}
\lean{LeeLib.Ch02.grad}
\uses{def:sharp-map, def:riemannian-metric}
\leanok



Let $(M,g)$ be a finite-dimensional Riemannian manifold and $f\colon M\to\mathbb R$
smooth. The \emph{gradient} of $f$ is
\[
\operatorname{grad}f=(df)^\sharp,
\qquad \operatorname{grad}f|_p=(df_p)^\sharp\in T_pM.
\]
\end{definition}

\begin{lemma}[Characterization of the gradient]
\label{lem:gradient-characterization}
\lean{LeeLib.Ch02.innerAt_grad}
\uses{def:gradient, def:sharp-map}
\leanok



Let $(M,g)$ be a finite-dimensional Riemannian manifold and $f\colon M\to\mathbb R$
smooth. For every $p\in M$ and $w\in T_pM$,
\[
df_p(w)=\langle\operatorname{grad}f|_p,w\rangle.
\]
Moreover, $\operatorname{grad}f|_p$ is the unique vector with this property.
\end{lemma}

\begin{proof}
\leanok
\uses{def:gradient,def:sharp-map,lem:flat-bijective}
\sketch Apply the defining identity for $\sharp$ to $\alpha=df_p$. If $v$ satisfies the
same identity, then $v^\flat=(\operatorname{grad}f|_p)^\flat$, so injectivity of
$\flat$ gives $v=\operatorname{grad}f|_p$.
\end{proof}

\begin{lemma}[The gradient is orthogonal to the kernel of the differential]
\label{lem:gradient-orthogonal-kernel}
\lean{LeeLib.Ch02.innerAt_grad_eq_zero_of_mem_ker}
\uses{def:gradient, lem:gradient-characterization}
\leanok



Let $(M,g)$ be a finite-dimensional Riemannian manifold, $f\colon M\to\mathbb R$
smooth, and $p\in M$. If $w\in T_pM$ satisfies $df_p(w)=0$, then
$\langle\operatorname{grad}f|_p,w\rangle=0$.
\end{lemma}

\begin{proof}
\leanok
\uses{lem:gradient-characterization}
\sketch Equation (2.14) identifies the two quantities.
\end{proof}

\begin{lemma}[Critical points are the zeros of the gradient]
\label{lem:gradient-eq-zero-iff}
\lean{LeeLib.Ch02.grad_eq_zero_iff}
\uses{def:gradient, lem:gradient-characterization, lem:flat-bijective}
\leanok



Let $(M,g)$ be a finite-dimensional Riemannian manifold, $f\colon M\to\mathbb R$
smooth, and $p\in M$. Then
\[
\operatorname{grad}f|_p=0\quad\Longleftrightarrow\quad df_p=0.
\]
Thus $p$ is regular precisely when $\operatorname{grad}f|_p\ne0$.
\end{lemma}

\begin{proof}
\leanok
\uses{lem:gradient-characterization,lem:flat-bijective}
\sketch If the gradient vanishes, (2.14) gives $df_p(w)=0$ for every $w$. Conversely,
$df_p=0$ makes $\operatorname{grad}f|_p$ orthogonal to every vector; equivalently,
its flat image is zero, and injectivity of $\flat$ proves that the gradient vanishes.
\end{proof}

\begin{lemma}[Orthogonality to the gradient characterizes the kernel of the differential]
\label{lem:gradient-orthogonal-iff}
\lean{LeeLib.Ch02.innerAt_grad_eq_zero_iff}
\uses{def:gradient, lem:gradient-characterization}
\leanok



Let $(M,g)$ be a finite-dimensional Riemannian manifold, $f\colon M\to\mathbb R$
smooth, $p\in M$, and $w\in T_pM$. Then
\[
\langle\operatorname{grad}f|_p,w\rangle=0
\quad\Longleftrightarrow\quad df_p(w)=0.
\]
\end{lemma}

\begin{proof}
\leanok
\uses{lem:gradient-characterization}
\sketch This is immediate from (2.14).
\end{proof}

\begin{definition}[The regular set]
\label{def:regular-set}
\lean{LeeLib.Ch02.regularSet}
\leanok


Let $M$ be a smooth manifold and $f\colon M\to\mathbb R$ smooth. The
\emph{regular set} of $f$ is
\[
\mathcal R=\{x\in M:df_x\ne0\}.
\]
Its points are \emph{regular}; the remaining points are \emph{critical}.
\end{definition}

\begin{lemma}[The regular set is where the gradient does not vanish]
\label{lem:regular-set-eq-gradient-nonvanishing}
\lean{LeeLib.Ch02.mem_regularSet_iff_grad_ne_zero}
\uses{def:regular-set, def:gradient, lem:gradient-eq-zero-iff}
\leanok



Let $(M,g)$ be a finite-dimensional Riemannian manifold and $f\colon M\to\mathbb R$
smooth. Then
\[
x\in\mathcal R\quad\Longleftrightarrow\quad\operatorname{grad}f|_x\ne0.
\]
\end{lemma}

\begin{proof}
\leanok
\uses{lem:gradient-eq-zero-iff}
\sketch Negate the equivalence in Lemma \ref{lem:gradient-eq-zero-iff}.
\end{proof}

\begin{lemma}[The non-vanishing locus of a differential is open]
\label{lem:mfderiv-nonvanishing-open}
\lean{LeeLib.Ch02.isOpen_setOf_mfderiv_ne_zero}
\leanok


Let $F\colon M\to N$ be a smooth map of smooth manifolds. Then
$\{x\in M:dF_x\ne0\}$ is open in $M$.
\end{lemma}

\begin{proof}
\leanok
\sketch For $x_0$ with $dF_{x_0}\ne0$, choose tangent-bundle trivializations $\Phi$
near $x_0$ and $\Psi$ near $F(x_0)$. On their common domain set
\[
D(x)=\Psi_{F(x)}\circ dF_x\circ\Phi_x^{-1}\in\operatorname{Hom}(E,E').
\]
The map $D$ is continuous, and the trivializations imply $D(x)=0$ exactly when
$dF_x=0$. Since the complement of zero in $\operatorname{Hom}(E,E')$ is open,
$dF$ remains nonzero near $x_0$.
\end{proof}

\begin{lemma}[The regular set is open]
\label{lem:regular-set-open}
\lean{LeeLib.Ch02.isOpen_regularSet}
\uses{def:regular-set, lem:mfderiv-nonvanishing-open}
\leanok



Let $M$ be a smooth manifold and $f\colon M\to\mathbb R$ smooth. Then the regular
set $\mathcal R$ of $f$ is open.
\end{lemma}

\begin{proof}
\leanok
\uses{lem:mfderiv-nonvanishing-open}
\sketch Apply Lemma \ref{lem:mfderiv-nonvanishing-open} to $F=f$.
\end{proof}

A level $f^{-1}(c)$ is regular when $df_x\ne0$ for every $x\in f^{-1}(c)$,
equivalently, when $c$ is a regular value of $f$.

\begin{proposition}[The gradient is normal to a regular level set]
\label{prop:gradient-normal-regular-level-set}
\lean{LeeLib.Ch02.innerAt_grad_mfderiv_levelSet_val_eq_zero}
\uses{prop:regular-level-set-hypersurface-codim-one, prop:regular-level-set-tangent-kernel, def:gradient, lem:gradient-orthogonal-kernel}
\leanok



Let $(M,g)$ be an $(n+1)$-dimensional Riemannian manifold without boundary,
$f\in C^\infty(M)$, and $c$ a regular value of $f$. If
$\iota_c\colon f^{-1}(c)\hookrightarrow M$ is the inclusion, then for every
$x_0\in f^{-1}(c)$ and $w\in T_{x_0}f^{-1}(c)$,
\[
\left\langle\operatorname{grad}f|_{x_0},d(\iota_c)_{x_0}(w)\right\rangle=0.
\]
Thus $\operatorname{grad}f$ is normal to $f^{-1}(c)$.
\end{proposition}

\begin{proof}
\leanok
\uses{prop:regular-level-set-tangent-kernel,lem:gradient-orthogonal-kernel}
\sketch Since $f\circ\iota_c$ is constant, the chain rule places
$d(\iota_c)_{x_0}(w)$ in $\ker df_{x_0}$. Apply
Lemma \ref{lem:gradient-orthogonal-kernel}.
\end{proof}

\begin{proposition}[The tangent space of a regular level set is the orthogonal complement of the gradient]
\label{prop:tangent-space-regular-level-set-gradient-orthogonal}
\lean{LeeLib.Ch02.exists_mfderiv_levelSet_val_eq_of_innerAt_grad_eq_zero}
\uses{prop:regular-level-set-tangent-kernel, lem:gradient-orthogonal-iff, def:gradient}
\leanok



In the setting of Proposition \ref{prop:gradient-normal-regular-level-set}, every
$u\in T_{x_0}M$ satisfying
$\langle\operatorname{grad}f|_{x_0},u\rangle=0$ belongs to the image of
$d(\iota_c)_{x_0}$. Hence
\[
d(\iota_c)_{x_0}\bigl(T_{x_0}f^{-1}(c)\bigr)
=\bigl(\operatorname{grad}f|_{x_0}\bigr)^\perp,
\]
and the normal space is the line spanned by $\operatorname{grad}f|_{x_0}$.
\end{proposition}

\begin{proof}
\leanok
\uses{prop:regular-level-set-tangent-kernel,lem:gradient-orthogonal-iff}
\sketch Orthogonality gives $u\in\ker df_{x_0}$ by
Lemma \ref{lem:gradient-orthogonal-iff}, and Proposition
\ref{prop:regular-level-set-tangent-kernel} identifies this kernel with the image
of $d(\iota_c)_{x_0}$. Proposition
\ref{prop:gradient-normal-regular-level-set} supplies the reverse inclusion.
Because this image is a hyperplane and the gradient is nonzero, its orthogonal
complement is the gradient line.
\end{proof}

\subsection*{Open submanifolds and restriction}

For an open submanifold $U\subseteq M$, tangent spaces are identified using the
restricted charts.

\begin{lemma}[The differential of an open-submanifold inclusion]
\label{lem:opens-inclusion-differential}
\lean{LeeLib.Ch02.mfderiv_opens_subtypeVal}
\leanok


Let $M$ be a smooth manifold and $U\subseteq M$ open. Under the canonical
identification $T_xU=T_xM$, the inclusion $\iota\colon U\hookrightarrow M$ satisfies
\[
d\iota_x=\operatorname{id}\colon T_xU\longrightarrow T_xM.
\]
In particular, $d\iota_x$ is an isomorphism.
\end{lemma}

\begin{proof}
\leanok
\sketch A chart for $U$ is the restriction of a chart for $M$, so the coordinate
representation of $\iota$ is the identity. Its differential is therefore the identity.
\end{proof}

\begin{lemma}[Restriction to an open set preserves the differential]
\label{lem:opens-restrict-differential}
\lean{LeeLib.Ch02.mfderiv_opens_restrict}
\uses{lem:opens-inclusion-differential}
\leanok



Let $F\colon M\to M'$ be smooth, $U\subseteq M$ open, and $x\in U$. Then
\[
d(F|_U)_x=dF_x.
\]
\end{lemma}

\begin{proof}
\leanok
\uses{lem:opens-inclusion-differential}
\sketch Write $F|_U=F\circ\iota$ and use the chain rule together with
$d\iota_x=\operatorname{id}$.
\end{proof}

\begin{lemma}[Every value is a regular value of $f|_{\mathcal{R}}$]
\label{lem:regular-restrict-regular-value}
\lean{LeeLib.Ch02.regularRestrict_regular}
\uses{def:regular-set, lem:regular-set-open, lem:opens-restrict-differential}
\leanok



Let $M$ be a smooth manifold, $f\in C^\infty(M)$, and regard its regular set
$\mathcal R$ as an open submanifold. Then
\[
d(f|_{\mathcal R})_x\ne0\qquad(x\in\mathcal R).
\]
Thus every $c\in\mathbb R$ is a regular value of $f|_{\mathcal R}$.
\end{lemma}

\begin{proof}
\leanok
\uses{lem:opens-restrict-differential,def:regular-set}
\sketch Lemma \ref{lem:opens-restrict-differential} gives
$d(f|_{\mathcal R})_x=df_x$, which is nonzero by the definition of $\mathcal R$.
\end{proof}

\subsection*{The regular part of an arbitrary level set}

Write $M_c=(f|_{\mathcal R})^{-1}(c)$. Its image in $M$ is
$f^{-1}(c)\cap\mathcal R$.

\begin{proposition}
\label{prop:regular-level-set-hypersurface}
\lean{LeeLib.Ch02.isManifold_regularLevelSet}
\uses{prop:regular-level-set-hypersurface-codim-one, def:regular-set, lem:regular-set-open, lem:regular-restrict-regular-value}
\leanok


\dcref{ch2:2.37}

Let $M$ be a smooth $(n+1)$-manifold without boundary, let
$f\in C^\infty(M)$, and let $\mathcal R$ be the regular set of $f$. For every
$c\in\mathbb R$, the set
\[
M_c=f^{-1}(c)\cap\mathcal R
\]
is a smooth $n$-manifold.
\end{proposition}

\begin{proof}
\leanok
\uses{prop:regular-level-set-hypersurface-codim-one,lem:regular-restrict-regular-value,lem:regular-set-open}
\sketch The open set $\mathcal R$ is an $(n+1)$-manifold, and $c$ is a regular value
of $f|_{\mathcal R}$. Apply Proposition
\ref{prop:regular-level-set-hypersurface-codim-one} to obtain the smooth
$n$-manifold $(f|_{\mathcal R})^{-1}(c)$.
\end{proof}

\begin{lemma}[The regular part, as a subset of $M$]
\label{lem:regular-part-image}
\lean{LeeLib.Ch02.range_regularLevelSetIncl}
\uses{prop:regular-level-set-hypersurface, def:regular-set}
\leanok



In the setting of Proposition \ref{prop:regular-level-set-hypersurface}, the image
of $M_c=(f|_{\mathcal R})^{-1}(c)$ under $M_c\hookrightarrow\mathcal R\hookrightarrow M$
is exactly $f^{-1}(c)\cap\mathcal R$.
\end{lemma}

\begin{proof}
\leanok
\sketch A point is in the image precisely when it lies in $\mathcal R$ and has
$f$-value $c$.
\end{proof}

\begin{lemma}[The regular part is embedded]
\label{lem:regular-part-embedded}
\lean{LeeLib.Ch02.isEmbedding_regularLevelSetIncl}
\uses{prop:regular-level-set-hypersurface}
\leanok



In the setting of Proposition \ref{prop:regular-level-set-hypersurface}, the inclusion
$M_c\hookrightarrow M$ is a topological embedding.
\end{lemma}

\begin{proof}
\leanok
\sketch Both $M_c\hookrightarrow\mathcal R$ and $\mathcal R\hookrightarrow M$ are
subspace inclusions and topological embeddings; so is their composite.
\end{proof}

\begin{lemma}[The inclusion of the regular part is smooth]
\label{lem:regular-part-inclusion-smooth}
\lean{LeeLib.Ch02.contMDiff_regularLevelSetIncl}
\uses{prop:regular-level-set-hypersurface, prop:regular-level-set-inclusion-smooth}
\leanok



In the setting of Proposition \ref{prop:regular-level-set-hypersurface}, the inclusion
$M_c\hookrightarrow M$ is smooth.
\end{lemma}

\begin{proof}
\leanok
\uses{prop:regular-level-set-inclusion-smooth}
\sketch Compose the smooth regular-level-set inclusion $M_c\hookrightarrow\mathcal R$
with the smooth open-submanifold inclusion $\mathcal R\hookrightarrow M$.
\end{proof}

\begin{lemma}[The inclusion of the regular part is an immersion]
\label{lem:regular-part-inclusion-immersion}
\lean{LeeLib.Ch02.injective_mfderiv_regularLevelSetIncl}
\uses{prop:regular-level-set-hypersurface, lem:opens-inclusion-differential, prop:regular-level-set-inclusion-smooth}
\leanok



In the setting of Proposition \ref{prop:regular-level-set-hypersurface}, the
differential of $M_c\hookrightarrow M$ is injective at every point. Together with
Lemma \ref{lem:regular-part-embedded}, this makes $M_c$ an embedded smooth
hypersurface in $M$.
\end{lemma}

\begin{proof}
\leanok
\uses{lem:opens-inclusion-differential,prop:regular-level-set-inclusion-smooth}
\sketch The chain rule factors this differential through $M_c\hookrightarrow\mathcal R$
and $\mathcal R\hookrightarrow M$. The second factor has identity differential,
while the first has injective differential by Proposition
\ref{prop:regular-level-set-inclusion-smooth}.
\end{proof}

\begin{lemma}[The tangent space of the regular part]
\label{lem:regular-part-tangent-kernel}
\lean{LeeLib.Ch02.range_mfderiv_regularLevelSetIncl}
\uses{prop:regular-level-set-hypersurface, prop:regular-level-set-tangent-kernel, lem:opens-restrict-differential}
\leanok



In the setting of Proposition \ref{prop:regular-level-set-hypersurface}, for each
$x\in M_c$ the image of the differential of $M_c\hookrightarrow M$ equals
$\ker df_x$; equivalently,
\[
T_xM_c=\ker df_x.
\]
\end{lemma}

\begin{proof}
\leanok
\uses{prop:regular-level-set-tangent-kernel,lem:opens-restrict-differential,lem:opens-inclusion-differential}
\sketch The open-submanifold inclusion has identity differential. Proposition
\ref{prop:regular-level-set-tangent-kernel} identifies the image in $T_x\mathcal R$
with $\ker d(f|_{\mathcal R})_x$, and Lemma
\ref{lem:opens-restrict-differential} identifies this kernel with $\ker df_x$.
\end{proof}

\begin{proposition}[$\operatorname{grad} f$ is normal to the regular part]
\label{prop:gradient-normal-regular-part}
\lean{LeeLib.Ch02.innerAt_grad_mfderiv_regularLevelSetIncl_eq_zero}
\uses{prop:regular-level-set-hypersurface, lem:regular-part-tangent-kernel, def:gradient, lem:gradient-orthogonal-kernel}
\leanok



Let $(M,g)$ be an $(n+1)$-dimensional Riemannian manifold without boundary,
$f\in C^\infty(M)$, and $c\in\mathbb R$. If $\iota\colon M_c\hookrightarrow M$ is
the inclusion, then for all $x\in M_c$ and $w\in T_xM_c$,
\[
\left\langle\operatorname{grad}f|_x,d\iota_x(w)\right\rangle=0.
\]
Thus $\operatorname{grad}f$ is normal to $M_c$.
\end{proposition}

\begin{proof}
\leanok
\uses{lem:regular-part-tangent-kernel,lem:gradient-orthogonal-kernel}
\sketch Constancy of $f$ on $M_c$ and the chain rule give
$d\iota_x(w)\in\ker df_x$. Apply Lemma
\ref{lem:gradient-orthogonal-kernel}.
\end{proof}

\begin{lemma}[$\operatorname{grad} f$ does not vanish on the regular part]
\label{lem:gradient-nonvanishing-regular-part}
\lean{LeeLib.Ch02.grad_ne_zero_regularLevelSet}
\uses{prop:regular-level-set-hypersurface, lem:regular-set-eq-gradient-nonvanishing, def:gradient}
\leanok



In the setting of Proposition \ref{prop:gradient-normal-regular-part},
$\operatorname{grad}f|_x\ne0$ for every $x\in M_c$. It therefore spans the normal
line to the hypersurface $M_c$ at $x$.
\end{lemma}

\begin{proof}
\leanok
\uses{lem:regular-set-eq-gradient-nonvanishing}
\sketch Every point of $M_c$ belongs to $\mathcal R$, so nonvanishing follows from
Lemma \ref{lem:regular-set-eq-gradient-nonvanishing}. Proposition
\ref{prop:gradient-normal-regular-part} and Lemma
\ref{lem:regular-part-tangent-kernel} then identify the normal line.
\end{proof}

Musical maps act on any tensor index. If the $i$th argument of a $(k,l)$-tensor
$F$ is a vector, raising that covariant index gives a $(k+1,l-1)$-tensor defined by
\[
F^\sharp(\alpha_1,\ldots,\alpha_{k+l})
=F(\alpha_1,\ldots,\alpha_{i-1},\alpha_i^\sharp,
   \alpha_{i+1},\ldots,\alpha_{k+l}).
\]
Its components are obtained by contracting the $i$th index with $g^{pq}$. If the
$i$th argument is a covector, lowering that contravariant index gives a
$(k-1,l+1)$-tensor
\[
F^\flat(\alpha_1,\ldots,\alpha_{k+l})
=F(\alpha_1,\ldots,\alpha_{i-1},\alpha_i^\flat,
   \alpha_{i+1},\ldots,\alpha_{k+l}),
\]
whose components are obtained by contraction with $g_{pq}$. When several indices
are eligible, the intended index is specified explicitly.

For example, if
\[
A=A_i{}^j{}_k\,\varepsilon^i\otimes E_j\otimes\varepsilon^k,
\]
then lowering the middle index produces a covariant $3$-tensor with components
\[
A_{ijk}=g_{jl}A_i{}^l{}_k.
\]

For a covariant $k$-tensor $h$ with $k\ge2$, raise its last index and define its
metric trace by
\[
\operatorname{tr}_gh=\operatorname{tr}(h^\sharp).
\]
This is a covariant $(k-2)$-tensor. Other raised or contracted index pairs are
specified when used. For a covariant $2$-tensor,
\[
\operatorname{tr}_gh=h_i{}^i=g^{ij}h_{ij}.
\]
In an orthonormal frame this equals the ordinary matrix trace; in a general frame
it need not.

\begin{proposition}
\label{prop:inverse-metric-ambiguity}


\dcref{ch2:2.38}
\uses{def:dual-coframe,def:frame-matrix,def:sharp-map}
Let $g$ be a Riemannian metric and $(E_i)$ a local frame. The components of
the contravariant tensor $g^{\sharp\sharp}$ are exactly the entries $(g^{ij})$
of the inverse matrix to $(g_{ij})$.
\end{proposition}

\section{Inner Products of Tensors}

The metric induces an inner product on $T_x^*M$ by
\[
\langle\omega,\eta\rangle_g
=\langle\omega^\sharp,\eta^\sharp\rangle_g.
\]
In a local frame,
\begin{align*}
\langle\omega,\eta\rangle
&=g_{kl}(g^{ki}\omega_i)(g^{lj}\eta_j)\\
&=\delta_l^i g^{lj}\omega_i\eta_j\\
&=g^{ij}\omega_i\eta_j.
\end{align*}
Equivalently,
\[
\langle\omega,\eta\rangle=\omega_i\eta^i=\omega^j\eta_j.
\]

\begin{proposition}
\label{prop:orthonormal-frames-and-coframes}


\dcref{ch2:2.39}
\uses{prop:inverse-metric-ambiguity,def:dual-coframe,def:sharp-map,def:frame-matrix}
Let $(M,g)$ be a Riemannian manifold, possibly with boundary, let $(E_i)$ be a
local frame, and let $(e^i)$ be its dual coframe. The following are
equivalent:
\begin{enumerate}
\item[(a)] $(E_i)$ is orthonormal.
\item[(b)] $(e^i)$ is orthonormal.
\item[(c)] $(e^i)^\sharp=E_i$ for every $i$.
\end{enumerate}
\end{proposition}

A \emph{smooth fiber metric} on a smooth vector bundle $E\to M$ is a fiberwise
inner product for which $\langle\sigma,\tau\rangle$ is smooth whenever $\sigma$ and
$\tau$ are local smooth sections.

\begin{proposition}[Inner Products of Tensors]
\label{prop:inner-products-of-tensors}
\dcref{ch2:2.40}
Let $(M,g)$ be an $n$-dimensional Riemannian manifold, possibly with boundary.
Each tensor bundle $T^{(k,l)}TM$ has a unique smooth fiber metric such that, for
vector or covector fields $\alpha_i,\beta_i$ of the appropriate types,
\begin{equation}
\left\langle\alpha_1\otimes\cdots\otimes\alpha_{k+l},
\beta_1\otimes\cdots\otimes\beta_{k+l}\right\rangle
=\prod_{a=1}^{k+l}\langle\alpha_a,\beta_a\rangle. \tag{2.15}
\end{equation}
If $(E_1,\ldots,E_n)$ is a local orthonormal frame and
$(e^1,\ldots,e^n)$ its dual coframe, then
\[
E_{i_1}\otimes\cdots\otimes E_{i_k}\otimes
e^{j_1}\otimes\cdots\otimes e^{j_l}
\]
over all indices is a local orthonormal frame for $T^{(k,l)}(T_pM)$. In an
arbitrary local frame,
\begin{equation}
\langle F,G\rangle
=g_{i_1r_1}\cdots g_{i_kr_k}g^{j_1s_1}\cdots g^{j_ls_l}
F_{j_1\cdots j_l}^{i_1\cdots i_k}
G_{s_1\cdots s_l}^{r_1\cdots r_k}. \tag{2.16}
\end{equation}
For covariant tensors $F,G$ this becomes
\[
\langle F,G\rangle=F_{j_1\cdots j_l}G^{j_1\cdots j_l},
\]
where
\[
G^{j_1\cdots j_l}
=g^{j_1s_1}\cdots g^{j_ls_l}G_{s_1\cdots s_l}.
\]
\end{proposition}
\section{The Volume Form and Integration}

\begin{definition}[Smooth Orientation]
\label{def:smooth-orientation}
\lean{LeeLib.Ch02.IsSmoothOrientation}
\leanok


A pointwise orientation of a smooth $n$-manifold $M$ assigns an orientation $o_x$ to each $T_xM$. It is \emph{smooth} if every point has a neighbourhood carrying a smooth local frame positively oriented for $o_x$ at every point. An oriented Riemannian manifold is a Riemannian manifold with a smooth orientation.
\end{definition}

\begin{lemma}[Gram--Schmidt Preserves Orientation]
\label{lem:gram-schmidt-preserves-orientation}
\lean{InnerProductSpace.gramSchmidtOrthonormalBasis_orientation}
\leanok


If $b=(b_1,\ldots,b_n)$ is a basis of an $n$-dimensional real inner product space, its Gram--Schmidt orthonormalization has the same orientation.
\end{lemma}

\begin{proof}
\leanok
\sketch Let $u_i$ be the unnormalized Gram--Schmidt vectors and $\widehat u_i=u_i/|u_i|$. The change-of-basis matrix is triangular with determinant $\prod_i\langle\widehat u_i,b_i\rangle=\prod_i|u_i|>0$, since $\langle u_i,b_i\rangle=|u_i|^2$. Thus the orientation is unchanged.
\end{proof}

\begin{proposition}[Existence of Oriented Orthonormal Frames]
\label{prop:oriented-orthonormal-frames}
\lean{LeeLib.Ch02.exists_orientedOrthonormalFrame_nhds}
\uses{def:smooth-orientation, lem:gram-schmidt-preserves-orientation, prop:existence-orthonormal-frames, def:riemannian-metric}
\leanok



If $(M,g)$ is a Riemannian $n$-manifold with smooth orientation $o$, every $p\in M$ has a neighbourhood with a smooth local frame that is orthonormal and positively oriented at every point.
\end{proposition}

\begin{proof}
\leanok
\uses{def:smooth-orientation,lem:gram-schmidt-preserves-orientation,prop:existence-orthonormal-frames}
\sketch Start with a smooth positively oriented frame supplied by Definition \ref{def:smooth-orientation}. Fibrewise Gram--Schmidt gives a smooth orthonormal frame by Proposition \ref{prop:existence-orthonormal-frames}, and Lemma \ref{lem:gram-schmidt-preserves-orientation} preserves its orientation.
\end{proof}

\begin{definition}[The Riemannian Volume Form]
\label{def:volume-form-pointwise}
\lean{LeeLib.Ch02.RiemannianMetric.volumeForm}
\uses{def:riemannian-metric, def:smooth-orientation}
\leanok



For a Riemannian $n$-manifold $(M,g)$ with pointwise orientation $o$, define $dV_g|_x$ to be the unique alternating $n$-form on $(T_xM,g_x,o_x)$ taking value $1$ on every positively oriented orthonormal basis. The family $dV_g:x\mapsto dV_g|_x$ is the pointwise Riemannian volume form.
\end{definition}

\begin{proposition}[The Determinant Formula]
\label{prop:volume-form-determinant}
\lean{LeeLib.Ch02.RiemannianMetric.volumeForm_apply_eq_det}
\uses{def:volume-form-pointwise, prop:oriented-orthonormal-frames}
\leanok



If $(E_1,\ldots,E_n)$ is a local orthonormal frame positively oriented at $x$ and $v_1,\ldots,v_n\in T_xM$, then
\[
dV_g|_x(v_1,\ldots,v_n)=\det\bigl(\langle E_i|_x,v_j\rangle_g\bigr)_{i,j=1}^n.
\]
\end{proposition}

\begin{proof}
\leanok
\uses{def:volume-form-pointwise}
\sketch The $E_i|_x$ form a positively oriented orthonormal basis, so $dV_g|_x$ is its determinant form; the coefficient of $v_j$ along $E_i|_x$ is $\langle E_i|_x,v_j\rangle_g$.
\end{proof}

\begin{proposition}[The Characterizing Property of $dV_g$]
\label{prop:volume-form-value-one}
\lean{LeeLib.Ch02.RiemannianMetric.volumeForm_apply_eq_one}
\uses{def:volume-form-pointwise, prop:oriented-orthonormal-frames}
\leanok



If $(E_1,\ldots,E_n)$ is local orthonormal and positively oriented at $x$, then $dV_g|_x(E_1|_x,\ldots,E_n|_x)=1$.
\end{proposition}

\begin{proof}
\leanok
\uses{def:volume-form-pointwise}
\sketch This is the defining normalization of the volume form; equivalently, apply Proposition \ref{prop:volume-form-determinant} to the identity matrix.
\end{proof}

\begin{proposition}[Uniqueness of $dV_g$]
\label{prop:volume-form-uniqueness}
\lean{LeeLib.Ch02.RiemannianMetric.volumeForm_unique}
\uses{def:volume-form-pointwise, prop:volume-form-value-one}
\leanok



If $\theta$ is an alternating $n$-form on $T_xM$ taking value $1$ on a positively oriented orthonormal basis $(E_i|_x)$, then $\theta=dV_g|_x$.
\end{proposition}

\begin{proof}
\leanok
\uses{prop:volume-form-value-one}
\sketch Every top-degree form is determined by its value on one basis. Both $\theta$ and $dV_g|_x$ equal the determinant form of $(E_i|_x)$ because both take value $1$ there.
\end{proof}

\begin{lemma}[Top-Degree Forms are Proportional]
\label{lem:top-forms-proportional}
\lean{ContinuousAlternatingMap.eq_smul_of_apply_basis_ne_zero}
\leanok


Let $b=(b_1,\ldots,b_n)$ be a basis of an $n$-dimensional real vector space and let $\omega,\theta$ be alternating $n$-forms with $\theta(b_1,\ldots,b_n)\ne0$. Then
\[
\omega=\frac{\omega(b_1,\ldots,b_n)}{\theta(b_1,\ldots,b_n)}\,\theta.
\]
\end{lemma}

\begin{proof}
\leanok
\sketch Express both forms as their value on $b$ times the determinant form of $b$ and divide by the nonzero value of $\theta$.
\end{proof}

\begin{definition}[The Wedge Product of Covectors]
\label{def:wedge-covectors}
\lean{LeeLib.Ch02.wedgeCovectors}
\leanok


For continuous linear functionals $f_1,\ldots,f_n$ on a real topological vector space, define
\[
(f_1\wedge\cdots\wedge f_n)(v_1,\ldots,v_n)=\det\bigl(f_i(v_j)\bigr)_{i,j=1}^n.
\]
This determinant convention yields a continuous alternating $n$-form.
\end{definition}

\begin{proposition}[The Coframe Characterization of $dV_g$]
\label{prop:volume-form-coframe}
\lean{LeeLib.Ch02.RiemannianMetric.volumeForm_eq_wedgeCovectors}
\uses{def:wedge-covectors, def:volume-form-pointwise, prop:volume-form-determinant, prop:oriented-orthonormal-frames}
\leanok



If $(E_1,\ldots,E_n)$ is a positively oriented local orthonormal frame at $x$ and $e^i=\langle E_i,\cdot\rangle_g$ is its dual coframe, then
\[
dV_g|_x=e^1\wedge\cdots\wedge e^n.
\]
\end{proposition}

\begin{proof}
\leanok
\uses{def:wedge-covectors,prop:volume-form-determinant}
\sketch Evaluate both forms on arbitrary $(v_j)$. The wedge evaluates to $\det(e^i(v_j))=\det(\langle E_i|_x,v_j\rangle_g)$, which is the determinant formula for $dV_g|_x$.
\end{proof}

\begin{proposition}[The Riemannian Volume Form is Smooth]
\label{prop:volume-form-smooth}
\lean{LeeLib.Ch02.RiemannianMetric.contMDiff_volumeForm}
\uses{def:volume-form-pointwise, def:smooth-orientation, prop:oriented-orthonormal-frames, lem:top-forms-proportional, prop:volume-form-determinant, prop:volume-form-value-one, lem:inner-product-vector-fields-smooth}
\leanok



For a Riemannian $n$-manifold with smooth orientation, the family $dV_g$ is a smooth $n$-form.
\end{proposition}

\begin{proof}
\leanok
\uses{prop:oriented-orthonormal-frames,lem:top-forms-proportional,prop:volume-form-determinant,prop:volume-form-value-one,lem:inner-product-vector-fields-smooth}
\sketch In a tangent-bundle trivialization, read $dV_g|_x$ as an $n$-form $\Phi(x)$ on $\mathbb R^n$. Fix a basis $b$ and a local oriented orthonormal frame $Z_i$. Since $\Phi(x_0)\ne0$, its value on the basis $b$ is nonzero. Lemma \ref{lem:top-forms-proportional} expresses $\Phi(x)$ as a scalar multiple of $\Phi(x_0)$, with scalar $\Phi(x)(b)/\Phi(x_0)(b)$. Proposition \ref{prop:volume-form-determinant} writes the numerator as the determinant of the smooth pairings $\langle Z_i,e^{-1}b_j\rangle_g$, hence the scalar and $\Phi$ are smooth.
\end{proof}

\begin{definition}[The Matrix of $g$ in a Frame]
\label{def:frame-matrix}
\lean{LeeLib.Ch02.RiemannianMetric.frameMatrix}
\leanok


For a local frame $(Y_1,\ldots,Y_n)$ on $u$, the matrix of $g$ at $x\in u$ is
\[
g_{ij}(x)=\langle Y_i|_x,Y_j|_x\rangle_g.
\]
For a coordinate frame this is the coordinate matrix $(g_{ij})$.
\end{definition}

\begin{lemma}[The Gram Determinant of a Frame is Nonzero]
\label{lem:frame-matrix-nonsingular}
\lean{LeeLib.Ch02.RiemannianMetric.frameMatrix_det_ne_zero}
\uses{def:frame-matrix}
\leanok



For every $x\in u$, the Gram determinant $\det(g_{ij}(x))$ of a local frame is nonzero.
\end{lemma}

\begin{proof}
\leanok
\uses{def:frame-matrix}
\sketch The Gram matrix is nonsingular exactly when its vectors are linearly independent; frame values are a basis of $T_xM$.
\end{proof}

\begin{definition}[The Dual Coframe of a Frame]
\label{def:dual-coframe}
\lean{LeeLib.Ch02.RiemannianMetric.dualCoframe}
\uses{def:frame-matrix, lem:frame-matrix-nonsingular}
\leanok



Let $(g^{ij}(x))=(g_{ij}(x))^{-1}$. The dual coframe associated to a local frame $(Y_i)$ is
\[
\varepsilon^i|_x=\sum_j g^{ij}(x)\langle Y_j|_x,\cdot\rangle_g.
\]
For a coordinate frame it is $(dx^i)$.
\end{definition}

\begin{lemma}[The Dual Coframe is Dual]
\label{lem:dual-coframe-duality}
\lean{LeeLib.Ch02.RiemannianMetric.dualCoframe_apply_frame}
\uses{def:dual-coframe, def:frame-matrix}
\leanok



The coframe of Definition \ref{def:dual-coframe} satisfies $\varepsilon^i|_x(Y_k|_x)=\delta^i_k$.
\end{lemma}

\begin{proof}
\leanok
\uses{def:dual-coframe,def:frame-matrix,lem:frame-matrix-nonsingular}
\sketch The value is $\sum_jg^{ij}(x)g_{jk}(x)$, the $(i,k)$ entry of an inverse matrix product.
\end{proof}

\begin{lemma}[The Gram Determinant Formula, Pointwise]
\label{lem:volume-form-gram-determinant}
\lean{LeeLib.Ch02.volumeFormL_apply_eq_sqrt_det_gram}
\uses{lem:gram-schmidt-preserves-orientation}
\leanok



If $(b_1,\ldots,b_n)$ is a positively oriented basis of an oriented inner product space with volume form $dV$, then
\[
dV(b_1,\ldots,b_n)=\sqrt{\det(\langle b_i,b_j\rangle)}.
\]
\end{lemma}

\begin{proof}
\leanok
\uses{lem:gram-schmidt-preserves-orientation}
\sketch Let $e_i$ be the oriented Gram--Schmidt orthonormalization and $A_{ij}=\langle e_i,b_j\rangle$. The Gram matrix is $A^{\mathsf T}A$, so its determinant is $(\det A)^2$. Orientation gives $\det A>0$, while the volume determinant formula gives $dV(b)=\det A$.
\end{proof}

\begin{proposition}[The Value of $dV_g$ on an Oriented Frame]
\label{prop:volume-form-frame-value}
\lean{LeeLib.Ch02.RiemannianMetric.volumeForm_apply_eq_sqrt_det_frameMatrix}
\uses{def:volume-form-pointwise, def:frame-matrix, lem:volume-form-gram-determinant}
\leanok



If $(Y_1,\ldots,Y_n)$ is positively oriented at $x$, then
\[
dV_g|_x(Y_1|_x,\ldots,Y_n|_x)=\sqrt{\det(g_{ij}(x))}.
\]
\end{proposition}

\begin{proof}
\leanok
\uses{lem:volume-form-gram-determinant,def:frame-matrix}
\sketch Apply Lemma \ref{lem:volume-form-gram-determinant} to the basis $(Y_i|_x)$; its Gram matrix is $(g_{ij}(x))$.
\end{proof}

\begin{proposition}[Lee 2.41(c): the Coordinate Formula for $dV_g$]
\label{prop:volume-form-coordinate-formula}
\lean{LeeLib.Ch02.RiemannianMetric.volumeForm_eq_sqrt_det_smul_wedgeCovectors}
\uses{def:volume-form-pointwise, def:frame-matrix, def:wedge-covectors, prop:volume-form-frame-value, lem:top-forms-proportional, lem:dual-coframe-duality}
\leanok



If $(Y_i)$ is positively oriented at $x$ and $(\varepsilon^i)$ is a dual coframe there, then
\[
dV_g|_x=\sqrt{\det(g_{ij}(x))}\,\varepsilon^1\wedge\cdots\wedge\varepsilon^n.
\]
Thus in oriented coordinates, $dV_g=\sqrt{\det(g_{ij})}\,dx^1\wedge\cdots\wedge dx^n$.
\end{proposition}

\begin{proof}
\leanok
\uses{prop:volume-form-frame-value,lem:top-forms-proportional,def:wedge-covectors,lem:dual-coframe-duality}
\sketch Both sides are top forms, so compare them on $(Y_i|_x)$. The dual wedge has value $1$ there by Lemma \ref{lem:dual-coframe-duality}; Proposition \ref{prop:volume-form-frame-value} gives the same scalar for $dV_g$.
\end{proof}

\begin{proposition}[The Riemannian Volume Form]
\label{prop:riemannian-volume-form}
\lean{LeeLib.Ch02.RiemannianMetric.riemannian_volumeForm}
\uses{def:volume-form-pointwise, prop:volume-form-smooth, prop:volume-form-value-one, prop:volume-form-uniqueness, prop:volume-form-determinant, prop:volume-form-coframe, def:wedge-covectors, prop:volume-form-coordinate-formula}
\leanok
\dcref{ch2:2.41}



On an oriented Riemannian $n$-manifold, with or without boundary, there is a unique smooth $n$-form $dV_g$ characterized equivalently by:
\begin{enumerate}[label=(\alph*)]
\item $dV_g=e^1\wedge\cdots\wedge e^n$ for every local oriented orthonormal coframe;
\item $dV_g(E_1,\ldots,E_n)=1$ for every local oriented orthonormal frame;
\item $dV_g=\sqrt{\det(g_{ij})}\,dx^1\wedge\cdots\wedge dx^n$ in every oriented coordinate chart.
\end{enumerate}
\end{proposition}

\begin{proof}
\leanok
\uses{def:volume-form-pointwise,prop:volume-form-smooth,prop:volume-form-uniqueness,prop:volume-form-coframe,prop:volume-form-value-one,prop:volume-form-coordinate-formula}
\sketch Definition \ref{def:volume-form-pointwise} and Proposition \ref{prop:volume-form-smooth} give a smooth form; Proposition \ref{prop:volume-form-uniqueness} gives uniqueness from (b). Propositions \ref{prop:volume-form-coframe}, \ref{prop:volume-form-value-one}, and \ref{prop:volume-form-coordinate-formula} establish (a)--(c), and each condition determines the same form by evaluation on an oriented orthonormal frame.
\end{proof}

For a compactly supported continuous real-valued function $f$ on an oriented Riemannian manifold, define
\[
\int_M f\,dV_g:=\int_M f dV_g,
\qquad \operatorname{Vol}(M)=\int_MdV_g
\]
when $M$ is compact. The same definitions apply to a regular domain with its induced metric. If $M$ is compact, oriented, and without boundary, $dV_g$ is not exact because its integral is positive while exact forms integrate to zero. Properties (a) and (c) of Proposition \ref{prop:riemannian-volume-form} are independent of the wedge convention; the determinant convention makes (b) equal to $1$, while the Alt convention gives $1/n!$.

\begin{definition}[The Pullback of a Differential Form]
\label{def:pullback-alternating}
\lean{LeeLib.Ch02.pullbackAlternating}
\leanok


For a smooth $\varphi:M\to\widetilde M$ and a $k$-form $w$ on $\widetilde M$, define
\[
((\varphi^*w)|_x)(v_1,\ldots,v_k)=w|_{\varphi(x)}(d\varphi_xv_1,\ldots,d\varphi_xv_k).
\]
\end{definition}

\begin{lemma}[A Metric-Preserving Map Pushes Orthonormal Frames Forward]
\label{lem:orthonormal-pushforward}
\lean{LeeLib.Ch02.IsMetricPreserving.orthonormal_pushforward}
\uses{def:metric-preserving}
\leanok



If $\varphi:(M,g)\to(\widetilde M,\widetilde g)$ preserves the metric and $(E_i)$ is $g$-orthonormal, then $(d\varphi_x(E_i|_x))$ is $\widetilde g$-orthonormal.
\end{lemma}

\begin{proof}
\leanok
\uses{def:metric-preserving}
\sketch Metric preservation gives $\langle d\varphi_xE_i,d\varphi_xE_j\rangle_{\widetilde g}=\langle E_i,E_j\rangle_g=\delta_{ij}$.
\end{proof}

\begin{lemma}[The Pushforward Frame is a Basis]
\label{lem:push-basis}
\lean{LeeLib.Ch02.IsMetricPreserving.pushBasis}
\uses{lem:orthonormal-pushforward}
\leanok



In the preceding situation, $(d\varphi_x(E_i|_x))$ is a basis of $T_{\varphi(x)}\widetilde M$.
\end{lemma}

\begin{proof}
\leanok
\uses{lem:orthonormal-pushforward}
\sketch It is an orthonormal, hence linearly independent, family of $n=\dim T_{\varphi(x)}\widetilde M$ vectors.
\end{proof}

\begin{proposition}\label{prop:isometry-preserves-volume-form}\dcref{ch2:2.42}


\uses{def:pullback-alternating,def:volume-form-pointwise,prop:volume-form-uniqueness,prop:volume-form-value-one,lem:orthonormal-pushforward,lem:push-basis}
If $\varphi:(M,g)\to(\widetilde M,\widetilde g)$ is an orientation-preserving
isometry of oriented Riemannian manifolds, then
$\varphi^*dV_{\widetilde g}=dV_g$.
\end{proposition}

\begin{proof}
\leanok
\uses{def:pullback-alternating,prop:volume-form-uniqueness,prop:volume-form-value-one,lem:orthonormal-pushforward,lem:push-basis}
\sketch Evaluate the pullback on a local oriented orthonormal frame $(E_i)$ at $x$. Its value is $dV_{\widetilde g}(d\varphi_xE_i)$, which is $1$ because the pushforward is a positively oriented orthonormal basis. Proposition \ref{prop:volume-form-uniqueness} then gives equality pointwise.
\end{proof}

\begin{definition}[Interior Multiplication]
\label{def:interior-multiplication}
\lean{LeeLib.Ch02.camContract}
\leanok


For a continuous alternating $(k+1)$-form $\mu$ on a real topological vector space and $x$ in the space, define the interior product $x\lrcorner\mu$ by
\[
(x\lrcorner\mu)(v_1,\ldots,v_k)=\mu(x,v_1,\ldots,v_k).
\]
\end{definition}

\begin{definition}[Unit Normal]
\label{def:unit-normal}
\lean{LeeLib.Ch02.IsUnitNormalAt}
\uses{def:normal-space}
\leanok



If $F:M\to\widetilde M$ immerses $M$ in $(\widetilde M,\widetilde g)$, a family $N_x\in T_{F(x)}\widetilde M$ is a \emph{unit normal} when
\[
\widetilde g_{F(x)}(N_x,dF_xv)=0\quad(v\in T_xM),\qquad \widetilde g_{F(x)}(N_x,N_x)=1.
\]
\end{definition}

\begin{lemma}[The Adapted Frame Along a Hypersurface]
\label{lem:adapted-frame-cons-normal}
\lean{LeeLib.Ch02.inner_consNormal}
\uses{def:unit-normal, def:metric-preserving, def:normal-space}
\leanok



If $F:M\to\widetilde M$ is an isometric immersion of an $n$-manifold into an $(n+1)$-manifold, $N_x$ is unit normal, and $(E_i)$ is a $g$-orthonormal basis of $T_xM$, then $(N_x,dF_xE_1,\ldots,dF_xE_n)$ is an orthonormal family in $T_{F(x)}\widetilde M$.
\end{lemma}

\begin{proof}
\leanok
\uses{def:unit-normal,def:metric-preserving}
\sketch Normality makes the cross terms zero, the unit-normal condition gives the first norm, and metric preservation gives $\widetilde g(dF_xE_i,dF_xE_j)=g(E_i,E_j)=\delta_{ij}$.
\end{proof}

\begin{lemma}[The Adapted Frame is an Ambient Basis]
\label{lem:adapted-frame-ambient-basis}
\lean{LeeLib.Ch02.consNormalBasis}
\uses{lem:adapted-frame-cons-normal}
\leanok



The adapted family of Lemma \ref{lem:adapted-frame-cons-normal} is a basis of $T_{F(x)}\widetilde M$.
\end{lemma}

\begin{proof}
\leanok
\uses{lem:adapted-frame-cons-normal}
\sketch It is an orthonormal independent family of $n+1$ vectors in an $(n+1)$-dimensional space.
\end{proof}

\begin{proposition}[Equation (2.17)]
\label{prop:hypersurface-volume-form-formula}
\lean{LeeLib.Ch02.volumeForm_eq_normalRestrict}
\uses{def:interior-multiplication, def:unit-normal, lem:adapted-frame-cons-normal, lem:adapted-frame-ambient-basis, def:volume-form-pointwise, prop:volume-form-uniqueness, prop:volume-form-value-one, def:metric-preserving}
\leanok



Let $F:M\to\widetilde M$ be an isometric immersion of an $n$-manifold into an oriented Riemannian $(n+1)$-manifold. If $N_x$ is unit normal and a local oriented orthonormal frame $(E_i)$ has adapted basis $(N_x,dF_xE_i)$ positively oriented, then
\begin{equation}
dV_g|_x=\bigl((N\lrcorner dV_{\widetilde g})|_M\bigr)_x. \tag{2.17}
\end{equation}
Here the right side maps $(v_1,\ldots,v_n)$ to $dV_{\widetilde g}|_{F(x)}(N_x,dF_xv_1,\ldots,dF_xv_n)$.
\end{proposition}

\begin{proof}
\leanok
\uses{def:interior-multiplication,lem:adapted-frame-cons-normal,lem:adapted-frame-ambient-basis,prop:volume-form-uniqueness,prop:volume-form-value-one}
\sketch On $(E_i)$, the contracted ambient form evaluates to $1$ because the adapted family is a positively oriented orthonormal ambient basis. Proposition \ref{prop:volume-form-uniqueness} identifies it with $dV_g|_x$.
\end{proof}

\begin{lemma}[A Nonvanishing Top Form Orients the Space]
\label{lem:nonvanishing-top-form-orientation}
\lean{LeeLib.Ch02.rayOfNeZero_eq_basis_orientation_iff}
\leanok


If $\omega$ is a nonzero alternating $n$-form on an $n$-dimensional real vector space, it determines the orientation whose positive bases $b$ satisfy $\omega(b)>0$.
\end{lemma}

\begin{proof}
\leanok
\sketch Write $\omega=\omega(b)\det_b$ for any basis $b$. Its ray agrees with the orientation of $b$ exactly when $\omega(b)>0$.
\end{proof}

\begin{definition}[The Orientation Induced by a Unit Normal]
\label{def:induced-orientation}
\lean{LeeLib.Ch02.inducedOrientation}
\uses{def:unit-normal, def:interior-multiplication, lem:nonvanishing-top-form-orientation, lem:adapted-frame-ambient-basis, def:volume-form-pointwise}



For a unit normal field $N$ along an isometric hypersurface immersion, the \emph{induced orientation} at $x$ is the orientation determined by the nonzero form $((N\lrcorner dV_{\widetilde g})|_M)_x$.
\end{definition}

\begin{lemma}[The Induced Orientation is Lee's]
\label{lem:induced-orientation-characterization}
\lean{LeeLib.Ch02.consNormalBasis_orientation_eq_iff}
\uses{def:induced-orientation, lem:adapted-frame-ambient-basis, lem:nonvanishing-top-form-orientation, prop:volume-form-value-one}
\leanok



For a $g$-orthonormal basis $(E_i)$ of $T_xM$, it is positively oriented for the induced orientation iff $(N_x,dF_xE_1,\ldots,dF_xE_n)$ is positively oriented in $\widetilde M$.
\end{lemma}

\begin{proof}
\leanok
\uses{def:induced-orientation,lem:nonvanishing-top-form-orientation,prop:volume-form-value-one}
\sketch The contracted form evaluates on $(E_i)$ as the ambient volume form on the adapted basis, which is $1$ or $-1$ according to its orientation. The positivity criterion of Lemma \ref{lem:nonvanishing-top-form-orientation} gives the equivalence.
\end{proof}

\begin{lemma}[Top Forms are Determined by One Basis]
\label{lem:top-form-eq-of-basis}
\lean{LeeLib.Ch02.eq_of_apply_basis, ContinuousAlternatingMap.ext_of_apply_basis_eq}
\leanok


If two alternating $n$-forms on an $n$-dimensional vector space agree on one basis, they are equal.
\end{lemma}

\begin{proof}
\leanok
\sketch Each form is its value on the basis times that basis's determinant form.
\end{proof}

\begin{proposition}[Equation (2.17) for the Induced Orientation]
\label{prop:hypersurface-volume-form-induced}
\lean{LeeLib.Ch02.volumeForm_inducedOrientation_eq_normalRestrict}
\uses{def:induced-orientation, lem:induced-orientation-characterization, lem:top-form-eq-of-basis, prop:volume-form-value-one, def:interior-multiplication, def:unit-normal}
\leanok



For an isometric immersion into an oriented Riemannian $(n+1)$-manifold, the orientation induced by a unit normal $N$ satisfies
\[
dV_g=(N\lrcorner dV_{\widetilde g})|_M,
\]
without any prior compatibility assumption between the orientations.
\end{proposition}

\begin{proof}
\leanok
\uses{def:induced-orientation,lem:induced-orientation-characterization,lem:top-form-eq-of-basis,prop:volume-form-value-one}
\sketch Compare the two forms on any $g$-orthonormal basis. If the basis is positively oriented for the induced orientation, both values are $1$; otherwise both are $-1$, by Lemma \ref{lem:induced-orientation-characterization}. Lemma \ref{lem:top-form-eq-of-basis} finishes.
\end{proof}

\begin{lemma}[Smoothness of a Pulled-Back Section]
\label{lem:pullback-comp-section}
\lean{LeeLib.Ch02.contMDiffAt_comp_section}
\uses{lem:ambient-tangent-bundle-riemannian}
\leanok



If $F:M\to\widetilde M$ is smooth and $Z$ is a smooth ambient tangent section near $F(x_0)$, then $x\mapsto Z_{F(x)}$ is a smooth section of $F^*T\widetilde M$ near $x_0$.
\end{lemma}

\begin{proof}
\leanok
\uses{lem:ambient-tangent-bundle-riemannian}
\sketch Pull a local trivialization of $T\widetilde M$ back along $F$; in these trivializations the section is the composite of the smooth coordinate representation of $Z$ with $F$.
\end{proof}

\begin{definition}[Smoothness of an Ambient Family at a Point]
\label{def:smooth-normal-at}
\lean{LeeLib.Ch02.IsSmoothNormalAt}
\uses{lem:ambient-tangent-bundle-riemannian}
\leanok



For an immersion $F:M\to\widetilde M$ and a family $N_x\in T_{F(x)}\widetilde M$, say $N$ \emph{varies smoothly at} $x_0$ when it is a smooth section of $F^*T\widetilde M$ at $x_0$.
\end{definition}

\begin{definition}[Smooth Unit Normal Field]
\label{def:smooth-unit-normal-field}
\lean{LeeLib.Ch02.IsSmoothNormalField}
\uses{def:unit-normal, def:smooth-normal-at, lem:ambient-tangent-bundle-riemannian, def:normal-space}
\leanok



In the preceding setting, $N$ \emph{varies smoothly} if it varies smoothly at every point. A smooth global unit normal field is a smoothly varying family that is a unit normal everywhere.
\end{definition}

\begin{lemma}[Negating One Member of a Local Frame]
\label{lem:frame-neg-member}
\lean{LeeLib.Ch02.isLocalFrameOn_negMember}
\leanok


If $(Y_i)_{i\in\iota}$ is a smooth local frame and one member is replaced by its negative, the result is again a smooth local frame.
\end{lemma}

\begin{proof}
\leanok
\sketch Negation is smooth and multiplying one basis vector by $-1$ preserves linear independence and spanning.
\end{proof}

\begin{lemma}[The Adapted Frame Varies Smoothly]
\label{lem:adapted-frame-smooth}
\lean{LeeLib.Ch02.contMDiffAt_consNormal}
\uses{def:smooth-normal-at, lem:pushforward-smooth, def:unit-normal, lem:ambient-tangent-bundle-riemannian}
\leanok



If $N$ varies smoothly along an immersion $F:M\to\widetilde M$ and $(E_i)$ is a smooth local frame of $TM$, every member of $(N,dFE_1,\ldots,dFE_n)$ is a smooth section of $F^*T\widetilde M$.
\end{lemma}

\begin{proof}
\leanok
\uses{def:smooth-normal-at,lem:pushforward-smooth}
\sketch $N$ is smooth by hypothesis and each $dF(E_i)$ is smooth by the pushforward-smoothness lemma.
\end{proof}

\begin{lemma}[Smoothness of the Value on a Frame]
\label{lem:normal-restrict-frame-smooth}
\lean{LeeLib.Ch02.contMDiffAt_normalRestrict_frame}
\uses{def:smooth-normal-at, lem:adapted-frame-smooth, prop:oriented-orthonormal-frames, prop:volume-form-determinant, def:interior-multiplication, def:smooth-orientation, lem:ambient-tangent-bundle-riemannian, lem:pullback-comp-section}
\leanok



If $N$ varies smoothly along an immersion into an oriented Riemannian $(n+1)$-manifold and $(E_i)$ is a smooth local frame of $TM$, then
\[
x\mapsto dV_{\widetilde g}|_{F(x)}(N_x,dF_xE_1|_x,\ldots,dF_xE_n|_x)
\]
is smooth.
\end{lemma}

\begin{proof}
\leanok
\uses{lem:adapted-frame-smooth,prop:oriented-orthonormal-frames,prop:volume-form-determinant,lem:pullback-comp-section,lem:ambient-tangent-bundle-riemannian}
\sketch Near $F(x_0)$ choose a smooth oriented orthonormal ambient frame. Proposition \ref{prop:volume-form-determinant} expresses the displayed value as a determinant of pairings between this frame and the adapted frame. All sections and pairings are smooth by Lemmas \ref{lem:pullback-comp-section} and \ref{lem:adapted-frame-smooth}, so the determinant is smooth.
\end{proof}

\begin{lemma}[Negating a Frame Member Negates the Value]
\label{lem:normal-restrict-neg-member}
\lean{LeeLib.Ch02.normalRestrict_negMember}
\uses{def:interior-multiplication, def:unit-normal}
\leanok



For the contracted form in Proposition \ref{prop:hypersurface-volume-form-formula}, replacing any argument $E_{i_0}$ by $-E_{i_0}$ negates its value.
\end{lemma}

\begin{proof}
\leanok
\sketch The contracted form is alternating and therefore linear in each tangent argument.
\end{proof}

\begin{proposition}[The Induced Orientation is Smooth]
\label{prop:hypersurface-orientable-of-normal}
\lean{LeeLib.Ch02.isSmoothOrientation_inducedOrientation}
\uses{def:induced-orientation, def:smooth-unit-normal-field, def:smooth-orientation, lem:normal-restrict-frame-smooth, lem:frame-neg-member, lem:normal-restrict-neg-member, lem:nonvanishing-top-form-orientation, prop:existence-orthonormal-frames, def:interior-multiplication, lem:adapted-frame-ambient-basis, lem:induced-orientation-characterization}
\leanok



If $n\ge1$ and an immersion $F:M^n\to\widetilde M^{n+1}$ into an oriented Riemannian manifold has a smooth global unit normal, then the induced orientation of $M$ is smooth; in particular $M$ is orientable.
\end{proposition}

\begin{proof}
\leanok
\uses{def:induced-orientation,lem:normal-restrict-frame-smooth,lem:frame-neg-member,lem:normal-restrict-neg-member,lem:nonvanishing-top-form-orientation,prop:existence-orthonormal-frames}
\sketch On a neighbourhood choose a smooth orthonormal frame $(E_i)$ and let $s(x)$ be the contracted ambient volume form evaluated on it. Lemma \ref{lem:normal-restrict-frame-smooth} makes $s$ continuous and it never vanishes. On the region where $s>0$ the frame is positively oriented; where $s<0$, negate one member using Lemmas \ref{lem:frame-neg-member} and \ref{lem:normal-restrict-neg-member}. These local frames prove smoothness.
\end{proof}

\begin{remark}
\label{rem:hypersurface-orientation-dimension}
\uses{prop:hypersurface-orientable-of-normal, def:smooth-orientation, def:induced-orientation}

The assumption $n\ge1$ in Proposition \ref{prop:hypersurface-orientable-of-normal} is necessary for this frame argument: a $0$-manifold has only the empty frame, while the induced orientation in a $1$-manifold can be the opposite ray determined by $dV_{\widetilde g}(N_x)=\pm1$.
\end{remark}

\begin{proposition}[Hypersurface Volume Form: the Sufficiency Half]
\label{prop:hypersurface-volume-form-if}
\lean{LeeLib.Ch02.exists_smoothOrientation_volumeForm_eq_normalRestrict}
\uses{prop:hypersurface-orientable-of-normal, prop:hypersurface-volume-form-induced, def:induced-orientation, def:smooth-unit-normal-field, def:smooth-orientation, def:interior-multiplication}
\leanok



If $n\ge1$ and an isometric hypersurface immersion has a smooth global unit normal $N$, then $M$ is orientable and has a smooth orientation satisfying $dV_g=(N\lrcorner dV_{\widetilde g})|_M$.
\end{proposition}

\begin{proof}
\leanok
\uses{prop:hypersurface-orientable-of-normal,prop:hypersurface-volume-form-induced}
\sketch Use the smooth induced orientation from Proposition \ref{prop:hypersurface-orientable-of-normal}; equation (2.17) is Proposition \ref{prop:hypersurface-volume-form-induced}.
\end{proof}

\begin{lemma}[A Hypersurface Has Exactly Two Unit Normals at Each Point]
\label{lem:two-unit-normals}
\lean{LeeLib.Ch02.eq_or_eq_neg_of_isUnitNormalVecAt}
\uses{def:unit-normal, lem:normal-space-orthonormal-frame, def:metric-preserving}
\leanok



If the normal space at $x$ is the line spanned by a unit normal $w_0$, every unit normal at $x$ is $w_0$ or $-w_0$.
\end{lemma}

\begin{proof}
\leanok
\uses{def:unit-normal}
\sketch Write $w=cw_0$. The unit equations give $1=\langle w,w\rangle=c^2\langle w_0,w_0\rangle=c^2$, hence $c=\pm1$.
\end{proof}

\begin{lemma}[The Normal Space of a Hypersurface Is a Line]
\label{lem:pointwise-unit-normal-span}
\lean{LeeLib.Ch02.exists_isUnitNormalVecAt_span}
\uses{def:unit-normal, lem:normal-space-orthonormal-frame}
\leanok



For a hypersurface immersion, each point has a unit normal $w_0$ with normal space $N_xM=\mathbb R\,w_0$.
\end{lemma}

\begin{proof}
\leanok
\uses{lem:normal-space-orthonormal-frame}
\sketch The local orthonormal normal-frame lemma supplies one normal vector because the codimension is $1$; its value at $x$ is a unit spanning vector.
\end{proof}

\begin{lemma}[Negating the Normal Negates the Contracted Form]
\label{lem:normal-restrict-neg-normal}
\lean{LeeLib.Ch02.normalRestrictVec_neg}
\uses{def:interior-multiplication, def:unit-normal, def:volume-form-pointwise}
\leanok



For every $w\in T_{F(x)}\widetilde M$,
\[
(((-w)\lrcorner dV_{\widetilde g})|_M)_x=-((w\lrcorner dV_{\widetilde g})|_M)_x.
\]
\end{lemma}

\begin{proof}
\leanok
\sketch Interior multiplication places $w$ in the first slot of a multilinear form, so negating $w$ negates the contracted form.
\end{proof}

\begin{lemma}[A Basis Realizing a Prescribed Orientation]
\label{lem:basis-realizing-orientation}
\lean{LeeLib.Ch02.exists_basis_orientation_eq}
\uses{lem:nonvanishing-top-form-orientation}
\leanok



If $n\ge1$ and $o$ is a pointwise orientation of $M$, every $T_xM$ has a basis positively oriented for $o(x)$.
\end{lemma}

\begin{proof}
\leanok
\sketch Start with any basis. If its orientation is opposite to $o(x)$, negate one member; $n\ge1$ makes this possible.
\end{proof}

\begin{definition}[A Unit Normal Inducing a Given Orientation]
\label{def:oriented-unit-normal}
\lean{LeeLib.Ch02.IsOrientedUnitNormalVecAt}
\uses{def:unit-normal, def:interior-multiplication, def:induced-orientation, lem:nonvanishing-top-form-orientation}
\leanok



A unit normal $w$ at $x$ \emph{induces} a pointwise orientation $o(x)$ if $((w\lrcorner dV_{\widetilde g})|_M)_x(b)>0$ for every basis $b$ positively oriented for $o(x)$.
\end{definition}

\begin{lemma}[One Basis Suffices]
\label{lem:oriented-unit-normal-iff-pos}
\lean{LeeLib.Ch02.isOrientedUnitNormalVecAt_iff_pos}
\uses{def:oriented-unit-normal, lem:nonvanishing-top-form-orientation, def:unit-normal, def:metric-preserving, lem:adapted-frame-ambient-basis}
\leanok



A unit normal induces $o(x)$ iff its contracted form is positive on one basis positively oriented for $o(x)$.
\end{lemma}

\begin{proof}
\leanok
\uses{lem:nonvanishing-top-form-orientation}
\sketch A nonzero top form determines an orientation by its ray, so positivity on one positive basis is equivalent to positivity on every basis in the same orientation.
\end{proof}

\begin{proposition}[An Orientation Selects a Unit Normal]
\label{prop:existsunique-oriented-unit-normal}
\lean{LeeLib.Ch02.existsUnique_isOrientedUnitNormalVecAt}
\uses{def:oriented-unit-normal, lem:two-unit-normals, lem:pointwise-unit-normal-span, lem:normal-restrict-neg-normal, lem:basis-realizing-orientation, lem:oriented-unit-normal-iff-pos, def:unit-normal, def:metric-preserving}
\leanok



If $n\ge1$ and $o$ is a pointwise orientation of an isometrically immersed hypersurface, exactly one unit normal at each $x$ induces $o(x)$.
\end{proposition}

\begin{proof}
\leanok
\uses{lem:two-unit-normals,lem:pointwise-unit-normal-span,lem:normal-restrict-neg-normal,lem:basis-realizing-orientation,lem:oriented-unit-normal-iff-pos}
\sketch The normal space is spanned by a unit $w_0$ and its only unit normals are $\pm w_0$. Evaluate the contracted form of $w_0$ on an $o(x)$-positive basis; the value is nonzero, and exactly one of it and its negative is positive.
\end{proof}

\begin{definition}[The Unit Normal Determined by an Orientation]
\label{def:unit-normal-of-orientation}
\lean{LeeLib.Ch02.unitNormalOfOrientation}
\uses{prop:existsunique-oriented-unit-normal, def:oriented-unit-normal}
\leanok



The unit normal determined by a pointwise orientation $o$ assigns to each point the unique unit normal inducing $o(x)$.
\end{definition}

\begin{proposition}[The Normal Determined by an Orientation Is Smooth]
\label{prop:unit-normal-of-orientation-smooth}
\lean{LeeLib.Ch02.isSmoothNormalField_unitNormalOfOrientation}
\uses{def:unit-normal-of-orientation, def:smooth-unit-normal-field, def:smooth-orientation, lem:normal-space-orthonormal-frame, lem:normal-restrict-frame-smooth, lem:oriented-unit-normal-iff-pos, lem:normal-restrict-neg-normal, prop:existsunique-oriented-unit-normal, def:interior-multiplication}
\leanok



If the ambient and hypersurface orientations are smooth, the unit normal determined by the hypersurface orientation is a smooth global unit normal field.
\end{proposition}

\begin{proof}
\leanok
\uses{lem:normal-space-orthonormal-frame,lem:normal-restrict-frame-smooth,lem:oriented-unit-normal-iff-pos,lem:normal-restrict-neg-normal,prop:existsunique-oriented-unit-normal}
\sketch Near $p$, choose a smooth positive tangent frame and a smooth local unit normal $W$. The smooth nonzero function $s(x)=((W\lrcorner dV_{\widetilde g})|_M)_x(E_1,\ldots,E_n)$ has constant sign on a smaller neighbourhood. The determined normal equals $W$ where $s>0$ and $-W$ where $s<0$, hence is smooth.
\end{proof}

\begin{proposition}[Hypersurface Volume Form: the Necessity Half]
\label{prop:hypersurface-volume-form-only-if}
\lean{LeeLib.Ch02.exists_isSmoothNormalField_volumeForm_eq_normalRestrict}
\uses{prop:unit-normal-of-orientation-smooth, def:unit-normal-of-orientation, prop:hypersurface-volume-form-induced, def:induced-orientation, lem:oriented-unit-normal-iff-pos, lem:basis-realizing-orientation, def:smooth-unit-normal-field, def:smooth-orientation, def:interior-multiplication}
\leanok



If $n\ge1$ and the immersed hypersurface has a smooth orientation, it has a smooth global unit normal $N$ and $dV_g=(N\lrcorner dV_{\widetilde g})|_M$.
\end{proposition}

\begin{proof}
\leanok
\uses{prop:unit-normal-of-orientation-smooth,prop:hypersurface-volume-form-induced,lem:oriented-unit-normal-iff-pos,lem:basis-realizing-orientation}
\sketch Take the smooth normal determined by the given orientation. Its induced orientation is the original one by the one-basis positivity criterion, so Proposition \ref{prop:hypersurface-volume-form-induced} gives (2.17).
\end{proof}

\begin{proposition}
\label{prop:hypersurface-volume-form}
\lean{LeeLib.Ch02.exists_isSmoothOrientation_iff_exists_isSmoothNormalField}
\uses{prop:hypersurface-volume-form-formula, prop:hypersurface-volume-form-induced, prop:hypersurface-volume-form-if, prop:hypersurface-volume-form-only-if, prop:hypersurface-orientable-of-normal, prop:unit-normal-of-orientation-smooth, def:induced-orientation, def:smooth-unit-normal-field, def:unit-normal, def:interior-multiplication, lem:normal-space-orthonormal-frame, lem:induced-orientation-characterization}
\leanok
\dcref{ch2:2.43}



For $n\ge1$, an isometric hypersurface immersion into an oriented Riemannian $(n+1)$-manifold is orientable iff it admits a smooth global unit normal field.
\end{proposition}

\begin{proof}
\leanok
\uses{prop:hypersurface-volume-form-if,prop:hypersurface-volume-form-only-if,prop:hypersurface-orientable-of-normal,prop:unit-normal-of-orientation-smooth,def:induced-orientation,def:smooth-unit-normal-field}
\sketch A smooth normal gives a smooth induced orientation by Proposition \ref{prop:hypersurface-orientable-of-normal}; a smooth orientation gives a smooth determined normal by Proposition \ref{prop:unit-normal-of-orientation-smooth}. In either direction equation (2.17) holds for the corresponding orientation and normal.
\end{proof}

\subsection{The density of an inner product space}

\begin{definition}[Density of a Vector Space]
\label{def:density-vector-space}
\lean{LeeLib.Ch02.IsDensity}
\leanok


For an $n$-dimensional real vector space, a \emph{density} is a function $\mu:V^n\to\mathbb R$ satisfying
\begin{equation}
\mu(fv_1,\ldots,fv_n)=|\det f|\,\mu(v_1,\ldots,v_n) \tag{B.14}
\end{equation}
for every linear $f:V\to V$. It is positive if it is positive on every basis.
\end{definition}

\begin{definition}[The Density of an Inner Product Space]
\label{def:density-pointwise}
\lean{LeeLib.Ch02.densityL}
\leanok


For an $n$-dimensional real inner product space, define
\[
|v_1\wedge\cdots\wedge v_n|:=\sqrt{\det(\langle v_i,v_j\rangle)_{i,j=1}^n}.
\]
The Gram determinant is nonnegative, so the square root is defined.
\end{definition}

\begin{proposition}[The Determinant Formula for the Density]
\label{prop:density-abs-det}
\lean{LeeLib.Ch02.densityL_eq_abs_det}
\uses{def:density-pointwise}
\leanok



For any orthonormal basis $(e_i)$,
\[
|v_1\wedge\cdots\wedge v_n|=\left|\det(\langle e_i,v_j\rangle)_{i,j=1}^n\right|.
\]
\end{proposition}

\begin{proof}
\leanok
\uses{def:density-pointwise}
\sketch If $M=(\langle e_i,v_j\rangle)$, orthonormal expansion gives the Gram matrix $M^{\mathsf T}M$. Its determinant is $(\det M)^2$, whose square root is $|\det M|$.
\end{proof}

\begin{proposition}[The Density of an Orthonormal Basis]
\label{prop:density-value-one}
\lean{LeeLib.Ch02.densityL_apply_orthonormalBasis}
\uses{prop:density-abs-det}
\leanok



An orthonormal basis $(e_1,\ldots,e_n)$ has density $|e_1\wedge\cdots\wedge e_n|=1$.
\end{proposition}

\begin{proof}
\leanok
\uses{prop:density-abs-det}
\sketch The determinant matrix in Proposition \ref{prop:density-abs-det} is the identity.
\end{proof}

\begin{proposition}[The Transformation Law]
\label{prop:density-transformation-law}
\lean{LeeLib.Ch02.densityL_comp}
\uses{prop:density-abs-det, def:density-vector-space}
\leanok



For linear $f:V\to V$,
\[
|fv_1\wedge\cdots\wedge fv_n|=|\det f|\,|v_1\wedge\cdots\wedge v_n|.
\]
Thus the inner-product density satisfies Definition \ref{def:density-vector-space}.
\end{proposition}

\begin{proof}
\leanok
\uses{prop:density-abs-det}
\sketch In an orthonormal basis, the coefficient matrix of $(fv_j)$ is $AM$, where $A$ represents $f$ and $M$ represents $(v_j)$. Proposition \ref{prop:density-abs-det} gives $|\det(AM)|=|\det f|\,|\det M|$.
\end{proof}

\begin{proposition}[Positivity of the Density]
\label{prop:density-positive}
\lean{LeeLib.Ch02.densityL_pos}
\uses{def:density-pointwise}
\leanok



If $(v_1,\ldots,v_n)$ is linearly independent, then $|v_1\wedge\cdots\wedge v_n|>0$.
\end{proposition}

\begin{proof}
\leanok
\uses{def:density-pointwise}
\sketch The Gram matrix is positive definite on a linearly independent family, hence has positive determinant; its defining square root is therefore positive.
\end{proof}

\begin{proposition}[Fibrewise Uniqueness of the Density]
\label{prop:density-uniqueness-pointwise}
\lean{LeeLib.Ch02.densityL_unique}
\uses{def:density-vector-space, prop:density-transformation-law, prop:density-value-one}
\leanok



If $\mu$ is a density on an inner product space and $\mu(e_1,\ldots,e_n)=1$ on an orthonormal basis, then $\mu(v_1,\ldots,v_n)=|v_1\wedge\cdots\wedge v_n|$ for all $v_i$.
\end{proposition}

\begin{proof}
\leanok
\uses{prop:density-transformation-law,prop:density-value-one}
\sketch Let $f$ be the linear map with $fe_i=v_i$. Both the density transformation law for $\mu$ and Proposition \ref{prop:density-transformation-law}, together with the value $1$ on $(e_i)$, give $|\det f|$.
\end{proof}

\begin{proposition}[The Density Is the Absolute Value of the Volume Form]
\label{prop:density-eq-abs-volume-form}
\lean{LeeLib.Ch02.densityL_eq_abs_volumeForm}
\uses{def:density-pointwise, prop:volume-form-determinant}
\leanok



If $n>0$ and $\omega$ is the volume form of an oriented inner product space, then
\[
|v_1\wedge\cdots\wedge v_n|=|\omega(v_1,\ldots,v_n)|.
\]
\end{proposition}

\begin{proof}
\leanok
\uses{prop:density-abs-det}
\sketch Choose an oriented orthonormal basis. The volume determinant formula and Proposition \ref{prop:density-abs-det} express the two sides as the absolute value of the same determinant.
\end{proof}

\subsection{Smooth densities and the Riemannian density}

\begin{definition}[Smooth Density]
\label{def:smooth-density}
\lean{LeeLib.Ch02.SmoothDensity}
\uses{def:density-vector-space}
\leanok



A \emph{smooth density} on a smooth $n$-manifold assigns a density $\mu_x$ to each $T_xM$ such that $x\mapsto\mu_x(Y_1|_x,\ldots,Y_n|_x)$ is smooth for every smooth local frame $(Y_i)$.
\end{definition}

\begin{definition}[The Riemannian Density, Pointwise]
\label{def:riemannian-density-pointwise}
\lean{LeeLib.Ch02.RiemannianMetric.densityAt}
\uses{def:density-pointwise, def:riemannian-metric}
\leanok



For a Riemannian manifold $(M,g)$, define the pointwise density by
\[
\mu_x(v_1,\ldots,v_n)=\sqrt{\det(\langle v_i,v_j\rangle_g)_{i,j=1}^n}.
\]
This uses no orientation.
\end{definition}

\begin{proposition}[Smoothness of the Riemannian Density]
\label{prop:density-smooth}
\lean{LeeLib.Ch02.RiemannianMetric.contMDiffOn_densityAt}
\uses{def:riemannian-density-pointwise, prop:density-positive, lem:inner-product-vector-fields-smooth}
\leanok



For a smooth local frame $(Y_i)$, $x\mapsto\mu_x(Y_1|_x,\ldots,Y_n|_x)$ is smooth; hence the pointwise Riemannian density is smooth.
\end{proposition}

\begin{proof}
\leanok
\uses{def:riemannian-density-pointwise,prop:density-positive,lem:inner-product-vector-fields-smooth}
\sketch The Gram entries are smooth pairings, so the determinant is smooth. It is strictly positive on frame values by Proposition \ref{prop:density-positive}; composing with the smooth square root on $(0,\infty)$ gives smoothness.
\end{proof}

\begin{proposition}[The Riemannian Density]\label{prop:riemannian-density}\dcref{ch2:2.44}
\lean{LeeLib.Ch02.riemannian_density_existsUnique}
\uses{def:smooth-density, def:riemannian-density-pointwise, prop:density-smooth, prop:density-transformation-law, prop:density-value-one, prop:density-uniqueness-pointwise, prop:density-positive, prop:existence-orthonormal-frames}
\leanok



Every Riemannian manifold has a unique smooth positive density $\mu$ satisfying
\begin{equation}
\mu(E_1,\ldots,E_n)=1 \tag{2.18}
\end{equation}
for every local orthonormal frame $(E_i)$.
\end{proposition}

\begin{proof}
\leanok
\uses{def:riemannian-density-pointwise,prop:density-smooth,prop:density-transformation-law,prop:density-value-one,prop:density-uniqueness-pointwise,prop:density-positive,prop:existence-orthonormal-frames}
\sketch Take the family of Definition \ref{def:riemannian-density-pointwise}; Proposition \ref{prop:density-transformation-law} makes each fibre a density and Proposition \ref{prop:density-smooth} makes the family smooth. Proposition \ref{prop:density-value-one} gives (2.18) and positivity. For uniqueness, evaluate any competing density on a local orthonormal frame and apply Proposition \ref{prop:density-uniqueness-pointwise} fibrewise.
\end{proof}

\begin{proposition}\label{prop:riemannian-density-formula}\dcref{ch2:2.45}


\uses{prop:riemannian-density,def:density-vector-space}
For every local orthonormal coframe,
$\mu=|\varepsilon^1\wedge\cdots\wedge\varepsilon^n|$. Pointwise, both sides
equal $|\det[\langle E_i,v_j\rangle]|$ on each $T_xM$.
\end{proposition}

On any Riemannian manifold, the density is defined without orientation; when an orientation exists it is $|dV_g|$. Thus compactly supported smooth functions can be integrated against $dV_g$ interpreted as a density on nonorientable manifolds or as the corresponding volume form on orientable ones.

\section{The Divergence and the Laplacian}

For an oriented Riemannian $n$-manifold, define the divergence of a smooth vector field $X$ by
\begin{equation}
d(X\lrcorner dV_g)=(\operatorname{div}X)dV_g. \tag{2.19}
\end{equation}
The same local equation defines divergence on nonorientable manifolds: reversing a local orientation changes both sides' signs. Define the Laplace--Beltrami operator by
\begin{equation}
\Delta u=\operatorname{div}(\operatorname{grad}u). \tag{2.20}
\end{equation}

\begin{proposition}
\label{prop:coordinate-representations-divergence-laplacian}
\dcref{ch2:2.46}
Let $(x^i)$ be smooth coordinates on $U\subseteq M$, write $g_{ij}$ for the metric matrix and $g^{ij}$ for its inverse, and let $\det g=\det(g_{kl})$. Then
\[
\operatorname{div}\left(X^i\frac{\partial}{\partial x^i}\right)=\frac{1}{\sqrt{\det g}}\frac{\partial}{\partial x^i}\left(\sqrt{\det g}\,X^i\right),
\]
\[
\Delta u=\frac{1}{\sqrt{\det g}}\frac{\partial}{\partial x^i}\left(\sqrt{\det g}\,g^{ij}\frac{\partial u}{\partial x^j}\right).
\]
For Euclidean $\mathbb R^n$ with standard coordinates these reduce to
\[
\operatorname{div}\left(X^i\frac{\partial}{\partial x^i}\right)=\sum_{i=1}^n\frac{\partial X^i}{\partial x^i},
\qquad
\Delta u=\sum_{i=1}^n\frac{\partial^2u}{(\partial x^i)^2}.
\]
\end{proposition}

\begin{proof}
\sketch In oriented coordinates, insert
$dV_g=\sqrt{\det g}\,dx^1\wedge\cdots\wedge dx^n$ into the defining equation
for divergence and differentiate. Substituting
$\operatorname{grad}u=g^{ij}(\partial_j u)\partial_i$ gives the Laplacian
formula. The Euclidean identities follow from $g_{ij}=\delta_{ij}$.
\end{proof}
\section{Lengths and Distances}

Throughout, $(M,g)$ is Riemannian with or without boundary. A curve is a continuous map $\gamma:I\to M$ from an interval; a curve segment has compact domain. Smoothness is understood in the category of manifolds with boundary, and $\gamma$ is regular when $\gamma'(t)\ne0$ everywhere.

An admissible curve $\gamma:[a,b]\to M$ is continuous and admits a partition
\[
a=a_0<a_1<\cdots<a_k=b
\]
such that every restriction to $[a_{i-1},a_i]$ is smooth and regular. A one-point map is also admissible. An admissible partition need only make the restrictions smooth; at an interior partition point the nonzero one-sided velocities $\gamma'(a_i^-)$ and $\gamma'(a_i^+)$ may differ.

For smooth $\gamma:I\to M$, a reparametrization is $\gamma\circ\varphi:I'\to M$ with $\varphi$ a diffeomorphism of intervals; it is forward or backward according as $\varphi$ is increasing or decreasing. For admissible curves, $\varphi:[c,d]\to[a,b]$ may instead be a homeomorphism that is a diffeomorphism on every subinterval of some partition. The classical length is
\[
L_g(\gamma)=\int_a^b|\gamma'(t)|_g\,dt.
\]

\begin{definition}[The Length of a Curve, extended form]
\label{def:length-of-curve}
\lean{Manifold.pathELength}
\uses{def:riemannian-metric}

\mathlibok

For a curve $\gamma$ and $a\le b$, define
\[
L_g(\gamma|_{[a,b]})=\int_{[a,b]}|\gamma'(t)|_g\,dt\in[0,\infty]
\]
using the Lebesgue integral. For an admissible curve on a compact interval this is finite and equals the classical length.
\end{definition}

\begin{proposition}[Additivity of Length]
\label{prop:length-additive}
\lean{Manifold.pathELength_add}
\uses{def:length-of-curve}

\mathlibok

\dcref{ch2:2.47}
For an admissible curve and $a\le c\le b$,
\[
L_g(\gamma|_{[a,c]})+L_g(\gamma|_{[c,b]})=L_g(\gamma|_{[a,b]}).
\]
\end{proposition}

\begin{proposition}[Parameter Independence of Length]
\label{prop:length-reparam-invariant}
\lean{Manifold.pathELength_comp_of_monotoneOn}
\uses{def:length-of-curve}

\mathlibok

\dcref{ch2:2.47}
If $\gamma:[a,b]\to M$ is admissible and $\widetilde\gamma=\gamma\circ\varphi$ is a forward reparametrization, then $L_g(\widetilde\gamma)=L_g(\gamma)$.
\end{proposition}

\begin{proposition}[Isometry Invariance of Length]
\label{prop:length-isometry-invariant}
\lean{LeeLib.Ch02.IsMetricPreserving.pathELength_comp}
\uses{def:length-of-curve, def:metric-preserving}
\leanok



\dcref{ch2:2.47}
Let $\varphi:(M,g)\to(\widetilde M,\widetilde g)$ be a smooth metric-preserving map between Riemannian manifolds with or without boundary. For every admissible $\gamma:[a,b]\to M$,
\[
L_{\widetilde g}(\varphi\circ\gamma)=L_g(\gamma).
\]
In particular, length is invariant under local isometries.
\end{proposition}
\begin{proof}
\leanok
\uses{def:metric-preserving,def:length-of-curve}
\sketch Metric preservation gives $|d\varphi_p(v)|_{\widetilde g}=|v|_g$. The chain rule therefore yields $|(\varphi\circ\gamma)'(t)|_{\widetilde g}=|\gamma'(t)|_g$ wherever $\gamma$ is differentiable. The integrands agree outside a finite, hence null, set.
\end{proof}

Propositions \ref{prop:length-additive}, \ref{prop:length-reparam-invariant}, and \ref{prop:length-isometry-invariant} are the three parts of Proposition 2.47.

For admissible $\gamma:[a,b]\to M$, its arc-length function is
\[
s(t)=L_g(\gamma|_{[a,t]})=\int_a^t|\gamma'(u)|_g\,du.
\]
It is continuous, is smooth where $\gamma$ is smooth, and satisfies $s'(t)=|\gamma'(t)|_g$. This norm is the speed. A curve has unit speed if the speed is $1$ wherever defined, and constant speed if it is constant. For a unit-speed admissible curve, $s(t)=t-a$; if its domain is $[0,b]$, it is parametrized by arc length.

\begin{proposition}[Unit-Speed Reparametrization of a Regular Curve]
\label{prop:regular-curve-unit-speed-reparam}
\lean{LeeLib.Ch02.regularCurve_arclengthReparametrization}
\uses{def:length-of-curve}
\leanok


\dcref{ch2:2.49}

Let $J\subseteq\mathbb R$ be open and let $\gamma:J\to M$ be smooth and regular on $[a,b]\subseteq J$. Put
\[
s(t)=\int_a^t|\gamma'(u)|_g\,du,\qquad L=s(b).
\]
There is a smooth increasing inverse $\psi:[0,L]\to[a,b]$ of $s$ such that $\widetilde\gamma=\gamma\circ\psi$ joins $\gamma(a)$ to $\gamma(b)$, has unit speed, and satisfies
\[
\int_0^\sigma|\widetilde\gamma'(u)|_g\,du=\sigma,\qquad 0\le\sigma\le L.
\]
\end{proposition}
\begin{proof}
\leanok
\uses{def:length-of-curve}
\sketch Regularity on the compact interval persists on an open neighborhood, where $s'(t)=|\gamma'(t)|_g>0$. Thus $s$ has a smooth increasing inverse with
\[
\psi'(\sigma)=\frac1{|\gamma'|_g\circ\psi(\sigma)}.
\]
The chain rule gives $|\widetilde\gamma'|_g=1$, hence the integral identity; $\psi(0)=a$ and $\psi(L)=b$ give the endpoints.
\end{proof}

\begin{proposition}
\label{prop:reparametrization-by-arc-length}
\uses{prop:regular-curve-unit-speed-reparam}
\dcref{ch2:2.49}

Let $(M,g)$ be Riemannian with or without boundary.
\begin{enumerate}
\item[(a)] Every regular curve has a unit-speed forward reparametrization.
\item[(b)] Every admissible curve has a unique forward reparametrization by arc length.
\end{enumerate}
\end{proposition}
\begin{proof}
\sketch For a regular $\gamma:I\to M$, fix $t_0\in I$ and set $s(t)=\int_{t_0}^t|\gamma'(u)|_g\,du$. Since $s'>0$, it is an increasing diffeomorphism onto an interval; composing with $s^{-1}$ gives unit speed.

For an admissible curve, apply this construction to each regular segment and concatenate the inverse parameter maps. Induction on the number of segments gives existence. If $\gamma\circ\varphi$ and $\gamma\circ\psi$ are two forward arc-length parametrizations on $[0,c]$, then $\eta=\varphi^{-1}\circ\psi$ is increasing, fixes $0$, and satisfies $1=\eta'$ away from finitely many break points. Continuity gives $\eta(s)=s$, proving uniqueness.
\end{proof}

\section{The Riemannian Distance Function}

\begin{proposition}
\label{prop:connected-manifold-admissible-curve}
\lean{LeeLib.Ch02.exists_contMDiffOn_path_of_preconnected}
\uses{prop:existence-riemannian-metric, prop:riemannian-distance-finite, def:riemannian-distance}
\leanok



\dcref{ch2:2.50}
Any two points of a connected smooth manifold, with or without boundary, can be joined by an admissible curve.
\end{proposition}
\begin{proof}
\leanok
\uses{prop:existence-riemannian-metric,prop:riemannian-distance-finite,def:riemannian-distance}
\sketch A continuous path exists by connectedness and path connectedness. Compactness permits a partition whose image segments lie in coordinate balls or boundary half-balls; replace each by a coordinate line segment. Alternatively, choose a metric by Proposition \ref{prop:existence-riemannian-metric}; Proposition \ref{prop:riemannian-distance-finite} and the distance definition provide a finite-length competitor.
\end{proof}

On a disconnected Riemannian manifold, $d_g(p,q)$ is also used when $p$ and $q$ belong to the same connected component.

\begin{definition}[The Riemannian Distance, extended form]
\label{def:riemannian-distance}
\lean{Manifold.riemannianEDist}
\uses{def:riemannian-metric, def:length-of-curve}

\mathlibok

For a Riemannian manifold $(M,g)$ and $p,q\in M$, set
\[
d_g(p,q)=\inf\{L_g(\gamma):\gamma:[0,1]\to M\text{ is }C^1,\ \gamma(0)=p,\ \gamma(1)=q\}\in[0,\infty],
\]
where the infimum of the empty set is $\infty$.
\end{definition}

\begin{remark}
\label{rem:riemannian-distance-c1-vs-admissible}
\uses{def:riemannian-distance}

Definition \ref{def:riemannian-distance} uses $C^1$ curves, while Lee uses admissible curves. Neither class contains the other. The infima agree: admissible curves can be reparametrized with zero velocity at corners without changing length, and $C^1$ curves can be approximated in length by regular curves. This equivalence is not used below, where distance always means Definition \ref{def:riemannian-distance}.
\end{remark}

\begin{proposition}[Connected Riemannian Manifolds Are $C^1$-Path Connected]
\label{prop:connected-manifold-c1-curve}
\lean{LeeLib.Ch02.exists_contMDiffOn_path}
\uses{def:riemannian-distance, prop:riemannian-distance-finite}
\leanok



Any two points of a connected Riemannian manifold $(M,g)$ can be joined by a $C^1$ curve $[0,1]\to M$. This is the $C^1$ form of Proposition \ref{prop:connected-manifold-admissible-curve} relative to Definition \ref{def:riemannian-distance}; existence of a metric on a smooth manifold follows from Proposition \ref{prop:existence-riemannian-metric}.
\end{proposition}
\begin{proof}
\leanok
\uses{def:riemannian-distance,prop:riemannian-distance-finite}
\sketch By Proposition \ref{prop:riemannian-distance-finite}, $d_g(p,q)<d_g(p,q)+1<\infty$. Definition \ref{def:riemannian-distance} then supplies a $C^1$ competitor from $p$ to $q$ of length below $d_g(p,q)+1$.
\end{proof}

\begin{lemma}[Symmetry of the Riemannian Distance]
\label{lem:riemannian-distance-symm}
\lean{Manifold.riemannianEDist_comm}
\uses{def:riemannian-distance}

\mathlibok

For all $p,q\in M$, $d_g(p,q)=d_g(q,p)$.
\end{lemma}

\begin{lemma}[Triangle Inequality for the Riemannian Distance]
\label{lem:riemannian-distance-triangle}
\lean{Manifold.riemannianEDist_triangle}
\uses{def:riemannian-distance}

\mathlibok

For all $p,q,r\in M$, $d_g(p,r)\le d_g(p,q)+d_g(q,r)$.
\end{lemma}

\begin{lemma}[Nearby Points Are at Small Distance]
\label{lem:riemannian-distance-locally-small}
\lean{eventually_riemannianEDist_lt}
\uses{def:riemannian-distance}

\mathlibok

For $p\in M$ and $c>0$, some neighborhood $U$ of $p$ satisfies $d_g(p,q)<c$ for all $q\in U$. This is Lemma \ref{lem:distance-comparison-neighborhood}(a).
\end{lemma}

\begin{lemma}[Points at Small Distance Are Nearby]
\label{lem:riemannian-distance-small-implies-nearby}
\lean{setOf_riemannianEDist_lt_subset_nhds'}
\uses{def:riemannian-distance}

\mathlibok

For every neighborhood $U$ of $p$, some $c>0$ satisfies
\[
\{q\in M:d_g(p,q)<c\}\subseteq U.
\]
This is Lemma \ref{lem:distance-comparison-neighborhood}(b).
\end{lemma}

\begin{proposition}[Finiteness of the Riemannian Distance]
\label{prop:riemannian-distance-finite}
\lean{LeeLib.Ch02.riemannianEDist_ne_top}
\uses{def:riemannian-distance, lem:riemannian-distance-locally-small, lem:riemannian-distance-symm, lem:riemannian-distance-triangle}
\leanok



If $(M,g)$ is connected, then $d_g(p,q)<\infty$ for all $p,q\in M$. This is the finite-length conclusion associated with Proposition \ref{prop:connected-manifold-admissible-curve}.
\end{proposition}
\begin{proof}
\leanok
\uses{lem:riemannian-distance-locally-small,lem:riemannian-distance-symm,lem:riemannian-distance-triangle}
\sketch Fix $p$ and let $S=\{q:d_g(p,q)<\infty\}$. It contains $p$, using the constant curve. If $q_0\in S$, Lemmas \ref{lem:riemannian-distance-locally-small} and \ref{lem:riemannian-distance-triangle} give a neighborhood on which $d_g(p,\cdot)$ remains finite, so $S$ is open. If $q_0\notin S$, the same local bound, Lemmas \ref{lem:riemannian-distance-triangle} and \ref{lem:riemannian-distance-symm}, show that a nearby $q$ cannot lie in $S$, for otherwise
\[
d_g(p,q_0)\le d_g(p,q)+d_g(q,q_0)<\infty.
\]
Thus $S$ is closed. Connectedness gives $S=M$.
\end{proof}

\begin{lemma}[Metric-Preserving Maps Do Not Increase Distance]
\label{lem:metric-preserving-distance-le}
\lean{LeeLib.Ch02.IsMetricPreserving.riemannianEDist_le}
\uses{def:metric-preserving, def:riemannian-distance, prop:length-isometry-invariant}
\leanok



If $\varphi:(M,g)\to(\widetilde M,\widetilde g)$ is smooth and metric-preserving, then
\[
d_{\widetilde g}(\varphi(x),\varphi(y))\le d_g(x,y),
\qquad x,y\in M.
\]
\end{lemma}
\begin{proof}
\leanok
\uses{prop:length-isometry-invariant,def:riemannian-distance}
\sketch For $r>d_g(x,y)$ choose a competitor $\gamma$ of length below $r$. Proposition \ref{prop:length-isometry-invariant} gives the same length for $\varphi\circ\gamma$, so $d_{\widetilde g}(\varphi(x),\varphi(y))<r$. Let $r$ decrease to $d_g(x,y)$.
\end{proof}

The inequality can be strict because the target may contain shorter curves outside the image.

\begin{proposition}[Isometry Invariance of the Riemannian Distance Function]
\label{prop:isometry-invariance-riemannian-distance}
\lean{LeeLib.Ch02.IsIsometry.riemannianEDist_comp}
\uses{def:isometry, def:riemannian-distance, lem:metric-preserving-distance-le, lem:isometry-symm}
\leanok



\dcref{ch2:2.51}
If $\varphi:(M,g)\to(\widetilde M,\widetilde g)$ is an isometry between connected Riemannian manifolds with or without boundary, then
\[
d_{\widetilde g}(\varphi(x),\varphi(y))=d_g(x,y),
\qquad x,y\in M.
\]
\end{proposition}
\begin{proof}
\leanok
\uses{lem:metric-preserving-distance-le,lem:isometry-symm}
\sketch Lemma \ref{lem:metric-preserving-distance-le} gives one inequality for $\varphi$. Lemma \ref{lem:isometry-symm} makes $\varphi^{-1}$ an isometry, and the same inequality for $\varphi^{-1}$ gives the reverse bound.
\end{proof}

\begin{proposition}[Linear Isometries of Euclidean Space]
\label{prop:euclidean-linear-isometry-distance}
\lean{LeeLib.Ch02.isIsometry_euclideanMetric_riemannianEDist}
\uses{ex:euclidean-metric, def:isometry, prop:isometry-invariance-riemannian-distance, prop:euclidean-riemannian-distance}
\leanok



Let $F,\widetilde F$ be inner product spaces with Euclidean metrics and let $e:F\to\widetilde F$ be a surjective linear isometry. Then $e$ is a Riemannian isometry and
\[
d_{\bar g}(e(x),e(y))=|x-y|,
\qquad x,y\in F.
\]
\end{proposition}
\begin{proof}
\leanok
\uses{prop:isometry-invariance-riemannian-distance,prop:euclidean-riemannian-distance,ex:euclidean-metric}
\sketch Since $de_x=e$, inner-product preservation gives $e^*\bar g=\bar g$; bijective linearity makes $e$ a diffeomorphism. Apply Proposition \ref{prop:isometry-invariance-riemannian-distance}, followed by Proposition \ref{prop:euclidean-riemannian-distance}.
\end{proof}

Local isometries need not preserve global Riemannian distance.
\uses{prop:local-isometry-distance}

\begin{lemma}[Smoothness of the Energy Function]
\label{lem:metric-energy-smooth}
\lean{LeeLib.Ch02.RiemannianMetric.contMDiff_innerAt_totalSpace}
\uses{def:riemannian-metric}
\leanok



For a Riemannian manifold, the function
\[
E:TM\to\mathbb R,\qquad E(x,v)=\langle v,v\rangle_g,
\]
is smooth. Consequently $(x,v)\mapsto|v|_g=E(x,v)^{1/2}$ is continuous.
\end{lemma}
\begin{proof}
\leanok
\uses{def:riemannian-metric}
\sketch Pull $TM$ back along its bundle projection and pair the tautological section with itself using the smooth metric; this is $E$. Compose the nonnegative function $E$ with the continuous square root.
\end{proof}

\begin{lemma}[Uniform Comparison over a Compact Set]
\label{lem:norm-comparison-compact-core}
\lean{LeeLib.Ch02.exists_forall_sqrt_comparison}
\leanok


Let $V$ be finite-dimensional normed, let $q_x$ be a bilinear form on $V$ for every $x\in V$, let $K\subseteq V$ be compact, and put $S=\{v:|v|=1\}$. Assume
\begin{enumerate}
\item[(i)] $(x,v)\mapsto q_x(v,v)$ is continuous on $K\times S$;
\item[(ii)] $q_x(v,v)>0$ for all $x\in K$ and $v\ne0$.
\end{enumerate}
Then some $c,C>0$ satisfy
\[
c|v|\le q_x(v,v)^{1/2}\le C|v|,
\qquad x\in K, v\in V.
\]
\end{lemma}
\begin{proof}
\leanok
\sketch On compact $K\times S$, the positive continuous function $q_x(v,v)$ has bounds $0<c_0\le q_x(v,v)\le C_0$. Set $c=c_0^{1/2}$ and $C=\max(C_0,1)^{1/2}$. For $v\ne0$, apply these bounds to $u=v/|v|$ and use $q_x(v,v)=|v|^2q_x(u,u)$; $v=0$ is immediate. If $K$ or $S$ is empty, the assertion is vacuous.
\end{proof}

The compact comparison applies to a metric read in a coordinate chart, whose domain need not be all of $\mathbb R^n$; this is the form used in Lemma \ref{lem:distance-comparison-neighborhood}.

\begin{lemma}[Norm Comparison on $\mathbb{R}^n$]
\label{lem:riemannian-euclidean-norm-comparison-flat}
\lean{LeeLib.Ch02.RiemannianMetric.exists_forall_normAt_le_euclidean}
\uses{lem:norm-comparison-compact-core, lem:metric-energy-smooth, def:riemannian-metric, ex:euclidean-metric}
\leanok



Let $g$ be Riemannian on $\mathbb R^n$ and $\widetilde g$ Euclidean. For every compact $K\subseteq\mathbb R^n$, some $c,C>0$ satisfy
\[
c|v|_{\widetilde g}\le|v|_g\le C|v|_{\widetilde g},
\qquad x\in K, v\in T_x\mathbb R^n.
\]
\end{lemma}
\begin{proof}
\leanok
\uses{lem:norm-comparison-compact-core, lem:metric-energy-smooth, ex:euclidean-metric}
\sketch Under $T\mathbb R^n\cong\mathbb R^n\times\mathbb R^n$, Lemma \ref{lem:metric-energy-smooth} gives continuity of $(x,v)\mapsto g_x(v,v)$, and positivity is that of $g$. Apply Lemma \ref{lem:norm-comparison-compact-core} to $q_x=g_x$ and use $|v|_{\widetilde g}=|v|$.
\end{proof}

\begin{lemma}
\label{lem:riemannian-euclidean-norm-comparison}
\lean{LeeLib.Ch02.RiemannianMetric.exists_forall_normAt_le_euclidean_opens}
\uses{def:riemannian-metric, ex:euclidean-metric, lem:norm-comparison-compact-core, lem:metric-energy-smooth, lem:riemannian-euclidean-norm-comparison-flat}
\leanok
\dcref{ch2:2.53}



Let $g$ be Riemannian on an open $W\subseteq\mathbb R^n$ or $\mathbb R^n_+$, and let $\widetilde g$ be Euclidean. For every compact $K\subseteq W$, some $c,C>0$ satisfy
\begin{equation}
c|v|_{\widetilde g}\le|v|_g\le C|v|_{\widetilde g},
\qquad x\in K, v\in T_x\mathbb R^n.
\tag{2.21}
\end{equation}
\end{lemma}
\begin{proof}
\leanok
\uses{lem:norm-comparison-compact-core, lem:metric-energy-smooth, ex:euclidean-metric}
\sketch Lemma \ref{lem:metric-energy-smooth} makes $F(x,v)=|v|_g$ continuous. In the global trivialization $TW\cong W\times\mathbb R^n$, the set
\[
L=\{(x,v):x\in K,\ |v|_{\widetilde g}=1\}=K\times S^{n-1}
\]
is compact, so $0<c\le|v|_g\le C$ on $L$. Homogeneity gives (2.21) for nonzero $v$, and zero is immediate. Equivalently, this is Lemma \ref{lem:norm-comparison-compact-core} for $q_x=g_x$.
\end{proof}

\begin{lemma}\label{lem:distance-comparison-neighborhood}\dcref{ch2:2.54}
\lean{LeeLib.Ch02.RiemannianMetric.exists_coordinate_nbhd_riemannianEDist_comparison}
\uses{def:riemannian-metric, def:riemannian-distance, lem:riemannian-distance-locally-small, lem:riemannian-distance-small-implies-nearby}
\leanok



Let $(M,g)$ be Riemannian with or without boundary, let $U\subseteq M$ be open, and let $p\in U$. There is a coordinate neighborhood $V\subset U$ and constants $C,D>0$ such that:
\begin{enumerate}
\item[(a)] $d_g(p,q)\le C d_{\bar g}(p,q)$ for $q\in V$, where $\bar g$ is Euclidean in the coordinates;
\item[(b)] $d_g(p,q)\ge D$ for $q\in M\setminus V$.
\end{enumerate}
\begin{proof}
  \leanok

\sketch Identify a coordinate neighborhood $W\subset U$ with an open subset of $\mathbb R^n$ or $\mathbb R^n_+$. Choose $\varepsilon>0$ so that the closed ball or half-ball $K$ of radius $\varepsilon$ lies in $W$, and let $V$ be its relative interior. From (2.21), choose $c,C>0$ on $K$. Any admissible curve in $V$ satisfies
\[
cL_{\bar g}(\gamma)\le L_g(\gamma)\le CL_{\bar g}(\gamma).\tag{2.22}
\]
For $q\in V$, the coordinate line segment gives $d_g(p,q)\le C d_{\bar g}(p,q)$. If $q\notin V$, any curve from $p$ to $q$ first reaches the relative boundary at a point of Euclidean distance $\varepsilon$ from $p$; its initial segment has $g$-length at least $c\varepsilon$. Taking the infimum gives (b) with $D=c\varepsilon$.
\end{proof}
\end{lemma}

\begin{theorem}[Riemannian Manifolds as Metric Spaces]
\label{thm:riemannian-manifolds-metric-spaces}
\lean{LeeLib.Ch02.MetricSpace.ofRiemannianMetric}
\uses{def:riemannian-distance, prop:riemannian-distance-finite, lem:riemannian-distance-symm, lem:riemannian-distance-triangle, lem:riemannian-distance-locally-small, lem:riemannian-distance-small-implies-nearby}
\leanok
\dcref{ch2:2.55}



Every connected Riemannian manifold with or without boundary carries the metric-space structure determined below. Its distance is $d_g$ by Lemma \ref{lem:riemannian-metric-space-dist}, and its topology is the manifold topology by Theorem \ref{thm:riemannian-distance-topology}.
\end{theorem}
\begin{proof}
\leanok
\uses{prop:riemannian-distance-finite,lem:riemannian-distance-symm,lem:riemannian-distance-triangle,lem:riemannian-distance-locally-small,lem:riemannian-distance-small-implies-nearby}
\sketch Nonnegativity and $d_g(p,p)=0$ follow from the definition; symmetry is Lemma \ref{lem:riemannian-distance-symm} and the triangle inequality is Lemma \ref{lem:riemannian-distance-triangle}. If $p\ne q$, choose an open neighborhood $U$ of $p$ excluding $q$; Lemma \ref{lem:riemannian-distance-small-implies-nearby} gives $c>0$ whose $d_g$-ball lies in $U$, hence $d_g(p,q)\ge c$. Proposition \ref{prop:riemannian-distance-finite} makes the extended distance real-valued.
\end{proof}

\begin{lemma}[The Metric Is the Riemannian Distance]
\label{lem:riemannian-metric-space-dist}
\lean{LeeLib.Ch02.MetricSpace.ofRiemannianMetric_dist}
\uses{thm:riemannian-manifolds-metric-spaces, def:riemannian-distance, prop:riemannian-distance-finite}
\leanok



For the metric space of Theorem \ref{thm:riemannian-manifolds-metric-spaces},
\[
\operatorname{dist}(p,q)=d_g(p,q).
\]
The right side from Definition \ref{def:riemannian-distance} is real-valued by Proposition \ref{prop:riemannian-distance-finite}.
\end{lemma}
\begin{proof}
\leanok
\uses{def:riemannian-distance,prop:riemannian-distance-finite}
\sketch The structure in Theorem \ref{thm:riemannian-manifolds-metric-spaces} takes the finite real value of the extended distance from Definition \ref{def:riemannian-distance}; finiteness is Proposition \ref{prop:riemannian-distance-finite}. The identity is therefore definitional.
\end{proof}

\begin{theorem}[The Metric Topology Is the Manifold Topology]
\label{thm:riemannian-distance-topology}
\lean{LeeLib.Ch02.MetricSpace.ofRiemannianMetric_topology}
\uses{thm:riemannian-manifolds-metric-spaces, lem:riemannian-distance-locally-small, lem:riemannian-distance-small-implies-nearby}
\leanok



For a connected Riemannian manifold with or without boundary, the topology induced by the metric $d_g$ of Theorem \ref{thm:riemannian-manifolds-metric-spaces} equals the manifold topology.
\end{theorem}
\begin{proof}
\leanok
\uses{lem:riemannian-distance-locally-small,lem:riemannian-distance-small-implies-nearby}
\sketch If $U$ is manifold-open and $p\in U$, Lemma \ref{lem:riemannian-distance-small-implies-nearby} puts a $d_g$-ball about $p$ inside $U$. Conversely, if a $d_g$-ball of radius $\delta$ about $p$ lies in $U'$, Lemma \ref{lem:riemannian-distance-locally-small} gives a manifold neighborhood $V$ of $p$ with $d_g(p,q)<\delta$ on $V$, so $V\subseteq U'$.
\end{proof}

\begin{proposition}[The Euclidean Distance]
\label{prop:euclidean-riemannian-distance}
\lean{LeeLib.Ch02.euclideanMetric_riemannianEDist}
\uses{def:riemannian-distance, ex:euclidean-metric}
\leanok



On $\mathbb R^n$ with the Euclidean metric $\bar g$ of Example \ref{ex:euclidean-metric},
\[
d_{\bar g}(x,y)=|x-y|,
\qquad x,y\in\mathbb R^n.
\]
This is Definition \ref{def:riemannian-distance} in the Euclidean case; its real-valued form is Lemma \ref{lem:euclidean-riemannian-distance-real}.
\end{proposition}
\begin{proof}
\leanok
\uses{def:riemannian-distance,ex:euclidean-metric}
\sketch The segment $\gamma(t)=x+t(y-x)$ has length $|y-x|$, giving $d_{\bar g}(x,y)\le|x-y|$. For any $C^1$ competitor,
\[
|y-x|=\left|\int_0^1\gamma'(t)\,dt\right|\le\int_0^1|\gamma'(t)|\,dt=L_{\bar g}(\gamma),
\]
so the reverse inequality follows on taking the infimum.
\end{proof}

\begin{lemma}[The Euclidean Distance, real-valued form]
\label{lem:euclidean-riemannian-distance-real}
\lean{LeeLib.Ch02.euclideanMetric_riemannianEDist_toReal}
\uses{prop:euclidean-riemannian-distance}
\leanok



With the notation of Proposition \ref{prop:euclidean-riemannian-distance}, $d_{\bar g}(x,y)$ is finite and equals $|x-y|$ as a real number.
\end{lemma}
\begin{proof}
\leanok
\uses{prop:euclidean-riemannian-distance}
\sketch Proposition \ref{prop:euclidean-riemannian-distance} identifies the distance with the finite real number $|x-y|$.
\end{proof}

A connected Riemannian manifold is metrically complete if every Cauchy sequence converges. A subset $A\subseteq M$ is bounded if $d_g(p,q)\le C$ for all $p,q\in A$ and some $C$; its diameter is
\[
\operatorname{diam}(A)=\sup\{d_g(p,q):p,q\in A\}.
\]
Every compact connected Riemannian manifold with or without boundary has finite diameter.
\section{Pseudo-Riemannian Metrics}

For a symmetric bilinear form $q$ on a finite-dimensional vector space $V$, define
\[
\widetilde q\colon V\longrightarrow V^*,
\qquad \widetilde q(v)(w)=q(v,w).
\]
The form $q$ is \emph{nondegenerate} when $\widetilde q$ is an isomorphism.

\begin{lemma}
\label{lem:nondegenerate-symmetric-bilinear}
\lean{LeeLib.Ch02.nondegenerate_tfae}
\leanok


\dcref{ch2:2.56}
For a symmetric bilinear form $q$ on a finite-dimensional vector space $V$, the
following conditions are equivalent:
\begin{enumerate}
\item[(a)] $q$ is nondegenerate.
\item[(b)] For each nonzero $v\in V$, some $w\in V$ satisfies $q(v,w)\ne0$.
\item[(c)] In any expression $q=q_{ij}e^ie^j$ relative to a basis $(e^j)$ of
$V^*$, the matrix $(q_{ij})$ is invertible.
\end{enumerate}
\end{lemma}

\begin{proof}
\leanok
\sketch Conditions (a) and (b) respectively say that $\widetilde q$ is an
isomorphism and is injective; these are equivalent because $V$ and $V^*$ have the
same finite dimension. If $(e_i)$ is dual to $(e^j)$, then
$\widetilde q(e_i)(e_j)=q_{ij}$, so $(q_{ij})$ is the matrix of $\widetilde q$,
proving the equivalence with (c).
\end{proof}

\begin{proof}
\leanok
\sketch This is the preceding lemma.
\end{proof}

A \emph{scalar product} is a nondegenerate symmetric bilinear form on a
finite-dimensional real vector space; an \emph{inner product} is a positive
definite scalar product. A \emph{scalar product space} is a vector space equipped
with such a form. Vectors $v,w$ are \emph{orthogonal} when
$\langle v,w\rangle=0$, and the conventional indefinite norm is
\[
|v|=|\langle v,v\rangle|^{1/2}.
\]
A nonzero vector may be null, so this function is not a norm in the usual sense.

\begin{proposition}
\label{prop:polarization-scalar-product}


\dcref{ch2:2.58}
\uses{lem:polarization-identity}
The polarization identity (2.2) holds for every scalar product.
\end{proposition}

\begin{proof}
\leanok
\sketch Bilinearity and symmetry give
\[
\langle v+w,v+w\rangle-\langle v-w,v-w\rangle=4\langle v,w\rangle,
\]
without using positive definiteness.
\end{proof}

For a subspace $S\subseteq V$, write
$S^\perp=\{v\in V:q(v,s)=0\text{ for all }s\in S\}$.

\begin{lemma}
\label{lem:orthogonal-complement}
\lean{LeeLib.Ch02.finrank_add_finrank_orthogonal_eq_finrank}
\leanok


\dcref{ch2:2.59}
If $(V,q)$ is a finite-dimensional scalar product space and $S\subseteq V$ is a
subspace, then
\[
\dim S+\dim S^\perp=\dim V.
\]
\end{lemma}

\begin{proof}
\leanok
\sketch The map
$\Phi\colon V\to S^*$, $\Phi(v)=\widetilde q(v)|_S$, has kernel $S^\perp$.
Every element of $S^*$ extends to $V^*$ and $\widetilde q$ is onto, so $\Phi$ is
surjective. Rank--nullity yields the formula.
\end{proof}

\begin{lemma}
\label{lem:orthogonal-complement-involution}
\lean{LeeLib.Ch02.orthogonal_orthogonal_eq_self}
\uses{lem:orthogonal-complement}
\leanok


\dcref{ch2:2.59}

For every subspace $S$ of a finite-dimensional scalar product space,
\[
(S^\perp)^\perp=S.
\]
\end{lemma}

\begin{proof}
\leanok
\uses{lem:orthogonal-complement}
\sketch The definition gives $S\subseteq(S^\perp)^\perp$, and Lemma
\ref{lem:orthogonal-complement} shows that the two spaces have equal dimension.
\end{proof}

An ordered tuple $(v_1,\ldots,v_k)$ is \emph{orthonormal} when
$\langle v_i,v_j\rangle=\pm\delta_{ij}$. A subspace is \emph{nondegenerate} when
the restricted scalar product is nondegenerate. A tuple is \emph{nondegenerate}
when each prefix $(v_1,\ldots,v_j)$ spans a nondegenerate $j$-dimensional
subspace.

\begin{lemma}\label{lem:nondegenerate-subspace}
\lean{LeeLib.Ch02.restrict_nondegenerate_tfae}
\uses{lem:orthogonal-complement-involution}
\leanok


\dcref{ch2:2.60}

Let $S$ be a subspace of a finite-dimensional scalar product space $V$. The
following are equivalent:
\begin{enumerate}[label=(\alph*)]
\item $S$ is nondegenerate.
\item $S^\perp$ is nondegenerate.
\item $S\cap S^\perp=\{0\}$.
\item $V=S\oplus S^\perp$.
\end{enumerate}
\end{lemma}

\begin{proof}
\leanok
\uses{lem:orthogonal-complement,lem:orthogonal-complement-involution}
\sketch The radical of $q|_S$ is $S\cap S^\perp$, proving (a)$\Leftrightarrow$(c).
Together with $\dim S+\dim S^\perp=\dim V$, trivial intersection is equivalent to
(d). Applying the same argument to $S^\perp$ and using
$(S^\perp)^\perp=S$ proves (b)$\Leftrightarrow$(d).
\end{proof}

\begin{proof}
  \uses{lem:orthogonal-complement, lem:orthogonal-complement-involution}
\leanok
\sketch This is the preceding lemma.
\end{proof}

\begin{lemma}[Completion of Nondegenerate Bases]\label{lem:completion-nondegenerate-bases}
\lean{LeeLib.Ch02.exists_isNondegenerateTuple_basis}
\uses{lem:nondegenerate-subspace}
\leanok


\dcref{ch2:2.62}

Let $(V,q)$ be an $n$-dimensional scalar product space and
$(v_1,\ldots,v_k)$ a nondegenerate $k$-tuple with $0\leq k<n$. It extends to a
nondegenerate basis $(v_1,\ldots,v_n)$ of $V$.
\end{lemma}

\begin{proof}
  \uses{lem:nondegenerate-subspace, lem:orthogonal-complement, prop:polarization-scalar-product, exer:polarization-scalar-product}
\leanok

\sketch Put $S=\operatorname{span}(v_1,\ldots,v_k)$. Then
$V=S\oplus S^\perp$, and $S^\perp$ is nonzero and nondegenerate. It contains a
vector $v_{k+1}$ with $q(v_{k+1},v_{k+1})\ne0$: otherwise polarization would make
$q|_{S^\perp}$ identically zero. Orthogonality and the nonzero self-pairing show
$S\oplus\operatorname{span}(v_{k+1})$ is nondegenerate. Iterate this construction
until the dimension is $n$.
\end{proof}

\begin{proposition}[Gram--Schmidt Algorithm for Scalar Products]
\label{prop:gram-schmidt-algorithm-scalar-products}
\lean{LeeLib.Ch02.exists_basis_isOrthonormal_prefixSpan_eq}
\leanok


\dcref{ch2:2.63}
Let $(V,q)$ be an $n$-dimensional scalar product space. Every nondegenerate basis
$(v_i)$ admits an orthonormal basis $(b_i)$ satisfying
\[
\operatorname{span}(b_1,\ldots,b_k)
=\operatorname{span}(v_1,\ldots,v_k)
\qquad(1\leq k\leq n).
\]
\end{proposition}

\begin{proof}
\leanok
\uses{lem:nondegenerate-symmetric-bilinear}
\sketch Assume $b_1,\ldots,b_k$ have been constructed with the required flag and
set
\[
z=v_{k+1}-\sum_{i=1}^k
\frac{\langle v_{k+1},b_i\rangle}{\langle b_i,b_i\rangle}b_i.
\]
The denominators are $\pm1$, $z$ is orthogonal to the preceding span, and adjoining
$z$ reproduces the next prefix. If $z=0$, the next prefix has dimension $k$; if
$\langle z,z\rangle=0$, then $z$ is orthogonal to the entire next prefix. Both
contradict its nondegeneracy. Therefore
\[
b_{k+1}=\frac{z}{\sqrt{|\langle z,z\rangle|}}
\]
is defined, orthonormal to the preceding vectors, and preserves the prefix span.
Induction gives the desired basis.
\end{proof}

\begin{proposition}
\label{prop:symmetric-bilinear-orthogonal-basis}
\lean{LinearMap.BilinForm.exists_orthogonal_basis}

\mathlibok
Let $q$ be a symmetric bilinear form on a finite-dimensional vector space $V$
over a field in which $2$ is invertible. Then $V$ has a $q$-orthogonal basis
$(b_i)$, meaning $q(b_i,b_j)=0$ for $i\ne j$.
\end{proposition}

\begin{proposition}
\label{prop:scalar-product-orthonormal-basis}
\lean{LeeLib.Ch02.exists_basis_isOrthonormal}
\uses{prop:symmetric-bilinear-orthogonal-basis}
\leanok



Every scalar product space has an orthonormal basis.
\end{proposition}

\begin{proof}
\leanok
\uses{prop:symmetric-bilinear-orthogonal-basis,lem:nondegenerate-symmetric-bilinear}
\sketch Choose a $q$-orthogonal basis $(b_i)$. Nondegeneracy implies
$q(b_i,b_i)\ne0$ for every $i$; otherwise $b_i$ would be orthogonal to the entire
basis. Then
\[
e_i=\frac{b_i}{\sqrt{|q(b_i,b_i)|}}
\]
is an orthonormal basis.
\end{proof}

\begin{corollary}
\label{cor:scalar-product-signature-form}
\lean{LeeLib.Ch02.exists_basis_isOrthonormal_signature}
\uses{prop:scalar-product-orthonormal-basis}
\leanok


\dcref{ch2:2.64}

If $(V,q)$ is an $n$-dimensional scalar product space, some basis $(\beta^i)$ of
$V^*$ and nonnegative integers $r,s$ with $r+s=n$ satisfy
\begin{equation}
q=(\beta^1)^2+\cdots+(\beta^r)^2
 -(\beta^{r+1})^2-\cdots-(\beta^{r+s})^2. \tag{2.23}
\end{equation}
\end{corollary}

\begin{proof}
\leanok
\uses{prop:scalar-product-orthonormal-basis}
\sketch For an orthonormal basis $(b_i)$ with dual basis $(\beta^i)$,
\[
q=\sum_iq(b_i,b_i)(\beta^i)^2,
\qquad q(b_i,b_i)\in\{1,-1\}.
\]
Reorder the positive entries first; their counts $r,s$ sum to $n$.
\end{proof}

\begin{proposition}[Sylvester's Law of Inertia]
\label{prop:sylvesters-law-inertia}
\lean{LeeLib.Ch02.IsOrthonormal.isGreatest_ncard}
\uses{cor:scalar-product-signature-form}
\leanok


\dcref{ch2:2.65}

If a finite-dimensional scalar product $q$ has the normal form (2.23), then $r$
is the maximum dimension of a subspace on which $q$ is positive definite.
Consequently $r$ and $s$ are independent of the chosen basis.
\end{proposition}

\begin{proof}
\leanok
\uses{cor:scalar-product-signature-form}
\sketch Let $(b_i)$ be the orthonormal basis dual to $(\beta^i)$. The span of
$b_1,\ldots,b_r$ is positive definite, so the maximum is at least $r$. On
$N=\operatorname{span}(b_{r+1},\ldots,b_{r+s})$, the form is nonpositive. Any
positive-definite subspace $W$ has $W\cap N=\{0\}$, hence
$\dim W+s\leq r+s$ and $\dim W\leq r$. Thus $r$ has a basis-free
characterization, and $s=\dim V-r$.
\end{proof}

\begin{lemma}[A Scalar Product Has Trivial Radical]
\label{lem:radical-scalar-product-trivial}
\lean{LeeLib.Ch02.radical_toQuadraticMap_eq_bot}
\uses{lem:nondegenerate-symmetric-bilinear}
\leanok



The radical of a scalar product $q$ is $\{0\}$; explicitly, if $q(v,v)=0$ and
$w\mapsto q(v,w)+q(w,v)$ vanishes identically, then $v=0$.
\end{lemma}

\begin{proof}
\leanok
\uses{lem:nondegenerate-symmetric-bilinear}
\sketch Symmetry turns the polar form into $w\mapsto2q(v,w)$. Its vanishing makes
$v$ orthogonal to all of $V$, so nondegeneracy gives $v=0$.
\end{proof}

\begin{corollary}[$r + s = n$, without a basis]
\label{cor:scalar-product-signature-sum}
\lean{LeeLib.Ch02.sigPos_add_sigNeg_eq_finrank}
\uses{lem:radical-scalar-product-trivial, prop:sylvesters-law-inertia}
\leanok



If an $n$-dimensional scalar product space has signature $(r,s)$, then
$r+s=n$.
\end{corollary}

\begin{proof}
\leanok
\uses{lem:radical-scalar-product-trivial,prop:sylvesters-law-inertia}
\sketch The positive index, negative index, and radical dimension of a quadratic
form sum to $\dim V$. Here they are $r$, $s$, and $0$.
\end{proof}

The integer $s$ in (2.23) is the \emph{index} of $q$, and $(r,s)$ is its
\emph{signature}.

\begin{definition}[Smooth Field of Scalar Products]
\label{def:smooth-field-scalar-products}
\lean{Bundle.ContMDiffPseudoMetric}
\uses{lem:nondegenerate-symmetric-bilinear}
\leanok



Let $\pi\colon V\to B$ be a smooth real vector bundle. A \emph{smooth field of
scalar products} is a family of bilinear forms $g_b$ on $V_b$ such that
\begin{enumerate}[label=(\alph*)]
\item $g_b$ is symmetric;
\item for every nonzero $v\in V_b$, some $w\in V_b$ satisfies $g_b(v,w)\ne0$;
\item $b\mapsto g_b$ is a smooth section of the bundle of bilinear forms on $V$.
\end{enumerate}
\end{definition}

\begin{lemma}[The Fibres are Scalar Product Spaces]
\label{lem:pseudo-metric-fibre-scalar-product}
\lean{Bundle.ContMDiffPseudoMetric.bilin_nondegenerate}
\uses{def:smooth-field-scalar-products, lem:nondegenerate-symmetric-bilinear}
\leanok



If $g$ is a smooth field of scalar products on a bundle with finite-dimensional
fibres, then $(V_b,g_b)$ is a scalar product space for every $b\in B$.
\end{lemma}

\begin{proof}
\leanok
\uses{def:smooth-field-scalar-products,lem:nondegenerate-symmetric-bilinear}
\sketch Symmetry and bilinearity follow from the definition, while its witness
condition is equivalent to nondegeneracy by Lemma
\ref{lem:nondegenerate-symmetric-bilinear}.
\end{proof}

\begin{lemma}[Smoothness of the Pairing of Smooth Sections]
\label{lem:pseudo-metric-pairing-smooth}
\lean{Bundle.ContMDiffPseudoMetric.contMDiffOn_pairing, Bundle.ContMDiffPseudoMetric.contMDiff_pairing}
\uses{def:smooth-field-scalar-products}
\leanok



Let $g$ be a smooth field of scalar products on $V\to B$, let $f\colon N\to B$
be smooth on $s\subseteq N$, and let $v,w$ be sections along $f$ that are smooth
on $s$. Then
\[
x\longmapsto g_{f(x)}(v_x,w_x)
\]
is smooth on $s$. In particular, this holds for globally smooth sections along a
smooth map.
\end{lemma}

\begin{proof}
\leanok
\uses{def:smooth-field-scalar-products}
\sketch In a local trivialization over $f(x_0)$, the field $g$ becomes a smooth map
to the fixed space of bilinear forms on the model fibre, and $v,w$ become smooth
model-fibre-valued maps. Compose these maps with the smooth evaluation
$(\beta,y,z)\mapsto\beta(y,z)$.
\end{proof}

\begin{definition}[Pseudo-Riemannian Metric]
\label{def:pseudo-riemannian-metric}
\lean{LeeLib.Ch02.PseudoRiemannianMetric}
\uses{def:smooth-field-scalar-products}
\leanok


\dcref{ch2:2.7}

A \emph{pseudo-Riemannian metric} on a smooth manifold $M$ is a smooth field of
scalar products on $TM$, equivalently a smooth symmetric $2$-tensor $g$ whose
restriction $g_p$ to every $T_pM$ is nondegenerate. A \emph{pseudo-Riemannian
manifold} is a pair $(M,g)$ with such a metric.
\end{definition}

Constant signature is imposed separately in Definition
\ref{def:pseudo-metric-has-signature}.

\begin{proposition}[Every Riemannian Metric is Pseudo-Riemannian]
\label{prop:riemannian-is-pseudo-riemannian}
\lean{LeeLib.Ch02.RiemannianMetric.toPseudoRiemannianMetric}
\uses{def:riemannian-metric, def:pseudo-riemannian-metric}
\leanok


\dcref{ch2:2.7}

Every Riemannian metric is pseudo-Riemannian.
\end{proposition}

\begin{proof}
\leanok
\uses{def:riemannian-metric,def:pseudo-riemannian-metric,lem:nondegenerate-symmetric-bilinear}
\sketch For $v\ne0$, positive definiteness gives $g_p(v,v)>0$. Taking $w=v$
verifies the witness condition for nondegeneracy; smoothness and symmetry are part
of the Riemannian metric definition.
\end{proof}

\begin{definition}[Signature of a Pseudo-Riemannian Metric at a Point]
\label{def:pseudo-metric-signature-at}
\lean{LeeLib.Ch02.PseudoRiemannianMetric.sigPosAt, LeeLib.Ch02.PseudoRiemannianMetric.sigNegAt}
\uses{def:pseudo-riemannian-metric, lem:pseudo-metric-fibre-scalar-product, cor:scalar-product-signature-form, prop:sylvesters-law-inertia}
\leanok



For a pseudo-Riemannian manifold $(M,g)$ and $p\in M$, let $(r_p,s_p)$ be the
signature of the scalar product $(T_pM,g_p)$. This is the \emph{signature of $g$
at $p$}, and $s_p$ is its \emph{index}. Equivalently, $r_p$ and $s_p$ are the
maximal dimensions of subspaces on which $g_p$ is respectively positive and
negative definite.
\end{definition}

\begin{proposition}[The Signature at a Point Sums to the Dimension]
\label{prop:pseudo-metric-signature-sum}
\lean{LeeLib.Ch02.PseudoRiemannianMetric.sigPosAt_add_sigNegAt}
\uses{def:pseudo-metric-signature-at, lem:pseudo-metric-fibre-scalar-product, cor:scalar-product-signature-sum}
\leanok



If $(r_p,s_p)$ is the signature of a pseudo-Riemannian metric at $p$, then
\[
r_p+s_p=\dim T_pM.
\]
\end{proposition}

\begin{proof}
\leanok
\uses{lem:pseudo-metric-fibre-scalar-product,cor:scalar-product-signature-sum}
\sketch Apply Corollary \ref{cor:scalar-product-signature-sum} to $(T_pM,g_p)$.
\end{proof}

\begin{definition}[Constant Signature]
\label{def:pseudo-metric-has-signature}
\lean{LeeLib.Ch02.HasSignature}
\uses{def:pseudo-metric-signature-at}
\leanok



A pseudo-Riemannian metric $g$ \emph{has signature $(r,s)$} if its signature is
$(r,s)$ at every point. A pseudo-Riemannian manifold of signature $(r,s)$ is a
pair $(M,g)$ with this property.
\end{definition}

On an $n$-manifold, a constant signature satisfies $r+s=n$.

\subsection{The Gram--Schmidt Algorithm as an Explicit Recursion}

\begin{definition}[The Gram--Schmidt Recursion for a Bilinear Form]
\label{def:gram-schmidt-bilin}
\lean{LeeLib.Ch02.gramSchmidtBilin}
\leanok


For a bilinear form $B$ on a real vector space $V$ and an $n$-tuple
$v=(v_1,\ldots,v_n)$, define $u=\mathrm{GS}_B(v)$ recursively by
\[
u_j=v_j-\sum_{i<j}\frac{B(u_i,v_j)}{B(u_i,u_i)}u_i.
\]
The recursion is used under Proposition \ref{prop:gram-schmidt-bilin-spec}, which
ensures that every denominator is nonzero.
\end{definition}

\begin{definition}[The Normalized Gram--Schmidt Recursion]
\label{def:gram-schmidt-bilin-normed}
\lean{LeeLib.Ch02.gramSchmidtBilinNormed}
\uses{def:gram-schmidt-bilin}
\leanok



For $u=\mathrm{GS}_B(v)$, define $b=\widehat{\mathrm{GS}}_B(v)$ by
\[
b_j=\frac{u_j}{\sqrt{|B(u_j,u_j)|}}.
\]
Then $B(b_j,b_j)=\pm1$ whenever the denominator is nonzero.
\end{definition}

\begin{lemma}[The New Vector is Orthogonal to its Predecessors]
\label{lem:gram-schmidt-bilin-orthogonal-earlier}
\lean{LeeLib.Ch02.gramSchmidtBilin_apply_of_lt}
\uses{def:gram-schmidt-bilin}
\leanok



Let $u=\mathrm{GS}_B(v)$ and $j=k+1$. Suppose
\[
B(u_i,u_i)\ne0\quad(i\leq k),
\qquad B(u_i,u_{i'})=0\quad(i\ne i',\ i,i'\leq k).
\]
Then $B(u_c,u_j)=0$ for every $c<j$.
\end{lemma}

\begin{proof}
\leanok
\uses{def:gram-schmidt-bilin}
\sketch Expanding the recursion gives
\[
B(u_c,u_j)=B(u_c,v_j)-\sum_{i<j}
\frac{B(u_i,v_j)}{B(u_i,u_i)}B(u_c,u_i).
\]
Orthogonality kills every summand except $i=c$, and that term cancels
$B(u_c,v_j)$.
\end{proof}

\begin{proposition}[The Gram--Schmidt Recursion: Simultaneous Induction]
\label{prop:gram-schmidt-bilin-spec}
\lean{LeeLib.Ch02.gramSchmidtBilin_spec}
\uses{def:gram-schmidt-bilin, lem:gram-schmidt-bilin-orthogonal-earlier, lem:nondegenerate-symmetric-bilinear}
\leanok



Let $B$ be symmetric and $(v_1,\ldots,v_n)$ nondegenerate. For
$u=\mathrm{GS}_B(v)$ and every $k\leq n$:
\begin{enumerate}[label=(\alph*)]
\item $B(u_i,u_i)\ne0$ for every $i\leq k$;
\item $B(u_i,u_{i'})=0$ for distinct $i,i'\leq k$;
\item $\operatorname{span}(u_1,\ldots,u_j)=
\operatorname{span}(v_1,\ldots,v_j)$ for every $j\leq k$.
\end{enumerate}
\end{proposition}

\begin{proof}
\leanok
\uses{def:gram-schmidt-bilin,lem:gram-schmidt-bilin-orthogonal-earlier,lem:nondegenerate-symmetric-bilinear}
\sketch Induct on $k$. For $j=k+1$, the preceding lemma makes $u_j$ orthogonal
to $S=\operatorname{span}(u_1,\ldots,u_k)$. The recursion gives
$u_j=v_j-\sigma$ for some $\sigma\in S$, so adjoining $u_j$ or $v_j$ produces
the same next prefix $T$. Nondegeneracy of the input flag gives $v_j\notin S$,
hence $u_j\ne0$. If $B(u_j,u_j)=0$, then $u_j$ is orthogonal to every vector in
$T=S+\operatorname{span}(u_j)$, contradicting nondegeneracy of $T$. This proves
the new denominator condition, orthogonality, and span identity simultaneously.
\end{proof}

\begin{lemma}[The Denominators Never Vanish]
\label{lem:gram-schmidt-bilin-nonvanishing}
\lean{LeeLib.Ch02.gramSchmidtBilin_apply_self_ne_zero}
\uses{prop:gram-schmidt-bilin-spec}
\leanok



Under the hypotheses of Proposition \ref{prop:gram-schmidt-bilin-spec},
$B(u_j,u_j)\ne0$ for every $j$.
\end{lemma}

\begin{proof}
\leanok
\uses{prop:gram-schmidt-bilin-spec}
\sketch Take $k=n$ in part (a) of Proposition
\ref{prop:gram-schmidt-bilin-spec}.
\end{proof}

\begin{lemma}[The Unnormalized Vectors are Orthogonal]
\label{lem:gram-schmidt-bilin-orthogonal}
\lean{LeeLib.Ch02.gramSchmidtBilin_apply_ne}
\uses{prop:gram-schmidt-bilin-spec}
\leanok



With the same hypotheses, $B(u_i,u_j)=0$ whenever $i\ne j$.
\end{lemma}

\begin{proof}
\leanok
\uses{prop:gram-schmidt-bilin-spec}
\sketch Take $k=n$ in part (b) of Proposition
\ref{prop:gram-schmidt-bilin-spec}.
\end{proof}

\begin{lemma}[The Unnormalized Vectors Reproduce the Flag]
\label{lem:gram-schmidt-bilin-flag}
\lean{LeeLib.Ch02.prefixSpan_gramSchmidtBilin}
\uses{prop:gram-schmidt-bilin-spec}
\leanok



With the same hypotheses,
\[
\operatorname{span}(u_1,\ldots,u_j)=\operatorname{span}(v_1,\ldots,v_j)
\qquad(j\leq n).
\]
\end{lemma}

\begin{proof}
\leanok
\uses{prop:gram-schmidt-bilin-spec}
\sketch Take $k=n$ in part (c) of Proposition
\ref{prop:gram-schmidt-bilin-spec}.
\end{proof}

\begin{lemma}[The Normalizing Scalar is Nonzero]
\label{lem:gram-schmidt-bilin-normalizer-ne-zero}
\lean{LeeLib.Ch02.gramSchmidtBilin_normalizer_ne_zero}
\uses{lem:gram-schmidt-bilin-nonvanishing, def:gram-schmidt-bilin-normed}
\leanok



For every $j$, the scalar $(\sqrt{|B(u_j,u_j)|})^{-1}$ is nonzero.
\end{lemma}

\begin{proof}
\leanok
\uses{lem:gram-schmidt-bilin-nonvanishing}
\sketch Nonvanishing of $B(u_j,u_j)$ makes $\sqrt{|B(u_j,u_j)|}$ strictly
positive.
\end{proof}

\begin{proposition}[The Recursion Produces an Orthonormal Tuple]
\label{prop:gram-schmidt-bilin-orthonormal}
\lean{LeeLib.Ch02.gramSchmidtBilin_isOrthonormal}
\uses{def:gram-schmidt-bilin-normed, lem:gram-schmidt-bilin-nonvanishing, lem:gram-schmidt-bilin-orthogonal}
\leanok



For symmetric $B$ and a nondegenerate tuple $v$, the tuple
$b=\widehat{\mathrm{GS}}_B(v)$ satisfies
\[
B(b_i,b_j)=0\quad(i\ne j),
\qquad B(b_j,b_j)=\pm1.
\]
\end{proposition}

\begin{proof}
\leanok
\uses{def:gram-schmidt-bilin-normed,lem:gram-schmidt-bilin-nonvanishing,lem:gram-schmidt-bilin-orthogonal}
\sketch Writing $b_j=\lambda_ju_j$ with
$\lambda_j=|B(u_j,u_j)|^{-1/2}$, orthogonality of the $u_j$ gives the off-diagonal
equalities, while
\[
B(b_j,b_j)=\frac{B(u_j,u_j)}{|B(u_j,u_j)|}\in\{1,-1\}.
\]
\end{proof}

\begin{lemma}[The Normalized Tuple Reproduces the Flag]
\label{lem:gram-schmidt-bilin-normed-flag}
\lean{LeeLib.Ch02.prefixSpan_gramSchmidtBilinNormed}
\uses{def:gram-schmidt-bilin-normed, lem:gram-schmidt-bilin-flag, lem:gram-schmidt-bilin-normalizer-ne-zero}
\leanok



For every $j\leq n$,
\[
\operatorname{span}(b_1,\ldots,b_j)=\operatorname{span}(v_1,\ldots,v_j).
\]
\end{lemma}

\begin{proof}
\leanok
\uses{lem:gram-schmidt-bilin-flag,lem:gram-schmidt-bilin-normalizer-ne-zero}
\sketch Each $b_i$ is a nonzero scalar multiple of $u_i$, so their prefix spans
coincide; apply Lemma \ref{lem:gram-schmidt-bilin-flag}.
\end{proof}

\begin{corollary}[Proposition \ref{prop:gram-schmidt-algorithm-scalar-products}, Recovered from the Recursion]
\label{cor:gram-schmidt-bilin-recovers}
\lean{LeeLib.Ch02.exists_isOrthonormal_prefixSpan_eq_of_gramSchmidt}
\uses{prop:gram-schmidt-bilin-orthonormal, lem:gram-schmidt-bilin-normed-flag}
\leanok



Every nondegenerate tuple for a symmetric bilinear form admits an orthonormal tuple
with the same prefix spans.
\end{corollary}

\begin{proof}
\leanok
\uses{prop:gram-schmidt-bilin-orthonormal,lem:gram-schmidt-bilin-normed-flag}
\sketch Take $b=\widehat{\mathrm{GS}}_B(v)$ and apply the preceding proposition and
lemma.
\end{proof}

\subsection{Smoothness of the Fibrewise Gram--Schmidt Algorithm}

\begin{lemma}[$(\sqrt{|\cdot|})^{-1}$ is Smooth Away from the Origin]
\label{lem:inv-sqrt-abs-smooth}
\lean{LeeLib.Ch02.contDiffAt_inv_sqrt_abs}
\leanok


For $t\in\mathbb R\setminus\{0\}$, the function
$s\mapsto(\sqrt{|s|})^{-1}$ is smooth at $t$.
\end{lemma}

\begin{proof}
\leanok
\sketch Near a positive $t$, $|s|=s$; near a negative $t$, $|s|=-s$. Compose
this local linear expression with the smooth square-root function on
$(0,\infty)$ and inversion on $\mathbb R\setminus\{0\}$.
\end{proof}

\begin{definition}[Pointwise Nondegenerate Smooth Family]
\label{def:local-nondegenerate-family}
\lean{LeeLib.Ch02.IsLocalNondegenerateOn}
\uses{def:smooth-field-scalar-products, lem:pseudo-metric-fibre-scalar-product}
\leanok



Let $g$ be a smooth field of scalar products on $V\to B$. A family of sections
$(X_1,\ldots,X_m)$ is \emph{nondegenerate on $u\subseteq B$} if
\begin{enumerate}[label=(\alph*)]
\item every $X_i$ is smooth on $u$;
\item $(X_1|_x,\ldots,X_m|_x)$ is a nondegenerate tuple in $(V_x,g_x)$ for
every $x\in u$.
\end{enumerate}
\end{definition}

\begin{definition}[The Fibrewise Gram--Schmidt Frame]
\label{def:pseudo-gram-schmidt-frame}
\lean{LeeLib.Ch02.pseudoGramSchmidtFrameAux, LeeLib.Ch02.pseudoGramSchmidtFrame}
\uses{def:gram-schmidt-bilin, def:gram-schmidt-bilin-normed, def:smooth-field-scalar-products}
\leanok



For a family $(X_1,\ldots,X_m)$ of sections, define fibrewise sections $U_j,E_j$
by
\[
U_j|_x=\mathrm{GS}_{g_x}(X_1|_x,\ldots,X_m|_x)_j,
\qquad
E_j|_x=\frac{U_j|_x}{\sqrt{|g_x(U_j|_x,U_j|_x)|}}.
\]
\end{definition}

\begin{lemma}[Nonvanishing of the Fibrewise Denominators]
\label{lem:pseudo-gram-schmidt-denominators}
\lean{LeeLib.Ch02.pseudoGramSchmidtFrameAux_apply_self_ne_zero}
\uses{def:pseudo-gram-schmidt-frame, def:local-nondegenerate-family, lem:gram-schmidt-bilin-nonvanishing, lem:pseudo-metric-fibre-scalar-product}
\leanok



If $(X_1,\ldots,X_m)$ is nondegenerate on $u$, then
\[
g_x(U_j|_x,U_j|_x)\ne0
\qquad(x\in u).
\]
\end{lemma}

\begin{proof}
\leanok
\uses{lem:gram-schmidt-bilin-nonvanishing,lem:pseudo-metric-fibre-scalar-product}
\sketch Apply Lemma \ref{lem:gram-schmidt-bilin-nonvanishing} in each scalar
product space $(V_x,g_x)$ to the nondegenerate tuple $(X_i|_x)$.
\end{proof}

\begin{proposition}[Orthonormality of the Gram--Schmidt Frame]
\label{prop:pseudo-gram-schmidt-frame-orthonormal}
\lean{LeeLib.Ch02.pseudoGramSchmidtFrame_isOrthonormal}
\uses{def:pseudo-gram-schmidt-frame, def:local-nondegenerate-family, prop:gram-schmidt-bilin-orthonormal, lem:pseudo-metric-fibre-scalar-product}
\leanok



Under the preceding hypotheses, $(E_1|_x,\ldots,E_m|_x)$ is orthonormal in
$(V_x,g_x)$ for every $x\in u$; equivalently,
$g_x(E_i|_x,E_j|_x)=\pm\delta_{ij}$.
\end{proposition}

\begin{proof}
\leanok
\uses{prop:gram-schmidt-bilin-orthonormal,lem:pseudo-metric-fibre-scalar-product}
\sketch Apply Proposition \ref{prop:gram-schmidt-bilin-orthonormal} fibrewise.
\end{proof}

\begin{lemma}[The Gram--Schmidt Frame Reproduces the Flag]
\label{lem:pseudo-gram-schmidt-frame-flag}
\lean{LeeLib.Ch02.prefixSpan_pseudoGramSchmidtFrame}
\uses{def:pseudo-gram-schmidt-frame, def:local-nondegenerate-family, lem:gram-schmidt-bilin-normed-flag, lem:pseudo-metric-fibre-scalar-product}



For $x\in u$ and $k\leq m$,
\[
\operatorname{span}(E_1|_x,\ldots,E_k|_x)
=\operatorname{span}(X_1|_x,\ldots,X_k|_x).
\]
\end{lemma}

\begin{proof}
\leanok
\uses{lem:gram-schmidt-bilin-normed-flag,lem:pseudo-metric-fibre-scalar-product}
\sketch Apply Lemma \ref{lem:gram-schmidt-bilin-normed-flag} in $(V_x,g_x)$.
\end{proof}

\begin{proposition}[Smoothness of the Unnormalized Gram--Schmidt Sections]
\label{prop:pseudo-gram-schmidt-frame-aux-smooth}
\lean{LeeLib.Ch02.contMDiffOn_pseudoGramSchmidtFrameAux}
\uses{def:pseudo-gram-schmidt-frame, def:local-nondegenerate-family, lem:pseudo-metric-pairing-smooth, lem:pseudo-gram-schmidt-denominators}



Each section $U_j$ is smooth on $u$.
\end{proposition}

\begin{proof}
\leanok
\uses{def:pseudo-gram-schmidt-frame,lem:pseudo-metric-pairing-smooth,lem:pseudo-gram-schmidt-denominators}
\sketch Use strong induction in the recursion
\[
U_j=X_j-\sum_{i<j}
\frac{g(U_i,X_j)}{g(U_i,U_i)}U_i.
\]
The inductive hypothesis and Lemma \ref{lem:pseudo-metric-pairing-smooth} make
each numerator and denominator smooth, and Lemma
\ref{lem:pseudo-gram-schmidt-denominators} makes every denominator nonvanishing.
Thus all coefficients and $U_j$ are smooth.
\end{proof}

\begin{proposition}[Smoothness of the Gram--Schmidt Frame]
\label{prop:pseudo-gram-schmidt-frame-smooth}
\lean{LeeLib.Ch02.contMDiffOn_pseudoGramSchmidtFrame}
\uses{prop:pseudo-gram-schmidt-frame-aux-smooth, lem:inv-sqrt-abs-smooth, lem:pseudo-metric-pairing-smooth, lem:pseudo-gram-schmidt-denominators, def:pseudo-gram-schmidt-frame}
\leanok



Each normalized section $E_j$ is smooth on $u$.
\end{proposition}

\begin{proof}
\leanok
\uses{prop:pseudo-gram-schmidt-frame-aux-smooth,lem:inv-sqrt-abs-smooth,lem:pseudo-metric-pairing-smooth,lem:pseudo-gram-schmidt-denominators}
\sketch The function $q_j(x)=g_x(U_j|_x,U_j|_x)$ is smooth and nowhere zero.
Lemma \ref{lem:inv-sqrt-abs-smooth} makes $|q_j|^{-1/2}$ smooth, and
$E_j=|q_j|^{-1/2}U_j$.
\end{proof}

\subsection{Orthonormal Frames}

\begin{lemma}[Nondegeneracy of a Tuple via its Gram Matrix]
\label{lem:nondegenerate-tuple-gram-matrix}
\lean{LeeLib.Ch02.isNondegenerate_span_iff_det_gramMatrix_ne_zero}
\uses{lem:nondegenerate-symmetric-bilinear}
\leanok



Let $(V,q)$ be a scalar product space and $(v_1,\ldots,v_k)$ a tuple. Its span
is nondegenerate of dimension $k$ if and only if the Gram matrix
\[
G=(\langle v_i,v_j\rangle)_{i,j=1}^k
\]
is invertible.
\end{lemma}

\begin{proof}
\leanok
\uses{lem:nondegenerate-symmetric-bilinear}
\sketch If $G$ is invertible, pairing a linear relation or a vector in the
restricted radical with all $v_j$ gives $Ga=0$, hence all coefficients vanish.
Conversely, when the span is nondegenerate of dimension $k$, the tuple is its
basis and $G$ is the matrix of the isomorphism induced by $q$ from the span to
its dual.
\end{proof}

\begin{lemma}[Nondegeneracy of a Tuple via its Leading Principal Minors]
\label{lem:nondegenerate-tuple-leading-minors}
\lean{LeeLib.Ch02.isNondegenerateTuple_iff_forall_det_gramMatrix_ne_zero}
\uses{lem:nondegenerate-tuple-gram-matrix}
\leanok



Let $G_j$ be the Gram matrix of the prefix $(v_1,\ldots,v_j)$. Then
$(v_1,\ldots,v_n)$ is nondegenerate if and only if
\[
\det G_j\ne0\qquad(j\leq n).
\]
\end{lemma}

\begin{proof}
\leanok
\uses{lem:nondegenerate-tuple-gram-matrix}
\sketch Apply Lemma \ref{lem:nondegenerate-tuple-gram-matrix} to every prefix.
\end{proof}

\begin{lemma}[A Nondegenerate Tuple is Linearly Independent]
\label{lem:nondegenerate-tuple-independent}
\lean{LeeLib.Ch02.IsNondegenerateTuple.linearIndependent}
\leanok


A nondegenerate $n$-tuple in a scalar product space is linearly independent.
\end{lemma}

\begin{proof}
\leanok
\sketch Its final prefix has dimension $n$ and is spanned by the $n$ vectors.
\end{proof}

\begin{lemma}[A Nondegenerate Tuple of Full Length is a Basis]
\label{lem:nondegenerate-tuple-basis}
\lean{LeeLib.Ch02.IsNondegenerateTuple.basis}
\uses{lem:nondegenerate-tuple-independent}
\leanok



In an $n$-dimensional scalar product space, every nondegenerate $n$-tuple is a
basis.
\end{lemma}

\begin{proof}
\leanok
\uses{lem:nondegenerate-tuple-independent}
\sketch It is an independent family of $n=\dim V$ vectors.
\end{proof}

\begin{lemma}[An Orthonormal Tuple is Nondegenerate]
\label{lem:orthonormal-tuple-nondegenerate}
\lean{LeeLib.Ch02.IsOrthonormal.isNondegenerateTuple}
\uses{lem:nondegenerate-tuple-leading-minors}
\leanok



Every orthonormal tuple in a scalar product space is nondegenerate.
\end{lemma}

\begin{proof}
\leanok
\uses{lem:nondegenerate-tuple-leading-minors}
\sketch Each leading Gram matrix is diagonal with entries $\pm1$, so every
leading determinant is nonzero. Apply Lemma
\ref{lem:nondegenerate-tuple-leading-minors}.
\end{proof}

\begin{lemma}[Smoothness of the Leading Gram Determinants]
\label{lem:gram-det-smooth}
\lean{LeeLib.Ch02.contMDiffOn_det_gramMatrix}
\uses{lem:pseudo-metric-pairing-smooth}
\leanok



Let $Y_1,\ldots,Y_k$ be sections smooth on $u$ for a smooth field $g$ of scalar
products. Then
\[
x\longmapsto\det(g_x(Y_i|_x,Y_l|_x))_{i,l=1}^k
\]
is smooth on $u$.
\end{lemma}

\begin{proof}
\leanok
\uses{lem:pseudo-metric-pairing-smooth}
\sketch Every matrix entry is smooth, and the determinant is a polynomial in
those entries.
\end{proof}

\begin{lemma}[Nondegeneracy of a Tuple of Smooth Sections is an Open Condition]
\label{lem:open-nondegenerate-tuple}
\lean{LeeLib.Ch02.exists_isLocalNondegenerateOn_nhds}
\uses{def:local-nondegenerate-family, lem:nondegenerate-tuple-leading-minors, lem:gram-det-smooth, lem:pseudo-metric-fibre-scalar-product}
\leanok



Let $u\subseteq B$ be open and let $X_1,\ldots,X_m$ be smooth on $u$. If their
values at $x_0\in u$ form a nondegenerate tuple, then some open
$w$ with $x_0\in w\subseteq u$ makes $(X_1,\ldots,X_m)$ nondegenerate on $w$.
\end{lemma}

\begin{proof}
\leanok
\uses{lem:nondegenerate-tuple-leading-minors,lem:gram-det-smooth,lem:pseudo-metric-fibre-scalar-product}
\sketch For each $k\leq m$, let $G_k(x)$ be the leading $k\times k$ Gram matrix.
Its determinant is continuous and nonzero at $x_0$. Intersect the finitely many
open neighborhoods on which all $\det G_k$ remain nonzero, then use Lemma
\ref{lem:nondegenerate-tuple-leading-minors} fibrewise.
\end{proof}

\begin{lemma}[Smooth Sections Through a Prescribed Basis]
\label{lem:sections-through-prescribed-basis}
\lean{LeeLib.AppendixA.exists_contMDiffOn_section_eq_basis}
\leanok


Let $V\to B$ have rank $m$, let $x_0\in B$, and let
$(w_1,\ldots,w_m)$ be a basis of $V_{x_0}$. There are an open neighborhood
$v$ of $x_0$ and sections $Y_j$ smooth on $v$ with $Y_j|_{x_0}=w_j$.
\end{lemma}

\begin{proof}
\leanok
\sketch In a local trivialization $\Phi$, set
\[
Y_j|_x=\Phi_x^{-1}(\Phi_{x_0}(w_j)).
\]
These sections are constant in trivialized coordinates.
\end{proof}

\begin{proposition}[Orthonormalization of a Nondegenerate Family of Vector Fields]
\label{prop:orthonormalize-nondegenerate-family}
\lean{LeeLib.Ch02.exists_pseudo_orthonormalFrame}
\uses{def:pseudo-riemannian-metric, def:local-nondegenerate-family, def:pseudo-gram-schmidt-frame, prop:pseudo-gram-schmidt-frame-smooth, prop:pseudo-gram-schmidt-frame-orthonormal, lem:pseudo-gram-schmidt-frame-flag}
\leanok



Let $(X_1,\ldots,X_m)$ be nondegenerate on $u$ in a pseudo-Riemannian manifold.
There are vector fields $(Y_1,\ldots,Y_m)$ smooth on $u$, pointwise orthonormal,
and satisfying
\[
\operatorname{span}(Y_1|_x,\ldots,Y_k|_x)
=\operatorname{span}(X_1|_x,\ldots,X_k|_x)
\quad(k\leq m).
\]
\end{proposition}

\begin{proof}
\leanok
\uses{def:pseudo-gram-schmidt-frame,prop:pseudo-gram-schmidt-frame-smooth,prop:pseudo-gram-schmidt-frame-orthonormal,lem:pseudo-gram-schmidt-frame-flag}
\sketch Take the fibrewise Gram--Schmidt frame. The cited proposition and lemma
give smoothness, orthonormality, and equality of prefix spans.
\end{proof}

\begin{proposition}[Orthonormal Frames for Pseudo-Riemannian Manifolds]
\label{prop:orthonormal-frames-pseudo-riemannian}
\lean{LeeLib.Ch02.exists_pseudo_orthonormalFrame_nhds}
\uses{def:pseudo-riemannian-metric, lem:completion-nondegenerate-bases, lem:sections-through-prescribed-basis, lem:open-nondegenerate-tuple, prop:orthonormalize-nondegenerate-family}
\leanok


\dcref{ch2:2.66}

Every point of a pseudo-Riemannian manifold has a neighborhood carrying a smooth
orthonormal frame.
\end{proposition}

\begin{proof}
\leanok
\uses{lem:completion-nondegenerate-bases,lem:sections-through-prescribed-basis,lem:open-nondegenerate-tuple,prop:orthonormalize-nondegenerate-family,lem:pseudo-metric-fibre-scalar-product}
\sketch Complete the empty tuple to a nondegenerate basis of $T_pM$. Extend this
basis to smooth local vector fields and shrink the neighborhood so their leading
Gram determinants remain nonzero. Proposition
\ref{prop:orthonormalize-nondegenerate-family} then produces $n$ smooth
orthonormal vector fields, which form a frame in each $n$-dimensional tangent
space.
\end{proof}

\begin{proof}
  \uses{prop:orthonormal-frames-pseudo-riemannian, lem:completion-nondegenerate-bases, lem:sections-through-prescribed-basis, lem:open-nondegenerate-tuple, prop:orthonormalize-nondegenerate-family, lem:pseudo-metric-fibre-scalar-product}
\leanok

\sketch This is the preceding proposition.
\end{proof}

For $r,s\geq0$, the \emph{pseudo-Euclidean space} $\mathbb R^{r,s}$ is
$\mathbb R^{r+s}$ with coordinates
$(\xi^1,\ldots,\xi^r,\tau^1,\ldots,\tau^s)$ and metric
\begin{equation}
\bar g^{(r,s)}=(d\xi^1)^2+\cdots+(d\xi^r)^2
 -(d\tau^1)^2-\cdots-(d\tau^s)^2. \tag{2.24}
\end{equation}

\begin{definition}[Lorentz Metric]
\label{def:lorentz-metric}
\lean{LeeLib.Ch02.IsLorentzMetric}
\uses{def:pseudo-metric-has-signature}
\leanok



A \emph{Lorentz metric} is a pseudo-Riemannian metric of index $1$, hence of
signature $(r,1)$ for some $r\geq0$. A pseudo-Riemannian manifold carrying such a
metric is a \emph{Lorentz manifold}.
\end{definition}

On an $n$-manifold a Lorentz metric has signature $(n-1,1)$. The metric
$\bar g^{(r,1)}$ on $\mathbb R^{r,1}$ is the \emph{Minkowski metric}; writing its
coordinates as $(\xi^1,\ldots,\xi^r,\tau)$,
\begin{equation}
\bar g^{(r,1)}=(d\xi^1)^2+\cdots+(d\xi^r)^2-(d\tau)^2. \tag{2.25}
\end{equation}

\begin{lemma}[Flipping the sign of a scalar product along a unit vector]
\label{lem:lorentz-form-pointwise}
\lean{LeeLib.Ch02.lorentzForm_isSymm_nondegenerate_signature}
\uses{prop:sylvesters-law-inertia, cor:scalar-product-signature-sum}
\leanok



Let $(V,\langle\cdot,\cdot\rangle)$ be an inner product space of dimension $n+1$
and let $u$ be a unit vector. The form
\[
B(v,w)=\langle v,w\rangle
 -2\langle v,u\rangle\langle w,u\rangle
\]
is symmetric, nondegenerate, and has signature $(n,1)$. Replacing $u$ by $-u$
does not change $B$.
\end{lemma}

\begin{proof}
\leanok
\uses{prop:sylvesters-law-inertia,cor:scalar-product-signature-sum}
\sketch Define $Rw=w-2\langle w,u\rangle u$. Then $R^2=\operatorname{id}$ and
$B(v,w)=\langle v,Rw\rangle$, proving nondegeneracy. Extend $u$ to an orthonormal
basis with first vector $u$; the matrix of $B$ is
$\operatorname{diag}(-1,1,\ldots,1)$, so its signature is $(n,1)$. The formula is
quadratic in $u$, hence sign-independent.
\end{proof}

\begin{definition}[Rank-1 tangent distribution]
\label{def:rank-one-distribution}
\lean{LeeLib.Ch02.IsRankOneDistribution}
\leanok


A \emph{rank-1 tangent distribution} on $M$ is a family of lines
$D_x\subseteq T_xM$ such that every point has a neighborhood $U$ with a smooth
nowhere-zero vector field $X$ satisfying $D_x=\operatorname{span}(X_x)$ for
$x\in U$.
\end{definition}

\begin{proposition}[A rank-1 distribution is a rank-1 subbundle]
\label{prop:rank-one-distribution-is-subbundle}
\lean{LeeLib.Ch02.hasLocalSubframes_of_isRankOneDistribution}
\uses{def:rank-one-distribution, lem:local-frame-criterion-for-subbundles, prop:existence-riemannian-metric}
\leanok



Every rank-1 tangent distribution is a smooth rank-1 subbundle of $TM$.
\end{proposition}

\begin{proof}
\leanok
\uses{def:rank-one-distribution,lem:local-frame-criterion-for-subbundles}
\sketch The local field in Definition \ref{def:rank-one-distribution} is a
one-element linearly independent family spanning $D$. Apply the local-frame
criterion with $k=1$.
\end{proof}

\begin{definition}[Tensor product of two $1$-forms]
\label{def:tensor-product-one-forms}
\lean{Bundle.formProduct}
\leanok


For covectors $\alpha,\beta\in V^*$, define
\[
(\alpha\otimes\beta)(v,w)=\alpha(v)\beta(w).
\]
For $1$-forms this definition is applied fibrewise.
\end{definition}

\begin{lemma}[The tensor product of two smooth $1$-forms is smooth]
\label{lem:tensor-product-one-forms-smooth}
\lean{Bundle.contMDiffAt_formProduct}
\uses{def:tensor-product-one-forms}
\leanok



If $\alpha,\beta$ are smooth sections of the dual of a smooth vector bundle
$E\to M$, then $x\mapsto\alpha_x\otimes\beta_x$ is a smooth field of bilinear
forms on $E$.
\end{lemma}

\begin{proof}
\leanok
\uses{def:tensor-product-one-forms}
\sketch In a local trivialization, $\alpha$ and $\beta$ are smooth maps into the
fixed model dual space. Compose them with the smooth bilinear tensor-product map.
\end{proof}

\begin{definition}[The metric with its sign flipped along a vector field]
\label{def:sign-flipped-form}
\lean{Bundle.flipFormOf}
\uses{def:tensor-product-one-forms}
\leanok



For a field $b$ of bilinear forms and a section $u$, define
\[
\flip(b,u)=b-2b(u,\cdot)\otimes b(u,\cdot),
\]
that is,
\[
\flip(b,u)_x(v,w)=b_x(v,w)-2b_x(u_x,v)b_x(u_x,w).
\]
\end{definition}

\begin{lemma}[The sign flip of a smooth form along a smooth section is smooth]
\label{lem:sign-flipped-form-smooth}
\lean{Bundle.contMDiffAt_flipFormOf}
\uses{def:sign-flipped-form, lem:tensor-product-one-forms-smooth}
\leanok



If $b$ and $u$ are smooth, then $\flip(b,u)$ is a smooth field of bilinear forms.
\end{lemma}

\begin{proof}
\leanok
\uses{lem:tensor-product-one-forms-smooth}
\sketch The $1$-form $b(u,\cdot)$ is smooth, so its tensor square is smooth by
Lemma \ref{lem:tensor-product-one-forms-smooth}; take the indicated linear
combination with $b$.
\end{proof}

\begin{lemma}[A line contains exactly two unit vectors]
\label{lem:two-unit-vectors-of-a-line}
\lean{LeeLib.Ch02.eq_or_eq_neg_of_span_eq}
\leanok


If unit vectors $u,u'$ in an inner product space span the same line, then
$u'=u$ or $u'=-u$.
\end{lemma}

\begin{proof}
\leanok
\sketch Write $u'=cu$. Unit length gives $c^2=1$.
\end{proof}

\begin{lemma}[Smooth local unit sections of a rank-1 distribution]
\label{lem:local-unit-section}
\lean{LeeLib.Ch02.IsRankOneDistribution.exists_localUnitSection}
\uses{def:rank-one-distribution, def:riemannian-metric}
\leanok



Let $(M,g)$ be Riemannian and $D$ a rank-1 tangent distribution. Every point has
a neighborhood $U$ carrying a smooth field $E$ such that
\[
g_x(E_x,E_x)=1,
\qquad\operatorname{span}(E_x)=D_x
\quad(x\in U).
\]
\end{lemma}

\begin{proof}
\leanok
\uses{def:rank-one-distribution}
\sketch Normalize a local nowhere-zero spanning field $X$ by
$E=X/\sqrt{g(X,X)}$. Positive definiteness makes the denominator smooth and
strictly positive, and normalization does not change the line.
\end{proof}

\begin{proposition}[Sufficiency in Theorem \ref{thm:lorentz-metrics}]
\label{prop:lorentz-metric-of-rank-one-distribution}
\lean{LeeLib.Ch02.exists_isLorentzMetric_of_isRankOneDistribution}
\uses{def:lorentz-metric, def:rank-one-distribution, prop:existence-riemannian-metric, lem:lorentz-form-pointwise}
\leanok



If $M$ admits a rank-1 tangent distribution, then it admits a Lorentz metric.
\end{proposition}

\begin{proof}
\leanok
\uses{lem:lorentz-form-pointwise,prop:existence-riemannian-metric,def:sign-flipped-form,lem:sign-flipped-form-smooth,lem:two-unit-vectors-of-a-line,lem:local-unit-section}
\sketch Choose a Riemannian metric $g$. For each $x$, choose either unit vector
$u_x$ in the line $D_x$ and set
\[
\bar g_x(v,w)=g_x(v,w)-2g_x(v,u_x)g_x(w,u_x).
\]
The pointwise sign-flip lemma makes this independent of the choice $u_x$ and gives
signature $(n,1)$ when $\dim M=n+1$. Locally choose a smooth unit section $E$ of
$D$. Since $E_x=\pm u_x$, the field $\bar g$ equals the smooth field
$\flip(g,E)$ locally; hence $\bar g$ is a Lorentz metric.
\end{proof}

\begin{definition}[Simple eigenpair]
\label{def:simple-eigenpair}
\lean{LeeLib.Ch02.IsSimpleEigenpair}
\leanok


Let $V$ be a finite-dimensional real inner product space. A triple
$(S,v,\lambda)$ is a \emph{simple eigenpair} if $S$ is symmetric, $v\ne0$,
$Sv=\lambda v$, and
\[
\ker(S-\lambda\operatorname{id})=\operatorname{span}(v).
\]
\end{definition}

\begin{definition}[Linearisation of the eigen-equation]
\label{def:eigen-linearisation}
\lean{LeeLib.Ch02.eigenLin}
\leanok


For $S\colon V\to V$, $v,v_0\in V$, and $\lambda\in\mathbb R$, define
\[
L_{S,v,\lambda,v_0}(w,\mu)
=\bigl(Sw-\lambda w-\mu v,\langle v_0,w\rangle\bigr)
\]
as a linear map $V\times\mathbb R\to V\times\mathbb R$.
\end{definition}

\begin{lemma}[Simplicity makes the linearisation injective]
\label{lem:eigen-linearisation-injective}
\lean{LeeLib.Ch02.eigenLin_injective}
\uses{def:simple-eigenpair, def:eigen-linearisation}
\leanok



If $(S,v,\lambda)$ is simple and $\langle v_0,v\rangle=1$, then
$L_{S,v,\lambda,v_0}$ is injective and therefore bijective.
\end{lemma}

\begin{proof}
\leanok
\uses{def:simple-eigenpair,def:eigen-linearisation}
\sketch If $L(w,\mu)=0$, pair $Sw-\lambda w=\mu v$ with $v$. Symmetry and
$Sv=\lambda v$ give $\mu\langle v,v\rangle=0$, hence $\mu=0$. Simplicity gives
$w=cv$, and $0=\langle v_0,w\rangle=c$. Thus the kernel is trivial; equal finite
dimensions give surjectivity.
\end{proof}

\begin{definition}[Augmented eigen-equation]
\label{def:augmented-eigen-equation}
\lean{LeeLib.Ch02.eigenAug}
\leanok


For $v_0\in V$, define
\[
\Phi_{v_0}(S,v,\lambda)
=\bigl(S,Sv-\lambda v,\langle v_0,v\rangle-1\bigr)
\]
on $\operatorname{End}(V)\times V\times\mathbb R$. Thus
$\Phi_{v_0}(S,v,\lambda)=(S,0,0)$ is the eigen-equation with normalization
$\langle v_0,v\rangle=1$.
\end{definition}

\begin{lemma}[The augmented equation has an invertible derivative at a simple eigenpair]
\label{lem:augmented-eigen-equation-derivative}
\lean{LeeLib.Ch02.hasFDerivAt_eigenAug}
\uses{def:augmented-eigen-equation, def:eigen-linearisation, def:simple-eigenpair, lem:eigen-linearisation-injective}
\leanok



The map $\Phi_{v_0}$ is smooth, with derivative
\[
D\Phi_{v_0}(S,v,\lambda)(T,w,\mu)
=\bigl(T,Tv+Sw-\lambda w-\mu v,\langle v_0,w\rangle\bigr).
\]
At a simple eigenpair satisfying $\langle v_0,v\rangle=1$, this derivative is a
linear isomorphism.
\end{lemma}

\begin{proof}
\leanok
\uses{lem:eigen-linearisation-injective}
\sketch The formula follows by differentiating the polynomial components. Relative
to $\operatorname{End}(V)\oplus(V\times\mathbb R)$, the derivative is block lower
triangular with diagonal blocks $\operatorname{id}$ and
$L_{S,v,\lambda,v_0}$; both are invertible.
\end{proof}

\begin{proposition}[Smooth selection of a simple eigenpair]
\label{prop:smooth-eigen-selection}
\lean{LeeLib.Ch02.exists_eigenSelection}
\uses{def:simple-eigenpair, def:augmented-eigen-equation, lem:augmented-eigen-equation-derivative}
\leanok



Let $(S_0,v_0,\lambda_0)$ be simple with $\langle v_0,v_0\rangle=1$. There are
an open $W\subseteq\operatorname{End}(V)$ containing $S_0$ and maps
$\mathcal V\colon W\to V$, $\Lambda\colon W\to\mathbb R$ such that
\begin{enumerate}
\item $\mathcal V(S_0)=v_0$ and $\Lambda(S_0)=\lambda_0$;
\item $S\mathcal V(S)=\Lambda(S)\mathcal V(S)$ and
$\langle v_0,\mathcal V(S)\rangle=1$ for all $S\in W$;
\item $\mathcal V$ and $\Lambda$ are continuous on $W$;
\item they are smooth at every $S$ for which
$(S,\mathcal V(S),\Lambda(S))$ is simple.
\end{enumerate}
\end{proposition}

\begin{proof}
\leanok
\uses{lem:augmented-eigen-equation-derivative,def:augmented-eigen-equation}
\sketch The inverse function theorem makes $\Phi_{v_0}$ a diffeomorphism from a
neighborhood of $(S_0,v_0,\lambda_0)$ onto a neighborhood of $(S_0,0,0)$. Let
\[
W=\{S:(S,0,0)\text{ lies in the latter neighborhood}\}
\]
and define $(\mathcal V(S),\Lambda(S))$ as the last two components of
$\Phi_{v_0}^{-1}(S,0,0)$. The inverse identity gives (1)--(3). At any selected
simple eigenpair, the derivative is again invertible, so the local inverse theorem
gives (4).
\end{proof}

\begin{lemma}[The index counts the negative eigenvalues]
\label{lem:signature-and-eigenvalue-signs}
\lean{LeeLib.Ch02.sigNeg_eq_ncard_neg_eigenvalues}
\uses{cor:scalar-product-signature-form, prop:sylvesters-law-inertia}
\leanok



Let $B$ be a symmetric nondegenerate form of signature $(r,s)$ on an inner product
space $V$, and let $A$ be determined by
$B(v,w)=\langle Av,w\rangle$. Then $A$ is symmetric and $s$ equals the number of
negative eigenvalues of $A$, counted with multiplicity.
\end{lemma}

\begin{proof}
\leanok
\uses{cor:scalar-product-signature-form,prop:sylvesters-law-inertia}
\sketch Symmetry of $B$ makes $A$ self-adjoint. In an orthonormal eigenbasis
$Au_i=\mu_i u_i$, the matrix of $B$ is $\operatorname{diag}(\mu_i)$. No $\mu_i$
vanishes by nondegeneracy. Rescaling $u_i$ by $|\mu_i|^{-1/2}$ gives an
indefinite orthonormal basis whose negative entries are exactly the negative
$\mu_i$. Sylvester's law identifies their number with $s$.
\end{proof}

\begin{lemma}[Index one yields a simple eigenpair]
\label{lem:simple-eigenpair-of-index-one}
\lean{LeeLib.Ch02.exists_isSimpleEigenpair_of_sigNeg_eq_one}
\uses{lem:signature-and-eigenvalue-signs, def:simple-eigenpair}
\leanok



In the preceding setting, if $B$ has index $1$, then $A$ has one negative
eigenvalue $\lambda$, counted with multiplicity; its eigenspace is a line, and
any nonzero vector $u$ in it makes $(A,u,\lambda)$ a simple eigenpair.
\end{lemma}

\begin{proof}
\leanok
\uses{lem:signature-and-eigenvalue-signs,def:simple-eigenpair}
\sketch The preceding lemma gives exactly one negative eigenvalue in an eigenbasis.
Hence its eigenspace has dimension one and equals the span of any nonzero
eigenvector in it.
\end{proof}

\begin{proposition}[Necessity in Theorem \ref{thm:lorentz-metrics}]
\label{prop:rank-one-distribution-of-lorentz-metric}
\lean{LeeLib.Ch02.exists_isRankOneDistribution_of_isLorentzMetric}
\uses{def:lorentz-metric, def:rank-one-distribution, prop:existence-riemannian-metric, lem:simple-eigenpair-of-index-one, prop:smooth-eigen-selection, prop:existence-orthonormal-frames}
\leanok



If $M$ admits a Lorentz metric, then it admits a rank-1 tangent distribution.
\end{proposition}

\begin{proof}
\leanok
\uses{prop:existence-riemannian-metric,lem:simple-eigenpair-of-index-one,prop:smooth-eigen-selection,prop:existence-orthonormal-frames,def:rank-one-distribution}
\sketch Let $\bar g$ be Lorentz, choose a Riemannian metric $g$, and define the
$g$-self-adjoint $A_x$ by
\[
\bar g_x(v,w)=g_x(A_xv,w).
\]
Let $D_x$ be the sum of the negative eigenspaces of $A_x$. Index one and Lemma
\ref{lem:simple-eigenpair-of-index-one} show that $D_x$ is the one-dimensional
eigenspace of the unique negative eigenvalue.

Near $x_0$, choose a smooth $g$-orthonormal frame $E_i$ and identify $T_xM$ with
$\mathbb R^n$ by $\iota_x(c)=\sum_i c_iE_i(x)$. The symmetric matrix
\[
S(x)_{ij}=\bar g_x(E_i,E_j)
\]
depends smoothly on $x$ and satisfies
$A_x\circ\iota_x=\iota_x\circ S(x)$. Choose a unit vector $v_0$ in the negative
eigenline of $S(x_0)$. Proposition \ref{prop:smooth-eigen-selection} provides
continuous local selections $\mathcal V,\Lambda$. After restricting to
\[
U'=\{x:S(x)\in W, \Lambda(S(x))<0\},
\]
the selected eigenpair is simple at every point, hence smooth. Therefore
$X_x=\iota_x\mathcal V(S(x))$ is a smooth nowhere-zero field spanning $D_x$ on
$U'$. Such a field exists near every $x_0$, so $D$ is a rank-1 distribution.
\end{proof}

\begin{theorem}
\label{thm:lorentz-metrics}
\lean{LeeLib.Ch02.exists_isLorentzMetric_iff_exists_isRankOneDistribution}
\uses{def:lorentz-metric, prop:existence-riemannian-metric, prop:lorentz-metric-of-rank-one-distribution, prop:rank-one-distribution-of-lorentz-metric}
\leanok
\dcref{ch2:2.69}



A smooth manifold admits a Lorentz metric if and only if it admits a rank-1
tangent distribution, equivalently a rank-1 subbundle of $TM$.
\end{theorem}

\begin{proof}
\leanok
\uses{prop:lorentz-metric-of-rank-one-distribution,prop:rank-one-distribution-of-lorentz-metric}
\sketch Apply Proposition \ref{prop:lorentz-metric-of-rank-one-distribution} in one
direction and Proposition \ref{prop:rank-one-distribution-of-lorentz-metric} in
the other.
\end{proof}

Every noncompact connected smooth manifold admits a Lorentz metric, while a compact connected smooth
manifold admits one if and only if its Euler characteristic is zero (see \cite{ON83}, p.~149). Thus no
even-dimensional sphere admits a Lorentz metric. For further treatments of pseudo-Riemannian metrics,
see \cite{ON83} and \cite{HE73}.
\section{Pseudo-Riemannian Submanifolds}

For a pseudo-Riemannian $(\widetilde M,\widetilde g)$, an immersion $F:M\to\widetilde M$ presents a pseudo-Riemannian submanifold when $F^*\widetilde g$ is nondegenerate with constant signature; the induced tensor is then the metric on $M$. It is a Riemannian submanifold when the pullback is positive definite. Immersion alone does not ensure nondegeneracy: the diagonal $\{(\xi,\tau):\xi=\tau\}\subset\mathbb R^{1,1}$ has identically zero induced tensor. For $p\in M$, the normal space consists of ambient vectors orthogonal to $dF_p(T_pM)$.

\begin{lemma}[Orthonormal Bases Adapted to a Nondegenerate Flag]
\label{lem:adapted-orthonormal-basis}
\lean{LeeLib.Ch02.exists_isOrthonormal_prefixSpan_eq_of_isNondegenerateTuple}
\uses{lem:completion-nondegenerate-bases, prop:gram-schmidt-algorithm-scalar-products}
\leanok



If $(V,q)$ has dimension $n$ and $(v_1,\ldots,v_k)$ is a nondegenerate tuple with $k\le n$, then $V$ has an orthonormal basis $(b_1,\ldots,b_n)$ satisfying
\[
\operatorname{span}(b_1,\ldots,b_j)=\operatorname{span}(v_1,\ldots,v_j),\qquad j\le k.
\]
\end{lemma}
\begin{proof}
\leanok
\uses{lem:completion-nondegenerate-bases,prop:gram-schmidt-algorithm-scalar-products}
\sketch Extend the tuple to a nondegenerate basis by Lemma \ref{lem:completion-nondegenerate-bases}, then apply Proposition \ref{prop:gram-schmidt-algorithm-scalar-products}, whose orthonormalization preserves the complete span flag.
\end{proof}

\begin{lemma}[Orthonormal Basis Adapted to a Non-Null Line]
\label{lem:adapted-orthonormal-basis-line}
\lean{LeeLib.Ch02.exists_isOrthonormal_prefixSpan_one_eq_span}
\uses{lem:adapted-orthonormal-basis}
\leanok



If $\dim V=n\ge1$ and $\langle x,x\rangle\ne0$, then $V$ has an orthonormal basis $(b_i)$ with $\operatorname{span}(b_1)=\operatorname{span}(x)$.
\end{lemma}
\begin{proof}
\leanok
\uses{lem:adapted-orthonormal-basis}
\sketch The hypothesis makes $x$ nonzero and its line nondegenerate: $cx$ orthogonal to $x$ forces $c=0$. Apply Lemma \ref{lem:adapted-orthonormal-basis} with $k=1$.
\end{proof}

\begin{lemma}[Nondegeneracy of a Line]
\label{lem:nondegenerate-line}
\lean{LeeLib.Ch02.restrict_nondegenerate_iff_of_finrank_eq_one}
\leanok


For a one-dimensional subspace $L$ and bilinear $q$, the restriction $q|_{L\times L}$ is nondegenerate if and only if $\langle v,v\rangle\ne0$ for every nonzero $v\in L$.
\end{lemma}
\begin{proof}
\leanok
\sketch A nonzero $v\in L$ spans $L$. Thus $\langle v,v\rangle=0$ makes $v$ orthogonal to all of $L$, while conversely any vector orthogonal to all of $L$ has zero square and hence vanishes under the stated condition.
\end{proof}

In dimensions at least two, absence of nonzero null vectors is stronger than nondegeneracy; for example, $\xi^2-\tau^2$ on $\mathbb R^2$ is nondegenerate but has null lines.

\begin{lemma}[Nondegeneracy of a Hyperplane]
\label{lem:hyperplane-nondegenerate}
\lean{LeeLib.Ch02.restrict_nondegenerate_iff_forall_orthogonal}
\uses{lem:nondegenerate-subspace, lem:nondegenerate-line, lem:orthogonal-complement}
\leanok



If $(V,q)$ is a scalar product space and $\dim W+1=\dim V$, then $q|_W$ is nondegenerate if and only if every nonzero $v\in W^\perp$ satisfies $\langle v,v\rangle\ne0$.
\end{lemma}
\begin{proof}
\leanok
\uses{lem:nondegenerate-subspace,lem:nondegenerate-line,lem:orthogonal-complement}
\sketch Lemma \ref{lem:nondegenerate-subspace} equates nondegeneracy on $W$ and $W^\perp$. Lemma \ref{lem:orthogonal-complement} makes the latter one-dimensional, so Lemma \ref{lem:nondegenerate-line} applies.
\end{proof}

\begin{lemma}[The Tail of an Orthonormal Basis Spans the Complement of its Head]
\label{lem:orthonormal-tail-spans-orthogonal}
\lean{LeeLib.Ch02.span_range_succ_eq_orthogonal}
\uses{lem:orthogonal-complement}
\leanok



For an orthonormal basis $(b_1,\ldots,b_n)$ of a finite-dimensional nondegenerate space,
\[
\operatorname{span}(b_2,\ldots,b_n)=\operatorname{span}(b_1)^\perp.
\]
\end{lemma}
\begin{proof}
\leanok
\uses{lem:orthogonal-complement}
\sketch Orthogonality gives inclusion. Both sides have dimension $n-1$, the second by Lemma \ref{lem:orthogonal-complement}, so equality follows.
\end{proof}

\begin{lemma}[Orthonormal Basis of the Complement of a Line]
\label{lem:orthonormal-basis-orthogonal-complement}
\lean{LeeLib.Ch02.exists_basis_isOrthonormal_orthogonal_span_singleton}
\uses{lem:orthonormal-tail-spans-orthogonal}
\leanok



Under the preceding hypotheses, $(b_2,\ldots,b_n)$ is an orthonormal basis of $\operatorname{span}(b_1)^\perp$ with the restricted form.
\end{lemma}
\begin{proof}
\leanok
\uses{lem:orthonormal-tail-spans-orthogonal}
\sketch Lemma \ref{lem:orthonormal-tail-spans-orthogonal} gives membership and spanning; linear independence and orthonormality are inherited from the ambient basis.
\end{proof}

\begin{lemma}[The Complement of a Non-Null Line is a Hyperplane]
\label{lem:orthogonal-complement-non-null-line-hyperplane}
\lean{LeeLib.Ch02.finrank_orthogonal_span_singleton}
\uses{lem:orthogonal-complement}
\leanok



If $(V,q)$ has dimension $n$ and $x\ne0$, then $\dim\operatorname{span}(x)^\perp=n-1$.
\end{lemma}
\begin{proof}
\leanok
\uses{lem:orthogonal-complement}
\sketch Since $\operatorname{span}(x)$ has dimension one, apply Lemma \ref{lem:orthogonal-complement}.
\end{proof}

\begin{lemma}[The Complement of a Non-Null Line is a Scalar Product Space]
\label{lem:orthogonal-complement-non-null-line-nondegenerate}
\lean{LeeLib.Ch02.restrict_orthogonal_span_singleton_nondegenerate}
\uses{lem:nondegenerate-subspace, lem:adapted-orthonormal-basis-line}
\leanok



If $(V,q)$ is a scalar product space and $\langle x,x\rangle\ne0$, then $q|_{\operatorname{span}(x)^\perp}$ is nondegenerate.
\end{lemma}
\begin{proof}
\leanok
\uses{lem:nondegenerate-subspace,lem:adapted-orthonormal-basis-line}
\sketch The line of $x$ is nondegenerate, as in Lemma \ref{lem:adapted-orthonormal-basis-line}; Lemma \ref{lem:nondegenerate-subspace} then gives nondegeneracy of its complement.
\end{proof}

\begin{proposition}[Signature of the Complement of a Non-Null Line]
\label{prop:hypersurface-signature-pointwise}
\lean{LeeLib.Ch02.sigPos_sigNeg_restrict_orthogonal_span_singleton}
\uses{lem:adapted-orthonormal-basis-line, lem:orthonormal-basis-orthogonal-complement, lem:orthogonal-complement-non-null-line-nondegenerate, prop:sylvesters-law-inertia}
\leanok



Let $(V,q)$ have signature $(r,s)$, let $\langle x,x\rangle\ne0$, and let $(r',s')$ be the signature of $\operatorname{span}(x)^\perp$, which is nondegenerate by Lemma \ref{lem:orthogonal-complement-non-null-line-nondegenerate}. If $\langle x,x\rangle>0$, then $r'+1=r$ and $s'=s$; if $\langle x,x\rangle<0$, then $r'=r$ and $s'+1=s$.
\end{proposition}
\begin{proof}
\leanok
\uses{lem:adapted-orthonormal-basis-line,lem:orthonormal-basis-orthogonal-complement,prop:sylvesters-law-inertia}
\sketch Choose an orthonormal basis with $b_1$ spanning $x$ by Lemma \ref{lem:adapted-orthonormal-basis-line}; its tail is an orthonormal basis of the complement by Lemma \ref{lem:orthonormal-basis-orthogonal-complement}. Sylvester's law counts signs, and $b_1=cx$ has the sign of $x$.
\end{proof}

\begin{lemma}[A Hyperplane is the Complement of its Normal Line]
\label{lem:hyperplane-eq-complement-normal-line}
\lean{LeeLib.Ch02.orthogonal_span_singleton_eq}
\uses{lem:orthogonal-complement, lem:orthogonal-complement-involution}
\leanok



If $W\subseteq V$ is a hyperplane and $0\ne x\in W^\perp$, then $\operatorname{span}(x)^\perp=W$.
\end{lemma}
\begin{proof}
\leanok
\uses{lem:orthogonal-complement,lem:orthogonal-complement-involution}
\sketch Lemma \ref{lem:orthogonal-complement} makes $W^\perp$ one-dimensional, hence equal to $\operatorname{span}(x)$; apply Lemma \ref{lem:orthogonal-complement-involution}.
\end{proof}

\begin{proposition}[Signature of a Hypersurface's Tangent Space]
\label{prop:hypersurface-signature-pointwise-hyperplane}
\lean{LeeLib.Ch02.sigPos_sigNeg_restrict_of_mem_orthogonal}
\uses{prop:hypersurface-signature-pointwise, lem:hyperplane-eq-complement-normal-line}
\leanok



Let $(V,q)$ have signature $(r,s)$, let $W$ be a hyperplane, and let $0\ne x\in W^\perp$. If $(r',s')$ is the signature of $q|_W$, then a positive $x$ gives $r'+1=r$, $s'=s$, while a negative $x$ gives $r'=r$, $s'+1=s$.
\end{proposition}
\begin{proof}
\leanok
\uses{prop:hypersurface-signature-pointwise,lem:hyperplane-eq-complement-normal-line}
\sketch Lemma \ref{lem:hyperplane-eq-complement-normal-line} identifies $W$ with $\operatorname{span}(x)^\perp$; apply Proposition \ref{prop:hypersurface-signature-pointwise}.
\end{proof}

\begin{lemma}[The Normal Space is an Orthogonal Complement]
\label{lem:pseudo-normal-space-eq-orthogonal}
\lean{LeeLib.Ch02.pseudoNormalSpace_eq_orthogonal}
\uses{def:pseudo-normal-space, def:pseudo-riemannian-metric}
\leanok



For smooth $F:M\to(\widetilde M,\widetilde g)$,
\[
N_xM=(dF_x(T_xM))^\perp
\]
inside $(T_{F(x)}\widetilde M,\widetilde g_{F(x)})$.
\end{lemma}
\begin{proof}
\leanok
\uses{def:pseudo-normal-space,def:pseudo-riemannian-metric}
\sketch Definition \ref{def:pseudo-normal-space} tests $\widetilde g(v,w)=0$, while the complement tests $\widetilde g(w,v)=0$; symmetry makes the conditions equivalent.
\end{proof}

\begin{lemma}[A Hypersurface's Tangent Space is a Hyperplane]
\label{lem:hypersurface-tangent-hyperplane}
\lean{LeeLib.Ch02.finrank_tangentRange_add_one}
\leanok


If $dF_x$ is injective and $\dim M+1=\dim\widetilde M$, then
\[
\dim dF_x(T_xM)+1=\dim T_{F(x)}\widetilde M.
\]
\end{lemma}
\begin{proof}
\leanok
\sketch Injectivity gives $\dim dF_x(T_xM)=\dim M$, and the ambient tangent space has dimension $\dim\widetilde M$.
\end{proof}

\begin{lemma}[The Normal Space of a Hypersurface is a Line]
\label{lem:pseudo-normal-space-line}
\lean{LeeLib.Ch02.finrank_pseudoNormalSpace}
\uses{def:pseudo-normal-space, lem:pseudo-normal-space-eq-orthogonal, lem:hypersurface-tangent-hyperplane, lem:orthogonal-complement}
\leanok



Under the preceding codimension-one and injectivity hypotheses, $\dim N_xM=1$.
\end{lemma}
\begin{proof}
\leanok
\uses{lem:pseudo-normal-space-eq-orthogonal,lem:hypersurface-tangent-hyperplane,lem:orthogonal-complement}
\sketch Combine $N_xM=dF_x(T_xM)^\perp$ from Lemma \ref{lem:pseudo-normal-space-eq-orthogonal}, the dimension identity of Lemma \ref{lem:orthogonal-complement}, and Lemma \ref{lem:hypersurface-tangent-hyperplane}.
\end{proof}

\begin{lemma}[The Sign of a Line is Decided by One of its Vectors]
\label{lem:line-sign-one-vector}
\lean{LeeLib.Ch02.pos_of_mem_of_finrank_eq_one}
\leanok


Let $N$ be a one-dimensional subspace, let $B$ be bilinear, and suppose $0\ne x\in N$ with $B(x,x)>0$. Then $B(v,v)>0$ for every nonzero $v\in N$.
\end{lemma}
\begin{proof}
\leanok
\sketch Every nonzero $v\in N$ is $ax$ with $a\ne0$, so $B(v,v)=a^2B(x,x)>0$.
\end{proof}

\begin{lemma}[One Positive Normal Makes Every Normal Positive]
\label{lem:one-normal-decides-sign}
\lean{LeeLib.Ch02.forall_pos_of_pos_mem_pseudoNormalSpace}
\uses{def:pseudo-normal-space, lem:pseudo-normal-space-line, lem:line-sign-one-vector}
\leanok



Under the hypotheses of Lemma \ref{lem:pseudo-normal-space-line}, if some nonzero $v_0\in N_xM$ has positive square, then every nonzero $v\in N_xM$ has positive square.
\end{lemma}
\begin{proof}
\leanok
\uses{lem:pseudo-normal-space-line,lem:line-sign-one-vector}
\sketch Lemma \ref{lem:pseudo-normal-space-line} makes $N_xM$ a line; apply Lemma \ref{lem:line-sign-one-vector} to $\widetilde g_{F(x)}$.
\end{proof}

\begin{lemma}[One Negative Normal Makes Every Normal Negative]
\label{lem:one-normal-decides-sign-neg}
\lean{LeeLib.Ch02.forall_neg_of_neg_mem_pseudoNormalSpace}
\uses{def:pseudo-normal-space, lem:pseudo-normal-space-line, lem:line-sign-one-vector}
\leanok



Under the same hypotheses, one negative nonzero normal implies that every nonzero normal is negative.
\end{lemma}
\begin{proof}
\leanok
\uses{lem:one-normal-decides-sign}
\sketch Apply Lemma \ref{lem:one-normal-decides-sign} to $-\widetilde g$, which has the same normal subspace.
\end{proof}

\begin{lemma}[A Hypersurface has a Nonzero Normal Vector at Every Point]
\label{lem:exists-nonzero-normal}
\lean{LeeLib.Ch02.exists_ne_zero_mem_pseudoNormalSpace}
\uses{def:pseudo-normal-space, lem:pseudo-normal-space-line}
\leanok



Under the hypotheses of Lemma \ref{lem:pseudo-normal-space-line}, $N_xM$ contains a nonzero vector.
\end{lemma}
\begin{proof}
\leanok
\uses{lem:pseudo-normal-space-line}
\sketch Lemma \ref{lem:pseudo-normal-space-line} gives $\dim N_xM=1>0$.
\end{proof}

\begin{lemma}[The Differential is an Isometry onto the Tangent Space]
\label{lem:differential-isometry-tangent-range}
\lean{LeeLib.Ch02.isometryEquivRestrictTangentRange}
\uses{def:induced-tensor-along-map, lem:pseudo-submanifold-immersion}
\leanok



If $g$ is induced by $\widetilde g$ along smooth $F:M\to\widetilde M$, then $dF_x$ is an isometry from $(T_xM,g_x)$ onto $dF_x(T_xM)$ with the restricted ambient form.
\end{lemma}
\begin{proof}
\leanok
\uses{def:induced-tensor-along-map,lem:pseudo-submanifold-immersion}
\sketch Lemma \ref{lem:pseudo-submanifold-immersion} makes $dF_x$ an isomorphism onto its image, and Definition \ref{def:induced-tensor-along-map} says it preserves the forms.
\end{proof}

\begin{lemma}[The Signature of a Submanifold is Computed in the Ambient Tangent Space]
\label{lem:submanifold-signature-in-ambient-fibre}
\lean{LeeLib.Ch02.sigPosAt_eq_sigPos_restrict_tangentRange, LeeLib.Ch02.sigNegAt_eq_sigNeg_restrict_tangentRange}
\uses{def:pseudo-metric-signature-at, lem:differential-isometry-tangent-range, prop:sylvesters-law-inertia}
\leanok



If $g$ is induced by $\widetilde g$ along $F$, then the signature of $g$ at $x$ equals that of $\widetilde g_{F(x)}$ restricted to $dF_x(T_xM)$.
\end{lemma}
\begin{proof}
\leanok
\uses{def:pseudo-metric-signature-at,lem:differential-isometry-tangent-range,prop:sylvesters-law-inertia}
\sketch By Definition \ref{def:pseudo-metric-signature-at} and Lemma \ref{lem:differential-isometry-tangent-range}, the spaces are isometric. An isometry transports orthonormal bases without changing sign counts, which determine signature by Proposition \ref{prop:sylvesters-law-inertia}.
\end{proof}

\begin{proposition}[Nondegeneracy of a Hypersurface]
\label{prop:pseudo-riemannian-hypersurface-nondegenerate}
\lean{LeeLib.Ch02.restrict_tangentRange_nondegenerate_iff}
\uses{def:pseudo-riemannian-metric, def:pseudo-normal-space, lem:pseudo-normal-space-eq-orthogonal, lem:hypersurface-tangent-hyperplane, lem:hyperplane-nondegenerate}
\leanok


\dcref{ch2:2.70}

Let $F:M\to\widetilde M$ be smooth with $\dim M+1=\dim\widetilde M$ and $dF_x$ injective. The restriction of $\widetilde g_{F(x)}$ to $dF_x(T_xM)$ is nondegenerate if and only if every nonzero $v\in N_xM$ is non-null.
\end{proposition}
\begin{proof}
\leanok
\uses{lem:pseudo-normal-space-eq-orthogonal,lem:hypersurface-tangent-hyperplane,lem:hyperplane-nondegenerate}
\sketch Apply Lemma \ref{lem:hyperplane-nondegenerate} to $W=dF_x(T_xM)$; Lemma \ref{lem:hypersurface-tangent-hyperplane} supplies codimension one and Lemma \ref{lem:pseudo-normal-space-eq-orthogonal} identifies $W^\perp=N_xM$.
\end{proof}

\begin{proposition}[A Single Normal Vector Decides the Signature at a Point]
\label{prop:hypersurface-signature-at-a-point}
\lean{LeeLib.Ch02.sigPosAt_sigNegAt_of_mem_pseudoNormalSpace}
\uses{def:pseudo-metric-has-signature, def:pseudo-normal-space, def:induced-tensor-along-map, lem:submanifold-signature-in-ambient-fibre, prop:hypersurface-signature-pointwise-hyperplane, lem:pseudo-normal-space-eq-orthogonal, lem:hypersurface-tangent-hyperplane, lem:pseudo-submanifold-immersion}
\leanok



Let $(\widetilde M,\widetilde g)$ have signature $(r,s)$ and let $F$ present a codimension-one pseudo-Riemannian submanifold with induced metric $g$. If $(r_x,s_x)$ is the signature at $x$ and $0\ne v\in N_xM$, then positive $v$ gives $r_x+1=r$, $s_x=s$, while negative $v$ gives $r_x=r$, $s_x+1=s$.
\end{proposition}
\begin{proof}
\leanok
\uses{def:pseudo-metric-has-signature,lem:submanifold-signature-in-ambient-fibre,prop:hypersurface-signature-pointwise-hyperplane,lem:pseudo-normal-space-eq-orthogonal,lem:hypersurface-tangent-hyperplane,lem:pseudo-submanifold-immersion}
\sketch In $V=T_{F(x)}\widetilde M$, put $W=dF_x(T_xM)$. Definition \ref{def:pseudo-metric-has-signature}, Lemmas \ref{lem:pseudo-submanifold-immersion}, \ref{lem:hypersurface-tangent-hyperplane}, and \ref{lem:pseudo-normal-space-eq-orthogonal} meet the hypotheses of Proposition \ref{prop:hypersurface-signature-pointwise-hyperplane}; transfer its result to $g_x$ by Lemma \ref{lem:submanifold-signature-in-ambient-fibre}.
\end{proof}

\begin{proposition}
\label{prop:pseudo-riemannian-hypersurface-signature}
\lean{LeeLib.Ch02.hasSignature_of_forall_normal_pos, LeeLib.Ch02.hasSignature_of_forall_normal_neg}
\uses{def:pseudo-riemannian-metric, def:pseudo-metric-has-signature, def:induced-tensor-along-map, def:pseudo-normal-space, lem:exists-nonzero-normal, prop:hypersurface-signature-at-a-point, lem:pseudo-submanifold-immersion}
\leanok


\dcref{ch2:2.70}

Let $(\widetilde M,\widetilde g)$ have signature $(r,s)$ and let $F$ present a codimension-one pseudo-Riemannian submanifold with induced $g$. If every nonzero normal is positive, then $g$ has signature $(r-1,s)$; if every nonzero normal is negative, then $g$ has signature $(r,s-1)$.
\end{proposition}
\begin{proof}
\leanok
\uses{def:pseudo-metric-has-signature,lem:exists-nonzero-normal,prop:hypersurface-signature-at-a-point,lem:pseudo-submanifold-immersion}
\sketch At each $p$, Lemma \ref{lem:pseudo-submanifold-immersion} and Lemma \ref{lem:exists-nonzero-normal} provide a nonzero normal. Proposition \ref{prop:hypersurface-signature-at-a-point} gives $(r-1,s)$ in the positive case, including $r\ge1$, and $(r,s-1)$ in the negative case, including $s\ge1$. Definition \ref{def:pseudo-metric-has-signature} makes the pointwise result global.
\end{proof}

\begin{lemma}[The Induced Tensor of a Non-Null Hypersurface is Nondegenerate]
\label{lem:pseudo-hypersurface-induced-tensor-nondegenerate}
\lean{LeeLib.Ch02.exists_pseudoPullbackForm_ne_zero}
\uses{def:pseudo-pullback-form, def:pseudo-normal-space, prop:pseudo-riemannian-hypersurface-nondegenerate}
\leanok



Let $F:M\to\widetilde M$ be a codimension-one immersion. If every nonzero normal is non-null, then $F^*\widetilde g$ is pointwise nondegenerate: for each $0\ne v\in T_pM$, some $w\in T_pM$ satisfies $(F^*\widetilde g)_p(v,w)\ne0$.
\end{lemma}
\begin{proof}
\leanok
\uses{def:pseudo-pullback-form,prop:pseudo-riemannian-hypersurface-nondegenerate}
\sketch Proposition \ref{prop:pseudo-riemannian-hypersurface-nondegenerate} makes the ambient restriction to $W=dF_p(T_pM)$ nondegenerate. For $u=dF_pv\ne0$, choose $u'=dF_pw\in W$ with $\widetilde g(u,u')\ne0$; Definition \ref{def:pseudo-pullback-form} gives the claim.
\end{proof}

\begin{corollary}[Signature of a Pseudo-Riemannian Hypersurface]
\label{cor:pseudo-riemannian-hypersurface-submanifold}
\lean{LeeLib.Ch02.hasSignature_pseudoPullbackMetric_of_forall_normal_pos, LeeLib.Ch02.hasSignature_pseudoPullbackMetric_of_forall_normal_neg}
\uses{def:pseudo-metric-has-signature, def:pseudo-normal-space, def:pseudo-induced-metric, lem:pseudo-hypersurface-induced-tensor-nondegenerate, lem:pseudo-riemannian-submanifold-induced-tensor, prop:pseudo-riemannian-hypersurface-signature}
\leanok


\dcref{ch2:2.70}

Let $(\widetilde M,\widetilde g)$ have signature $(r,s)$ and let $F:M\to\widetilde M$ be a codimension-one immersion. If all nonzero normals are positive, $F^*\widetilde g$ is pseudo-Riemannian of signature $(r-1,s)$; if all are negative, it has signature $(r,s-1)$.
\end{corollary}
\begin{proof}
\leanok
\uses{def:pseudo-induced-metric,lem:pseudo-hypersurface-induced-tensor-nondegenerate,lem:pseudo-riemannian-submanifold-induced-tensor,prop:pseudo-riemannian-hypersurface-signature}
\sketch Lemma \ref{lem:pseudo-hypersurface-induced-tensor-nondegenerate} and Definition \ref{def:pseudo-induced-metric} construct the metric. Lemma \ref{lem:pseudo-riemannian-submanifold-induced-tensor} identifies it as induced, and Proposition \ref{prop:pseudo-riemannian-hypersurface-signature} computes its signature.
\end{proof}

For $f$ on a pseudo-Riemannian manifold, $\operatorname{grad}f=(df)^\sharp$.

\begin{definition}[The gradient of a functional relative to a scalar product]
\label{def:bilinear-gradient}
\lean{LeeLib.Ch02.bilinGrad}
\uses{lem:nondegenerate-symmetric-bilinear}
\leanok



For a finite-dimensional $V$, nondegenerate symmetric $B$, and $a\in V^*$, the $B$-gradient $a^\sharp$ is the unique vector satisfying
\[
B(a^\sharp,w)=a(w),\qquad w\in V.
\]
\end{definition}
\begin{proof}
\leanok
\uses{lem:nondegenerate-symmetric-bilinear}
\sketch The map $V\to V^*$, $v\mapsto B(v,\cdot)$, is injective by nondegeneracy and hence bijective by equal finite dimensions; take the inverse image of $a$.
\end{proof}

\begin{lemma}[The $B$-gradient is orthogonal to the kernel of the functional]
\label{lem:bilinear-gradient-orthogonal-ker}
\lean{LeeLib.Ch02.bilinGrad_mem_orthogonal_ker}
\uses{def:bilinear-gradient}
\leanok



With the preceding notation, $a^\sharp\in(\ker a)^\perp$.
\end{lemma}
\begin{proof}
\leanok
\uses{def:bilinear-gradient}
\sketch For $w\in\ker a$, symmetry and Definition \ref{def:bilinear-gradient} give $B(w,a^\sharp)=B(a^\sharp,w)=a(w)=0$.
\end{proof}

\begin{proposition}[Signature drop on the kernel of a functional]
\label{prop:embedding-criterion-pointwise}
\lean{LeeLib.Ch02.sigPos_sigNeg_restrict_ker_bilinGrad}
\uses{def:bilinear-gradient, lem:bilinear-gradient-orthogonal-ker, prop:hypersurface-signature-pointwise-hyperplane}
\leanok



Let $B$ on $V$ have signature $(r,s)$ and let $0\ne a\in V^*$. If $B(a^\sharp,a^\sharp)>0$, then $B|_{\ker a}$ has signature $(r-1,s)$; if it is negative, the signature is $(r,s-1)$. This is the fibrewise form of Corollary \ref{cor:pseudo-riemannian-embedding-criterion}.
\end{proposition}
\begin{proof}
\leanok
\uses{def:bilinear-gradient,lem:bilinear-gradient-orthogonal-ker,prop:hypersurface-signature-pointwise-hyperplane}
\sketch Rank--nullity makes $\ker a$ a hyperplane. Lemma \ref{lem:bilinear-gradient-orthogonal-ker} puts the nonzero $a^\sharp$ in its complement, so Proposition \ref{prop:hypersurface-signature-pointwise-hyperplane} applies.
\end{proof}

\begin{lemma}[Nondegeneracy on the kernel of a functional]
\label{lem:embedding-criterion-nondegenerate-pointwise}
\lean{LeeLib.Ch02.nondegenerate_restrict_ker_bilinGrad}
\uses{def:bilinear-gradient, lem:bilinear-gradient-orthogonal-ker, lem:hyperplane-eq-complement-normal-line, lem:orthogonal-complement-non-null-line-nondegenerate}
\leanok



If $a\ne0$ and $B(a^\sharp,a^\sharp)\ne0$, then $B|_{\ker a}$ is nondegenerate.
\end{lemma}
\begin{proof}
\leanok
\uses{lem:bilinear-gradient-orthogonal-ker,lem:hyperplane-eq-complement-normal-line,lem:orthogonal-complement-non-null-line-nondegenerate}
\sketch Lemmas \ref{lem:bilinear-gradient-orthogonal-ker} and \ref{lem:hyperplane-eq-complement-normal-line} give $\ker a=\operatorname{span}(a^\sharp)^\perp$; apply Lemma \ref{lem:orthogonal-complement-non-null-line-nondegenerate}.
\end{proof}

\begin{corollary}
\label{cor:pseudo-riemannian-embedding-criterion}
\lean{LeeLib.Ch02.hasSignature_pseudoPullbackMetric_levelSet_of_pseudoGrad_pos, LeeLib.Ch02.hasSignature_pseudoPullbackMetric_levelSet_of_pseudoGrad_neg, LeeLib.Ch02.form_pseudoGrad_mfderiv_levelSet_val_eq_zero}
\uses{def:pseudo-riemannian-metric, def:pseudo-metric-has-signature, prop:embedding-criterion-pointwise, def:bilinear-gradient, prop:regular-level-set-hypersurface-codim-one, prop:regular-level-set-tangent-kernel}
\leanok


\dcref{ch2:2.71}

Let $(\bar M,\bar g)$ have signature $(r,s)$ and let $M=f^{-1}(c)$. If $\bar g(\operatorname{grad}f,\operatorname{grad}f)>0$ on $M$, then $M$ is an embedded pseudo-Riemannian submanifold of signature $(r-1,s)$; if it is negative, the signature is $(r,s-1)$. In both cases $\operatorname{grad}f$ is normal to $M$.
\end{corollary}
\begin{proof}
\leanok
\uses{prop:embedding-criterion-pointwise,prop:regular-level-set-hypersurface-codim-one,prop:regular-level-set-tangent-kernel}
\sketch Non-null gradient makes $df_p\ne0$. Propositions \ref{prop:regular-level-set-hypersurface-codim-one} and \ref{prop:regular-level-set-tangent-kernel} give an embedded hypersurface with $T_pM=\ker df_p$. Apply Proposition \ref{prop:embedding-criterion-pointwise} to $a=df_p$; Lemma \ref{lem:bilinear-gradient-orthogonal-ker} gives normality.
\end{proof}

\begin{lemma}[The Ambient Metric Along a Submanifold]
\label{lem:pullback-pseudo-metric-bundle}
\lean{Bundle.ContMDiffPseudoMetric.pullback}
\uses{def:smooth-field-scalar-products}
\leanok



If $g$ is a smooth field of scalar products on $V\to B$ and $F:B'\to B$ is smooth, then $(F^*g)_x=g_{F(x)}$ is a smooth field of scalar products on $F^*V$.
\end{lemma}
\begin{proof}
\leanok
\uses{def:smooth-field-scalar-products}
\sketch Symmetry and nondegeneracy hold fibrewise. In pulled-back trivializations, the coordinate field is the smooth coordinate field of $g$ composed with $F$.
\end{proof}

\begin{definition}[The Tensor Induced Along a Map]
\label{def:induced-tensor-along-map}
\lean{LeeLib.Ch02.IsPullbackAlong}
\uses{def:pseudo-riemannian-metric}
\leanok



For smooth $F:M\to\bar M$ and pseudo-Riemannian metrics $g,\bar g$, say $g$ is induced by $\bar g$ along $F$, and write $g=F^*\bar g$, when
\[
g_p(v,w)=\bar g_{F(p)}(dF_pv,dF_pw)
\]
for all $p,v,w$.
\end{definition}

\begin{lemma}[A Pseudo-Riemannian Submanifold is Immersed]
\label{lem:pseudo-submanifold-immersion}
\lean{LeeLib.Ch02.injective_mfderiv_of_isPullbackAlong}
\uses{def:induced-tensor-along-map}
\leanok



If the pseudo-Riemannian metric $g$ is induced along smooth $F$, then every $dF_p$ is injective.
\end{lemma}
\begin{proof}
\leanok
\uses{def:induced-tensor-along-map}
\sketch If $dF_pv=0$, then Definition \ref{def:induced-tensor-along-map} gives $g_p(v,w)=0$ for every $w$; nondegeneracy forces $v=0$.
\end{proof}

\begin{lemma}[A Riemannian Submanifold is a Pseudo-Riemannian Submanifold]
\label{lem:riemannian-submanifold-induced-tensor}
\lean{LeeLib.Ch02.isPullbackAlong_pullbackMetric}
\uses{def:induced-tensor-along-map, lem:pullback-immersion-metric, def:induced-metric}
\leanok



For a smooth immersion $F:M\to\bar M$ into a Riemannian manifold, the induced metric of Definition \ref{def:induced-metric} is induced along $F$ in the sense of Definition \ref{def:induced-tensor-along-map}.
\end{lemma}
\begin{proof}
\leanok
\uses{def:induced-tensor-along-map,lem:pullback-immersion-metric,def:induced-metric}
\sketch Lemma \ref{lem:pullback-immersion-metric} gives positive definiteness, and the pullback formula is exactly Definition \ref{def:induced-tensor-along-map}.
\end{proof}

\begin{definition}[Transport of a Field of Bilinear Forms by a Family of Linear Maps]
\label{def:bilinear-comp-of-family}
\lean{LeeLib.Ch02.bilinearCompOf}
\leanok


For a map $F:M\to\widetilde M$, bilinear forms $b_y$, and linear maps $A_p:T_pM\to T_{F(p)}\widetilde M$, define
\[
(A^*b)_p(v,w)=b_{F(p)}(A_pv,A_pw).
\]
\end{definition}

\begin{lemma}[Transport by a Smooth Family of Linear Maps is Smooth]
\label{lem:bilinear-comp-of-family-smooth}
\lean{LeeLib.Ch02.contMDiffAt_bilinearCompOf}
\uses{def:bilinear-comp-of-family}
\leanok



If $b$ is smooth, $F$ is smooth at $p_0$, and the tangent-coordinate representation of $A_p$ is smooth at $p_0$, then $p\mapsto(A^*b)_p$ is smooth at $p_0$.
\end{lemma}
\begin{proof}
\leanok
\uses{def:bilinear-comp-of-family}
\sketch In tangent coordinates the field is $p\mapsto B(F(p))(D(p)\cdot,D(p)\cdot)$, a smooth expression linear in $B$ and quadratic in the smooth family $D$.
\end{proof}

This transport also covers horizontal lifts used in Theorem \ref{thm:group-action-riemannian-submersion}, where $A_p$ need not be a differential.

\begin{definition}[The Pullback of a Field of Bilinear Forms]
\label{def:pullback-form-of-family}
\lean{LeeLib.Ch02.pullbackFormOf}
\uses{def:bilinear-comp-of-family}
\leanok



For $F:M\to\widetilde M$, the pullback of $b$ is transport by $A_p=dF_p$:
\[
(F^*b)_p(v,w)=b_{F(p)}(dF_pv,dF_pw).
\]
\end{definition}

\begin{lemma}[The Pullback of a Smooth Field of Forms is Smooth]
\label{lem:pullback-form-of-family-smooth}
\lean{LeeLib.Ch02.contMDiff_pullbackFormOf}
\uses{def:pullback-form-of-family, lem:bilinear-comp-of-family-smooth}
\leanok



If $F$ and the field $b$ are smooth, then $F^*b$ is a smooth field of bilinear forms on $TM$.
\end{lemma}
\begin{proof}
\leanok
\uses{def:pullback-form-of-family,lem:bilinear-comp-of-family-smooth}
\sketch Apply Lemma \ref{lem:bilinear-comp-of-family-smooth} pointwise to $A_p=dF_p$, whose tangent-coordinate representation varies smoothly.
\end{proof}

\begin{definition}[The Pullback of an Indefinite Metric]
\label{def:pseudo-pullback-form}
\lean{LeeLib.Ch02.pseudoPullbackForm}
\uses{def:pullback-form-of-family, def:pseudo-riemannian-metric}
\leanok



For pseudo-Riemannian $(\widetilde M,\widetilde g)$ and $F:M\to\widetilde M$, define
\[
(F^*\widetilde g)_p(v,w)=\widetilde g_{F(p)}(dF_pv,dF_pw).
\]
\end{definition}

\begin{lemma}[The Pullback Tensor is Symmetric]
\label{lem:pseudo-pullback-form-symm}
\lean{LeeLib.Ch02.pseudoPullbackForm_symm}
\uses{def:pseudo-pullback-form, def:pseudo-riemannian-metric}
\leanok



With the preceding notation, $(F^*\widetilde g)_p(v,w)=(F^*\widetilde g)_p(w,v)$.
\end{lemma}
\begin{proof}
\leanok
\uses{def:pseudo-pullback-form,def:pseudo-riemannian-metric}
\sketch Expand Definition \ref{def:pseudo-pullback-form} and use symmetry of $\widetilde g_{F(p)}$.
\end{proof}

\begin{lemma}[The Pullback Tensor is Smooth]
\label{lem:pseudo-pullback-form-smooth}
\lean{LeeLib.Ch02.pseudoPullbackForm_contMDiff}
\uses{def:pseudo-pullback-form, lem:pullback-form-of-family-smooth, def:smooth-field-scalar-products}
\leanok



If $F$ is smooth, then $p\mapsto(F^*\widetilde g)_p$ is smooth.
\end{lemma}
\begin{proof}
\leanok
\uses{def:pseudo-pullback-form,lem:pullback-form-of-family-smooth,def:smooth-field-scalar-products}
\sketch Definition \ref{def:smooth-field-scalar-products} makes $\widetilde g$ a smooth field, so Lemma \ref{lem:pullback-form-of-family-smooth} applies.
\end{proof}

\begin{definition}[The Metric Induced by an Indefinite Ambient Metric]
\label{def:pseudo-induced-metric}
\lean{LeeLib.Ch02.pseudoPullbackMetric}
\uses{def:pseudo-pullback-form, lem:pseudo-pullback-form-symm, lem:pseudo-pullback-form-smooth, def:pseudo-riemannian-metric, def:smooth-field-scalar-products}
\leanok



Let $F:M\to(\widetilde M,\widetilde g)$ be smooth. If $F^*\widetilde g$ is pointwise nondegenerate, meaning that each $0\ne v\in T_pM$ has $w\in T_pM$ with $(F^*\widetilde g)_p(v,w)\ne0$, then it is a pseudo-Riemannian metric on $M$, called the metric induced by $F$.
\end{definition}
\begin{proof}
\leanok
\uses{lem:pseudo-pullback-form-symm,lem:pseudo-pullback-form-smooth,def:pseudo-riemannian-metric,def:smooth-field-scalar-products}
\sketch Lemmas \ref{lem:pseudo-pullback-form-symm} and \ref{lem:pseudo-pullback-form-smooth}, together with the nondegeneracy hypothesis, verify the three conditions of Definition \ref{def:smooth-field-scalar-products}; apply Definition \ref{def:pseudo-riemannian-metric}.
\end{proof}

\begin{lemma}[A Pseudo-Riemannian Pullback is an Induced Tensor]
\label{lem:pseudo-riemannian-submanifold-induced-tensor}
\lean{LeeLib.Ch02.isPullbackAlong_pseudoPullbackMetric}
\uses{def:pseudo-induced-metric, def:pseudo-pullback-form, def:induced-tensor-along-map}
\leanok



Under the hypotheses of Definition \ref{def:pseudo-induced-metric}, the metric $F^*\widetilde g$ is induced along $F$ in the sense of Definition \ref{def:induced-tensor-along-map}.
\end{lemma}
\begin{proof}
\leanok
\uses{def:pseudo-induced-metric,def:pseudo-pullback-form,def:induced-tensor-along-map}
\sketch Definition \ref{def:pseudo-induced-metric} supplies the metric, and Definition \ref{def:pseudo-pullback-form} is exactly the identity required by Definition \ref{def:induced-tensor-along-map}.
\end{proof}

\begin{lemma}[Extension of a Nondegenerate Family to a Nondegenerate Frame]
\label{lem:extend-nondegenerate-family}
\lean{LeeLib.Ch02.exists_isLocalNondegenerateOn_extend}
\uses{def:local-nondegenerate-family, lem:completion-nondegenerate-bases, lem:nondegenerate-tuple-independent, lem:nondegenerate-tuple-basis, lem:sections-through-prescribed-basis, lem:open-nondegenerate-tuple}
\leanok



Let $g$ be a smooth field of scalar products on a rank-$m$ bundle $V\to B$. If $X_1,\ldots,X_k$ are nondegenerate sections on open $u$ and $x_0\in u$, then $k\le m$ and, on some $x_0\in v\subseteq u$, they extend to nondegenerate sections $Y_1,\ldots,Y_m$ with $Y_i=X_i$ for $i\le k$.
\end{lemma}
\begin{proof}
\leanok
\uses{def:local-nondegenerate-family,lem:completion-nondegenerate-bases,lem:nondegenerate-tuple-independent,lem:nondegenerate-tuple-basis,lem:sections-through-prescribed-basis,lem:open-nondegenerate-tuple}
\sketch At $x_0$, use Lemmas \ref{lem:nondegenerate-tuple-independent}, \ref{lem:completion-nondegenerate-bases}, and \ref{lem:nondegenerate-tuple-basis} to extend the tuple to a basis $(w_j)$. Lemma \ref{lem:sections-through-prescribed-basis} realizes the extra $w_j$ by local sections; retain the original first $k$ sections and shrink using Lemma \ref{lem:open-nondegenerate-tuple}.
\end{proof}

\begin{proposition}[Pseudo-Riemannian Adapted Orthonormal Frames]
\label{prop:pseudo-riemannian-adapted-frame}
\lean{LeeLib.Ch02.exists_adapted_pseudo_orthonormalFrame}
\uses{def:pseudo-riemannian-metric, def:induced-tensor-along-map, lem:pullback-pseudo-metric-bundle, prop:orthonormal-frames-pseudo-riemannian, lem:extend-nondegenerate-family, prop:orthonormalize-nondegenerate-family, lem:orthonormal-tuple-nondegenerate, lem:pseudo-submanifold-immersion}
\leanok


\dcref{ch2:2.72}

If $F:M\to(\bar M,\bar g)$ presents a pseudo-Riemannian or Riemannian submanifold with induced $g$, then every $p\in M$ has a neighborhood with an orthonormal frame of $F^*T\bar M$ whose first $n=\dim M$ members span $dF_x(T_xM)$ at every $x$.
\end{proposition}
\begin{proof}
\leanok
\uses{def:induced-tensor-along-map,lem:pullback-pseudo-metric-bundle,prop:orthonormal-frames-pseudo-riemannian,lem:extend-nondegenerate-family,prop:orthonormalize-nondegenerate-family,lem:orthonormal-tuple-nondegenerate,lem:pseudo-submanifold-immersion}
\sketch Equip $F^*T\bar M$ with Lemma \ref{lem:pullback-pseudo-metric-bundle}; Lemma \ref{lem:pseudo-submanifold-immersion} gives injectivity. Push forward a local orthonormal frame of $(M,g)$ from Proposition \ref{prop:orthonormal-frames-pseudo-riemannian}; inducedness makes it orthonormal and Lemma \ref{lem:orthonormal-tuple-nondegenerate} makes it nondegenerate. Extend by Lemma \ref{lem:extend-nondegenerate-family}, orthonormalize by Proposition \ref{prop:orthonormalize-nondegenerate-family}, and use preservation of the span flag plus dimension $n$ to obtain adaptedness.
\end{proof}

\begin{definition}[The Normal Space of a Pseudo-Riemannian Submanifold]
\label{def:pseudo-normal-space}
\lean{LeeLib.Ch02.pseudoNormalSpace}
\uses{def:pseudo-riemannian-metric, lem:pullback-pseudo-metric-bundle}
\leanok



For an immersion $F:M^n\to(\bar M^m,\bar g)$, define
\[
N_xM=\{v\in T_{F(x)}\bar M:\bar g_{F(x)}(v,w)=0\text{ for all }w\in dF_x(T_xM)\}.
\]
\end{definition}

\begin{lemma}[Local Frames for the Normal Space]
\label{lem:pseudo-normal-space-frame}
\lean{LeeLib.Ch02.exists_pseudo_localFrame_normalSpace}
\uses{def:pseudo-normal-space, prop:pseudo-riemannian-adapted-frame, lem:orthonormal-tuple-nondegenerate}
\leanok



If $F:M^n\to\bar M^m$ presents a pseudo-Riemannian or Riemannian submanifold, then every point has a neighborhood with $m-n$ smooth sections whose values are orthonormal and span $N_xM$.
\end{lemma}
\begin{proof}
\leanok
\uses{def:pseudo-normal-space,prop:pseudo-riemannian-adapted-frame,lem:orthonormal-tuple-nondegenerate}
\sketch For an adapted orthonormal frame $(b_1,\ldots,b_m)$ from Proposition \ref{prop:pseudo-riemannian-adapted-frame}, take $N_j=b_{n+j}$. They are normal by orthogonality. Expanding $w\in N_xM$ in the full basis and pairing with $b_a$ for $a\le n$ gives $c_a\langle b_a,b_a\rangle=0$; since each square is $\pm1$, all tangential coefficients vanish.
\end{proof}

\begin{corollary}[Normal Bundle to a Pseudo-Riemannian Submanifold]
\label{cor:normal-bundle-pseudo-riemannian-submanifold}
\lean{LeeLib.Ch02.contMDiffVectorBundle_pseudoNormalSpace}
\uses{def:pseudo-riemannian-metric, def:pseudo-normal-space, prop:pseudo-riemannian-adapted-frame, lem:pseudo-normal-space-frame, lem:local-frame-criterion-for-subbundles, prop:existence-riemannian-metric}
\leanok


\dcref{ch2:2.73}

For an embedded or immersed pseudo-Riemannian or Riemannian submanifold $F:M^n\to\bar M^m$, the normal spaces form a smooth rank-$(m-n)$ vector subbundle of $F^*T\bar M$, the normal bundle.
\end{corollary}
\begin{proof}
\leanok
\uses{def:pseudo-normal-space,lem:pseudo-normal-space-frame,lem:local-frame-criterion-for-subbundles,prop:existence-riemannian-metric}
\sketch Choose an auxiliary Riemannian metric on $\bar M$ by Proposition \ref{prop:existence-riemannian-metric}; Lemma \ref{lem:ambient-tangent-bundle-riemannian} makes $F^*T\bar M$ a Riemannian vector bundle. Lemma \ref{lem:pseudo-normal-space-frame} supplies smooth local bases of $N_xM$, so Lemma \ref{lem:local-frame-criterion-for-subbundles} yields the rank-$(m-n)$ subbundle. The auxiliary metric does not alter the normal spaces defined by $\bar g$.
\end{proof}

\section{Other Generalizations of Riemannian Metrics}

\subsection{Sub-Riemannian Metrics}

A sub-Riemannian metric on $M$ is a smooth fibre metric on a smooth tangent distribution $S\subseteq TM$. Only horizontal curves, those satisfying $\gamma'(t)\in S_{\gamma(t)}$, have defined length. One generally requires an accessibility condition ensuring that nearby points can be joined by horizontal curves. For a CR submanifold $M\subseteq\mathbb C^n$, a standard distribution is
\[
S_p=T_pM\cap iT_pM,
\]
equipped with a fibre metric.

\subsection{Finsler Metrics}

A norm on a real vector space satisfies
\begin{enumerate}
\item $|cv|=|c|\,|v|$;
\item $|v|\ge0$, with equality exactly when $v=0$;
\item $|v+w|\le|v|+|w|$.
\end{enumerate}
A Finsler metric on a smooth manifold is a continuous $F:TM\to\mathbb R$, smooth away from the zero section, whose restriction to every $T_pM$ is a norm. The norm induced by a Riemannian metric is a Finsler metric.

% The following source references are retained from condensed dependency bridges.
% \ref{cor:normal-bundle-pseudo-riemannian-submanifold}
% \ref{def:bilinear-comp-of-family}
% \ref{def:induced-metric}
% \ref{def:induced-metric}
% \ref{def:induced-tensor-along-map}
% \ref{def:pseudo-induced-metric}
% \ref{def:pseudo-pullback-form}
% \ref{def:pseudo-pullback-form}
% \ref{def:pseudo-pullback-form}
% \ref{def:pullback-form-of-family}
% \ref{lem:hyperplane-nondegenerate}
% \ref{lem:hyperplane-nondegenerate}
% \ref{lem:hyperplane-nondegenerate}
% \ref{lem:local-frame-criterion-for-subbundles}
% \ref{lem:local-frame-criterion-for-subbundles}
% \ref{lem:local-frame-criterion-for-subbundles}
% \ref{lem:orthogonal-complement-involution}
% \ref{lem:pseudo-normal-space-frame}
% \ref{lem:pseudo-normal-space-line}
% \ref{lem:pullback-immersion-metric}
% \ref{lem:pullback-immersion-metric}
% \ref{lem:riemannian-submanifold-induced-tensor}
% \ref{lem:riemannian-submanifold-induced-tensor}
% \ref{prop:adapted-orthonormal-frames-immersed}
% \ref{prop:existence-riemannian-metric}
% \ref{prop:hypersurface-signature-at-a-point}
% \ref{prop:hypersurface-signature-pointwise-hyperplane}
% \ref{prop:hypersurface-signature-pointwise}
% \ref{prop:normal-bundle}
% \ref{prop:pseudo-riemannian-adapted-frame}
% \ref{prop:pseudo-riemannian-adapted-frame}
% \ref{prop:pseudo-riemannian-hypersurface-nondegenerate}
% \ref{prop:pseudo-riemannian-hypersurface-signature}
% \ref{prop:pseudo-riemannian-hypersurface-signature}
% \ref{prop:pseudo-riemannian-hypersurface-signature}
% \ref{prop:pseudo-riemannian-hypersurface-signature}
% \ref{prop:pseudo-riemannian-hypersurface-signature}
% \ref{prop:pseudo-riemannian-hypersurface-signature}
% \ref{prop:pseudo-riemannian-hypersurface-signature}
% \ref{prop:pseudo-riemannian-hypersurface-signature}
% \ref{prop:pseudo-riemannian-hypersurface-signature}
% \ref{prop:pseudo-riemannian-hypersurface-signature}
% \ref{prop:pseudo-riemannian-hypersurface-signature}
% \ref{prop:pseudo-riemannian-hypersurface-signature}
% \ref{prop:pseudo-riemannian-hypersurface-signature}
% \ref{prop:sylvesters-law-inertia}
\section{Further Results}

\begin{proposition}\label{prop:distance-preserving-map-linear-isometry}\dcref{ch2:2-2}


Let $V,W$ be finite-dimensional real inner product spaces of equal dimension.
If $F:V\to W$ satisfies $F(0)=0$ and
$|F(x)-F(y)|=|x-y|$ for all $x,y$, then $F$ is a linear isometry.
\end{proposition}

\begin{definition}[The open submanifold metric]
\label{def:open-submanifold-metric}
\lean{LeeLib.Ch02.openSubmanifoldMetric}
\uses{def:riemannian-metric, def:induced-metric, lem:mfderiv-open-inclusion}
\leanok



For an open $W\subseteq\mathbb R^n$ and a Riemannian metric $g$ on $\mathbb R^n$, the metric induced on $W$ is $\iota^*g$ for the inclusion $\iota:W\hookrightarrow\mathbb R^n$.
\end{definition}

\begin{lemma}[The differential of the inclusion of an open subset]
\label{lem:mfderiv-open-inclusion}
\lean{LeeLib.Ch02.mfderiv_opens_subtypeVal}
\leanok


An open $W\subseteq\mathbb R^n$ is a smooth manifold with $T_xW\cong\mathbb R^n$, and under this identification $d\iota_x=\operatorname{id}$; hence $\iota$ is an immersion.
\end{lemma}

\begin{proof}
\leanok
\sketch In the single inclusion chart of $W$, the coordinate representative is the identity on a neighbourhood of each point, so its differential is the identity.
\end{proof}

\begin{lemma}[Tangent directions to the sphere are orthogonal to the radius]
\label{lem:sphere-tangent-orthogonal-radius}
\lean{LeeLib.Ch02.inner_coe_sphere_mfderiv_eq_zero}
\uses{ex:spheres}
\leanok



For $\iota:S^{n-1}\hookrightarrow\mathbb R^n$, $\langle\omega,d\iota_\omega(w)\rangle=0$ for every $\omega\in S^{n-1}$ and $w\in T_\omega S^{n-1}$.
\end{lemma}

\begin{proof}
\leanok
\sketch The image of $d\iota_\omega$ is $(\mathbb R\omega)^\perp$.
\end{proof}

\begin{definition}[The polar map]
\label{def:polar-map}
\lean{LeeLib.Ch02.polarMap}
\uses{def:warped-product, ex:spheres, def:open-submanifold-metric}
\leanok



Define $\Phi:\mathbb R^+\times S^{n-1}\to\mathbb R^n$ by $\Phi(\rho,\omega)=\rho\omega$.
\end{definition}

\begin{lemma}[The differential of the polar map]
\label{lem:mfderiv-polar-map}
\lean{LeeLib.Ch02.mfderiv_polarMap_apply}
\uses{def:polar-map, lem:mfderiv-open-inclusion}
\leanok



For $(r,\omega)\in\mathbb R^+\times S^{n-1}$ and $(a,w)\in T_r\mathbb R^+\oplus T_\omega S^{n-1}$,
\[
d\Phi_{(r,\omega)}(a,w)=a\omega+r\,d\iota_\omega(w).
\]
\end{lemma}

\begin{proof}
\leanok
\sketch Differentiate the two factors separately: the radial map contributes $a\omega$ and the fixed-radius scaling contributes $r,d\iota_\omega(w)$.
\end{proof}

\begin{proposition}\label{prop:warped-product-isometry}\dcref{ch2:2-4}


\uses{ex:warped-products,def:polar-map,lem:mfderiv-polar-map,lem:sphere-tangent-orthogonal-radius,def:metric-preserving,def:open-submanifold-metric,prop:warped-product-constant-one}
For $\rho(r)=r$, the map
$\Phi:\mathbb R^+\times_\rho S^{n-1}\to\mathbb R^n\setminus\{0\}$,
$\Phi(\rho,\omega)=\rho\omega$, is an isometry.
\end{proposition}

\begin{proof}
\leanok
\uses{def:polar-map,lem:mfderiv-polar-map,lem:sphere-tangent-orthogonal-radius,def:metric-preserving}
\sketch Using Lemma \ref{lem:mfderiv-polar-map}, the Euclidean pairing of two differentials is
\[
aa'+ar\langle\omega,d\iota(w')\rangle+a'r\langle d\iota(w),\omega\rangle+r^2\langle d\iota(w),d\iota(w')\rangle
=aa'+r^2\mathring g_\omega(w,w')
\]
by Lemma \ref{lem:sphere-tangent-orthogonal-radius}; this is the warped metric. The map is bijective with inverse $x\mapsto(|x|,x/|x|)$.
\end{proof}

\begin{proposition}\label{prop:grassmann-stiefel-manifolds}\dcref{ch2:2-7}
For $0\le k\le n$, let $G_k(\mathbb R^n)$ be the Grassmannian and let $V_k(\mathbb R^n)$ be the set of orthonormal ordered $k$-frames, viewed as matrices in $M(n\times k,\mathbb R)$.
\begin{enumerate}
\item[(a)] $V_k(\mathbb R^n)$ is a smooth submanifold of dimension
$k(2n-k-1)/2$ (the Stiefel manifold).
\item[(b)] The span map $\pi:V_k(\mathbb R^n)\to G_k(\mathbb R^n)$ is a
surjective smooth submersion.
\item[(c)] With the induced Euclidean metric on $V_k(\mathbb R^n)$, the right
$O(k)$-action is isometric, vertical, and fiber-transitive. Consequently,
$G_k(\mathbb R^n)$ carries a unique metric making $\pi$ a Riemannian
submersion.
\end{enumerate}
\end{proposition}

\begin{exercise}\label{prob:grassmann-stiefel-manifolds}
\dcref{ch2:2-7}
For $0 \le k \le n$, the set $G_k(\mathbb{R}^n)$ of $k$-dimensional linear subspaces of $\mathbb{R}^n$ is called a \textbf{Grassmann manifold} or \textbf{Grassmannian}. The group $GL(n,\mathbb{R})$ acts transitively on $G_k(\mathbb{R}^n)$ in an obvious way, and $G_k(\mathbb{R}^n)$ has a unique smooth manifold structure making this action smooth (see [LeeSM, Example 21.21]).

(a) Let $V_k(\mathbb{R}^n)$ denote the set of orthonormal ordered $k$-tuples of vectors in $\mathbb{R}^n$. By arranging the vectors in $k$ columns, we can view $V_k(\mathbb{R}^n)$ as a subset of the vector space $M(n \times k, \mathbb{R})$ of all $n \times k$ real matrices. Prove that $V_k(\mathbb{R}^n)$ is a smooth submanifold of $M(n \times k, \mathbb{R})$ of dimension $k(2n - k - 1)/2$, called a \textbf{Stiefel manifold}. [Hint: Consider the map $\Phi: M(n \times k, \mathbb{R}) \to M(k \times k, \mathbb{R})$ given by $\Phi(A) = A^T A$.]

(b) Show that the map $\pi: V_k(\mathbb{R}^n) \to G_k(\mathbb{R}^n)$ that sends a $k$-tuple to its span is a surjective smooth submersion.

(c) Give $V_k(\mathbb{R}^n)$ the Riemannian metric induced from the Euclidean metric on $M(n \times k, \mathbb{R})$. Show that the right action of $O(k)$ on $V_k(\mathbb{R}^n)$ by matrix multiplication on the right is isometric, vertical, and transitive on fibers of $\pi$, and thus there is a unique metric on $G_k(\mathbb{R}^n)$ such that $\pi$ is a Riemannian submersion. [Hint: It might help to note that the Euclidean inner product on $M(n \times k, \mathbb{R})$ can be written in the form $(A,B) = \operatorname{tr}(A^T B)$.]
\end{exercise}


\begin{proposition}\label{prop:gradient-characterization}\dcref{ch2:2-10}


\uses{def:gradient,def:regular-set,lem:gradient-characterization,lem:transversal-representative-eq-gradient}
Let $(M,g)$ be Riemannian, $f\in C^\infty(M)$, and let $X$ be nowhere zero.
Then
\[
X=\operatorname{grad}f\quad\Longleftrightarrow\quad Xf=|X|_g^2\text{ and }X\perp\text{ every regular level set of }f.
\]
\end{proposition}

\begin{proof}
\leanok
\sketch If $X=\operatorname{grad}f$, the gradient characterization gives both conditions. Conversely, $df_p(X_p)=|X_p|^2>0$ makes every point regular; the orthogonality hypothesis and Lemma \ref{lem:transversal-representative-eq-gradient} then give $X_p=\operatorname{grad}f|_p$.
\end{proof}

\begin{lemma}[A transversal orthogonal representative is the gradient]
\label{lem:transversal-representative-eq-gradient}
\lean{LeeLib.Ch02.eq_grad_of_ne_zero_of_innerAt_self_of_orthogonal_ker}
\uses{def:gradient, lem:gradient-characterization}
\leanok



If $v\in T_pM$ is nonzero, $\langle v,v\rangle_g=df_p(v)$, and $\langle v,w\rangle_g=0$ for all $w\in\ker df_p$, then $v=\operatorname{grad}f|_p$.
\end{lemma}

\begin{proof}
\leanok
\sketch Since $df_p(v)>0$, for arbitrary $w$ set $t=df_p(w)/df_p(v)$. Then $w-tv\in\ker df_p$, so orthogonality gives $\langle v,w\rangle_g=t\langle v,v\rangle_g=df_p(w)$, which characterizes the gradient.
\end{proof}

\begin{definition}
\label{def:inner-k-covectors}
\lean{LeeLib.Ch02.innerForms}
\leanok


For an orthonormal basis $(e_m)_{m\in\iota'}$ and $\omega,\eta\in\Lambda^k(V^*)$, define
\[
\langle\omega,\eta\rangle_e=\frac1{k!}\sum_{s:\{1,\ldots,k\}\to\iota'}\omega(e_{s(1)},\ldots,e_{s(k)})\eta(e_{s(1)},\ldots,e_{s(k)}).
\]
\end{definition}

\begin{lemma}
\label{lem:inner-k-covectors-posdef}
\lean{LeeLib.Ch02.innerForms_self_pos}
\uses{def:inner-k-covectors}
\leanok



The displayed form is a symmetric bilinear inner product on $\Lambda^k(V^*)$.
\end{lemma}
\begin{proof}
\leanok
\sketch Symmetry and bilinearity are termwise. The self-pairing is a sum of squares divided by $k!$; a nonzero alternating tensor has a nonzero value on some basis tuple.
\end{proof}

\begin{lemma}
\label{lem:inner-k-covectors-wedge}
\lean{LeeLib.Ch02.innerForms_wedgeCovectors}
\uses{def:inner-k-covectors, lem:alternating-cauchy-binet}
\leanok



For covectors $\omega^1,\ldots,\omega^k,\eta^1,\ldots,\eta^k$,
\[
\langle\omega^1\wedge\cdots\wedge\omega^k,\eta^1\wedge\cdots\wedge\eta^k\rangle_e=\det(\langle\omega^i,\eta^j\rangle).
\]
\end{lemma}
\begin{proof}
\leanok
\uses{lem:alternating-cauchy-binet}
\sketch Apply the alternating Cauchy--Binet identity to $A_{im}=\omega^i(e_m)$ and use $\sum_m\omega^i(e_m)e_m=(\omega^i)^\sharp$; the factors of $k!$ cancel.
\end{proof}

\begin{lemma}
\label{lem:k-covector-wedge-expansion}
\lean{LeeLib.Ch02.factorial_smul_eq_sum_wedgeCovectors}
\uses{lem:alternating-cauchy-binet}
\leanok



For every $\omega\in\Lambda^k(V^*)$,
\[
k!\,\omega=\sum_{u:\{1,\ldots,k\}\to\iota'}\omega(e_{u(1)},\ldots,e_{u(k)})\,e_{u(1)}^\flat\wedge\cdots\wedge e_{u(k)}^\flat.
\]
\end{lemma}
\begin{proof}
\leanok
\uses{lem:alternating-cauchy-binet}
\sketch Evaluate both sides on basis tuples and apply alternating Cauchy--Binet with the matrix $\langle e_{v(i)},e_m\rangle$.
\end{proof}

\begin{lemma}
\label{lem:span-wedge-covectors}
\lean{LeeLib.Ch02.span_range_wedgeCovectors}
\uses{lem:k-covector-wedge-expansion}
\leanok



Wedges of $k$ covectors span $\Lambda^k(V^*)$.
\end{lemma}
\begin{proof}
\leanok
\uses{lem:k-covector-wedge-expansion}
\sketch The expansion lemma expresses every $\omega$ as a linear combination of such wedges, after division by $k!$.
\end{proof}

\begin{proposition}
\label{prop:inner-k-covectors-unique}
\lean{LeeLib.Ch02.eq_innerForms_of_wedgeCovectors}
\uses{def:inner-k-covectors, lem:span-wedge-covectors, lem:inner-k-covectors-wedge}
\leanok



The form $\langle\cdot,\cdot\rangle_e$ is the unique bilinear form satisfying
\[
B(\omega^1\wedge\cdots\wedge\omega^k,\eta^1\wedge\cdots\wedge\eta^k)=\det(\langle\omega^i,\eta^j\rangle).
\]
\end{proposition}
\begin{proof}
\leanok
\uses{lem:span-wedge-covectors,lem:inner-k-covectors-wedge}
\sketch Lemma \ref{lem:inner-k-covectors-wedge} proves existence, and Lemma \ref{lem:span-wedge-covectors} makes the displayed values determine any bilinear form.
\end{proof}

\begin{proposition}
\label{prop:inner-k-covectors-frame-independent}
\lean{LeeLib.Ch02.innerForms_eq_innerForms}
\uses{prop:inner-k-covectors-unique, lem:inner-k-covectors-wedge}
\leanok



The form $\langle\cdot,\cdot\rangle_e$ is independent of the orthonormal basis $e$.
\end{proposition}
\begin{proof}
\leanok
\uses{prop:inner-k-covectors-unique}
\sketch Apply Proposition \ref{prop:inner-k-covectors-unique} to the forms obtained from any two orthonormal bases.
\end{proof}

\begin{definition}[The Fibre Metric on $\Lambda^k T^*M$, Pointwise]
\label{def:fiber-metric-forms-pointwise}
\lean{LeeLib.Ch02.RiemannianMetric.innerFormsAt}
\uses{def:riemannian-metric, def:inner-k-covectors, prop:inner-k-covectors-frame-independent}
\leanok



For a Riemannian manifold $(M,g)$ and $p\in M$, define the fibre metric on $\Lambda^k(T_p^*M)$ to be $\langle\cdot,\cdot\rangle_e$ for any $g_p$-orthonormal basis $e$ of $T_pM$.
\end{definition}

\begin{lemma}[The Fibre Metric is an Inner Product]
\label{lem:fiber-metric-forms-inner-product}
\lean{LeeLib.Ch02.RiemannianMetric.innerFormsAt_comm, LeeLib.Ch02.RiemannianMetric.innerFormsAt_add_left, LeeLib.Ch02.RiemannianMetric.innerFormsAt_smul_left, LeeLib.Ch02.RiemannianMetric.innerFormsAt_self_nonneg, LeeLib.Ch02.RiemannianMetric.innerFormsAt_self_pos}
\uses{def:fiber-metric-forms-pointwise, lem:inner-k-covectors-posdef}
\leanok



The fibre metric is an inner product on each $\Lambda^k(T_p^*M)$.
\end{lemma}
\begin{proof}
\leanok
\uses{lem:inner-k-covectors-posdef}
\sketch It is the inner product of Definition \ref{def:inner-k-covectors} for an orthonormal basis of $(T_pM,g_p)$.
\end{proof}

\begin{lemma}[The Fibre Metric Satisfies (2.26)]
\label{lem:fiber-metric-forms-wedge}
\lean{LeeLib.Ch02.RiemannianMetric.innerFormsAt_wedgeCovectors}
\uses{def:fiber-metric-forms-pointwise, lem:inner-k-covectors-wedge, def:sharp-map, def:wedge-covectors}
\leanok



For covectors at $p$,
\[
\langle\omega^1\wedge\cdots\wedge\omega^k,\eta^1\wedge\cdots\wedge\eta^k\rangle_g=\det(\langle\omega^i,\eta^j\rangle_g),
\]
where $\langle a,b\rangle_g=\langle a^\sharp,b^\sharp\rangle_g$.
\end{lemma}
\begin{proof}
\leanok
\uses{lem:inner-k-covectors-wedge}
\sketch This is Lemma \ref{lem:inner-k-covectors-wedge} applied to $(T_pM,g_p)$ and the Riesz representatives from Definition \ref{def:sharp-map}.
\end{proof}

\begin{lemma}[Fibrewise Uniqueness]
\label{lem:fiber-metric-forms-unique-pointwise}
\lean{LeeLib.Ch02.RiemannianMetric.eq_innerFormsAt_of_wedgeCovectors}
\uses{def:fiber-metric-forms-pointwise, prop:inner-k-covectors-unique}
\leanok



Any bilinear form on $\Lambda^k(T_p^*M)$ satisfying the determinant identity of Lemma \ref{lem:fiber-metric-forms-wedge} equals the fibre metric.
\end{lemma}
\begin{proof}
\leanok
\uses{prop:inner-k-covectors-unique}
\sketch Apply the pointwise uniqueness Proposition \ref{prop:inner-k-covectors-unique} to $(T_pM,g_p)$.
\end{proof}

\begin{lemma}[The Fibre Metric Against a Local Orthonormal Frame]
\label{lem:fiber-metric-forms-frame}
\lean{LeeLib.Ch02.RiemannianMetric.innerFormsAt_eq_sum_frame}
\uses{def:fiber-metric-forms-pointwise, def:inner-k-covectors, prop:inner-k-covectors-frame-independent, prop:existence-orthonormal-frames}
\leanok



If $(Y_1,\ldots,Y_n)$ is a smooth local orthonormal frame and $x$ lies in its domain, then
\[
\langle\omega,\eta\rangle_g=\frac1{k!}\sum_{s:\{1,\ldots,k\}\to\{1,\ldots,n\}}\omega(Y_{s(1)},\ldots,Y_{s(k)})\eta(Y_{s(1)},\ldots,Y_{s(k)}).
\]
\end{lemma}
\begin{proof}
\leanok
\uses{prop:inner-k-covectors-frame-independent}
\sketch The frame values are an orthonormal basis, so the formula is Definition \ref{def:inner-k-covectors} transported fibrewise.
\end{proof}

\begin{lemma}[A Smooth $k$-Form Field Applied to Smooth Vector Fields]
\label{lem:form-field-applied-to-vector-fields-smooth}
\lean{LeeLib.Ch02.contMDiffAt_apply_section}
\leanok


If $\omega$ is a smooth section of $\Lambda^kT^*M$ at $x_0$ and $X_1,\ldots,X_k$ are smooth vector fields there, then $x\mapsto\omega_x(X_1|_x,\ldots,X_k|_x)$ is smooth at $x_0$.
\end{lemma}
\begin{proof}
\leanok
\sketch In a local trivialization this is evaluation of a smooth family of alternating tensors on smooth vector-valued functions.
\end{proof}

\begin{lemma}[The Fibre Metric Pairs Smooth Fields Smoothly]
\label{lem:fiber-metric-forms-smooth}
\lean{LeeLib.Ch02.RiemannianMetric.contMDiffAt_innerFormsAt, LeeLib.Ch02.RiemannianMetric.contMDiff_innerFormsAt}
\uses{def:fiber-metric-forms-pointwise, lem:fiber-metric-forms-frame, prop:existence-orthonormal-frames, lem:form-field-applied-to-vector-fields-smooth}
\leanok



If $\omega,\eta$ are smooth sections of $\Lambda^kT^*M$, then $x\mapsto\langle\omega_x,\eta_x\rangle_g$ is smooth.
\end{lemma}
\begin{proof}
\leanok
\uses{lem:fiber-metric-forms-frame,prop:existence-orthonormal-frames,lem:form-field-applied-to-vector-fields-smooth}
\sketch Use the finite-frame sum of Lemma \ref{lem:fiber-metric-forms-frame}; each factor is a smooth evaluation.
\end{proof}

\begin{proposition}[The Fibre Metric on $\Lambda^k T^*M$]
\label{prop:fiber-metric-forms}
\lean{LeeLib.Ch02.RiemannianMetric.riemannian_fiberMetric_forms, LeeLib.Ch02.RiemannianMetric.innerFormsAt_self_pos, LeeLib.Ch02.RiemannianMetric.innerFormsAt_comm, LeeLib.Ch02.RiemannianMetric.contMDiff_innerFormsAt}
\uses{def:fiber-metric-forms-pointwise, lem:fiber-metric-forms-inner-product, lem:fiber-metric-forms-wedge, lem:fiber-metric-forms-unique-pointwise, lem:fiber-metric-forms-smooth, def:wedge-covectors}
\leanok



The fibre metric is the unique family of bilinear forms satisfying the determinant identity on wedges; each fibre is an inner product and pairings of smooth form fields are smooth.
\end{proposition}
\begin{proof}
\leanok
\uses{lem:fiber-metric-forms-wedge,lem:fiber-metric-forms-unique-pointwise,lem:fiber-metric-forms-inner-product,lem:fiber-metric-forms-smooth}
\sketch The preceding lemmas give the identity, inner-product properties, smoothness, and pointwise uniqueness.
\end{proof}

\begin{definition}[The Wedge Product of Alternating Tensors]
\label{def:wedge-alternating-forms}
\lean{LeeLib.Ch02.wedgeSum, LeeLib.Ch02.wedge}
\leanok


For continuous alternating $\omega\in\Lambda^k(V^*)$ and $\xi\in\Lambda^l(V^*)$, define
\[
(\omega\wedge\xi)(v_1,\ldots,v_{k+l})=\frac1{k!\,l!}\sum_{\sigma\in S_{k+l}}\operatorname{sgn}(\sigma)\,\omega(v_{\sigma(1)},\ldots,v_{\sigma(k)})\xi(v_{\sigma(k+1)},\ldots,v_{\sigma(k+l)}).
\]
\end{definition}

\begin{lemma}[Bilinearity of the Wedge Product]
\label{lem:wedge-bilinear}
\lean{LeeLib.Ch02.wedge_add_left, LeeLib.Ch02.wedge_smul_left, LeeLib.Ch02.wedge_add_right, LeeLib.Ch02.wedge_smul_right}
\uses{def:wedge-alternating-forms}
\leanok



The wedge product is bilinear in both arguments.
\end{lemma}
\begin{proof}
\leanok
\sketch Distribute the finite permutation sum; every summand is linear in each factor.
\end{proof}

\begin{lemma}[Commutativity of the Wedge over a Disjoint Union]
\label{lem:wedge-comm-sum}
\lean{LeeLib.Ch02.wedgeSum_comm}
\uses{def:wedge-alternating-forms}
\leanok



When arguments are indexed by a disjoint union, swapping the two index blocks gives $\omega\wedge\xi=\xi\wedge\omega$ after re-indexing. The usual $(-1)^{kl}$ appears only after both blocks are placed in one ordered index set.
\end{lemma}
\begin{proof}
\leanok
\sketch Re-index the permutation sum by conjugation with the block swap; permutation signs and the normalization agree.
\end{proof}

\begin{lemma}[The Wedge of Wedges of Covectors]
\label{lem:wedge-covectors-concat}
\lean{LeeLib.Ch02.wedge_wedgeCovectors, LeeLib.Ch02.wedgeSum_wedgeCovectors}
\uses{def:wedge-alternating-forms, def:wedge-covectors, lem:laplace-expansion-block}
\leanok



\[
(f_1\wedge\cdots\wedge f_k)\wedge(g_1\wedge\cdots\wedge g_l)=f_1\wedge\cdots\wedge f_k\wedge g_1\wedge\cdots\wedge g_l.
\]
\end{lemma}
\begin{proof}
\leanok
\uses{lem:laplace-expansion-block}
\sketch Evaluate on a tuple. The left side is the block Laplace expansion of the determinant defining the concatenated wedge (Lemma \ref{lem:laplace-expansion-block}).
\end{proof}

\begin{lemma}[The Defining Identity of the Hodge Star]
\label{lem:hodge-scalar-identity}
\lean{LeeLib.Ch02.wedgeRef, LeeLib.Ch02.sum_wedgeCovectors_flatL_mul, LeeLib.Ch02.sum_wedgeRef_mul_wedgeRef}
\uses{def:wedge-alternating-forms, def:wedge-covectors, def:inner-k-covectors, lem:inner-k-covectors-wedge, lem:perm-block-collapse, lem:perm-extension-count}
\leanok



For the block data $(\iota_a,\iota_b,\varepsilon)$ and $B(\eta,\xi)=(\eta\wedge\xi)(e\circ\varepsilon)$,
\[
\sum_{t:\iota_b\to\iota'}B(\eta,E^t)B(\omega,E^t)=l!\,\langle\omega,\eta\rangle_e.
\]
\end{lemma}
\begin{proof}
\leanok
\sketch Expand both wedges by the permutation definition. The sum over $t$ is an orthonormal determinant sum, and the remaining permutation sum collapses by the block-permutation lemmas to the defining sum for $l!\,\langle\omega,\eta\rangle_e$.
\end{proof}

For a finite-dimensional real inner product space $V$, $k\ge1$, and $\omega^i,\eta^j\in V^*$, write $a^\sharp$ for the Riesz representative and
\begin{equation}
\langle a,b\rangle=\langle a^\sharp,b^\sharp\rangle. \tag{2.26a}
\end{equation}

\begin{proposition}[The Pointwise Hodge Star]
\label{prop:hodge-star-pointwise}
\lean{LeeLib.Ch02.hodgeStarSum, LeeLib.Ch02.hodgeStarSum_add, LeeLib.Ch02.hodgeStarSum_smul, LeeLib.Ch02.wedgeSum_hodgeStarSum, LeeLib.Ch02.eq_hodgeStarSum_of_forall_wedgeSum_eq, LeeLib.Ch02.eq_hodgeStarSum_iff_forall_wedgeSum_eq, LeeLib.Ch02.hodgeStarSum_congr}
\uses{lem:hodge-scalar-identity, lem:top-form-eq-of-basis, lem:wedge-bilinear, lem:wedge-comm-sum, def:inner-k-covectors, lem:inner-k-covectors-posdef, prop:inner-k-covectors-frame-independent}
\leanok



For a bijection $\varepsilon:\iota_a\sqcup\iota_b\to\iota'$ and $\eta\in\Lambda^{\iota_a}(V^*)$, set
\[
{*}\eta=\frac1{l!}\sum_{t:\iota_b\to\iota'}B(\eta,E^t)E^t.
\]
The map is linear, satisfies $\omega\wedge {*}\eta=\langle\omega,\eta\rangle_eE^\varepsilon$, is uniquely characterized by this identity, and is independent of orthonormal block data with the same reference volume form.
\end{proposition}
\begin{proof}
\leanok
\sketch Linearity is immediate. Compare the two top forms on $e\circ\varepsilon$ and use Lemma \ref{lem:hodge-scalar-identity}. For uniqueness, the difference of two candidates wedges to zero with every form; the mirrored scalar identity and positive definiteness force it to vanish. Equal reference volumes and frame independence give the last assertion.
\end{proof}

\begin{proposition}[The Hodge Star in Degree Zero]
\label{prop:hodge-star-degree-zero}
\lean{LeeLib.Ch02.wedgeSum_isEmpty_left, LeeLib.Ch02.innerForms_isEmpty, LeeLib.Ch02.hodgeStarSum_isEmpty}
\uses{prop:hodge-star-pointwise, def:wedge-alternating-forms, def:inner-k-covectors}
\leanok



When $\iota_a=\varnothing$, $0$-covectors are scalars, the degree-zero inner product is multiplication, and $*f=f\,dV$.
\end{proposition}
\begin{proof}
\leanok
\sketch The empty-index permutation sum has one term, so wedging by a scalar is scalar multiplication and the pointwise star formula is the scalar multiple of the reference volume form.
\end{proof}

\begin{lemma}[The Reproducing Property of the Star]
\label{lem:hodge-star-reproducing}
\lean{LeeLib.Ch02.hodgeStarSum_apply_basis}
\uses{prop:hodge-star-pointwise, lem:k-covector-wedge-expansion}
\leanok



\[
({*}\eta)(e_{t(1)},\ldots,e_{t(l)})=B(\eta,E^t)\qquad(t:\iota_b\to\iota').
\]
\end{lemma}
\begin{proof}
\leanok
\sketch Evaluate the defining sum on $e\circ t$ and use the expansion Lemma \ref{lem:k-covector-wedge-expansion}; the inner sum is $l!E^t$.
\end{proof}

\begin{proposition}[The Hodge Star is a Fibrewise Isometry]
\label{prop:hodge-star-isometry}
\lean{LeeLib.Ch02.innerForms_hodgeStarSum}
\uses{prop:hodge-star-pointwise, lem:hodge-scalar-identity, lem:wedge-comm-sum, def:inner-k-covectors, lem:k-covector-wedge-expansion}
\leanok



\[
\langle {*}\omega,{*}\eta\rangle_e=\langle\omega,\eta\rangle_e.
\]
\end{proposition}
\begin{proof}
\leanok
\sketch Apply the scalar identity in the mirrored block structure and use the star characterization to identify reference values with coefficients of $\omega$ and $\eta$.
\end{proof}

\begin{proposition}[The Involution Property, Sign-Free]
\label{prop:hodge-star-involution}
\lean{LeeLib.Ch02.hodgeStarSum_hodgeStarSum}
\uses{prop:hodge-star-pointwise, lem:hodge-star-reproducing, lem:wedge-comm-sum, lem:k-covector-wedge-expansion, lem:top-form-eq-of-basis}
\leanok



For the mirrored block structure $\varepsilon'=\varepsilon\circ\mathrm{swap}$, $*'(*\eta)=\eta$.
\end{proposition}
\begin{proof}
\leanok
\sketch Verify the mirrored star identity on coframe wedges using Lemma \ref{lem:hodge-star-reproducing}, then extend by the wedge expansion and uniqueness.
\end{proof}

\begin{proposition}[The Hodge Star in the Classical Grading]
\label{prop:hodge-star-graded}
\lean{LeeLib.Ch02.hodgeStar, LeeLib.Ch02.hodgeStar_add, LeeLib.Ch02.hodgeStar_smul, LeeLib.Ch02.wedge_hodgeStar, LeeLib.Ch02.eq_hodgeStar_of_forall_wedge_eq, LeeLib.Ch02.eq_hodgeStar_iff_forall_wedge_eq, LeeLib.Ch02.innerForms_hodgeStar}
\uses{prop:hodge-star-pointwise, prop:hodge-star-isometry, def:wedge-alternating-forms}
\leanok



For an orthonormal basis and $k+l=n$, the classical star $*:\Lambda^k(V^*)\to\Lambda^{n-k}(V^*)$ satisfies
\[
\omega\wedge {*}\eta=\langle\omega,\eta\rangle_e(e^1\wedge\cdots\wedge e^n),
\]
and is unique and fibrewise isometric.
\end{proposition}
\begin{proof}
\leanok
\sketch Re-index the block construction along concatenation; the characterization, uniqueness, and isometry transfer from the pointwise star propositions.
\end{proof}

\begin{lemma}[Iterating the Cyclic Rotation]
\label{lem:finrotate-pow-val}
\lean{LeeLib.Ch02.finRotate_pow_val}
\leanok


If $\rho(i)=i+1\pmod n$ on $\{0,\ldots,n-1\}$, then $\rho^m(i)=i+m\pmod n$.
\end{lemma}
\begin{proof}
\leanok
\sketch Induct on $m$.
\end{proof}

\begin{lemma}[The Block Swap is a Cyclic Rotation]
\label{lem:blockswap-eq-rotate}
\lean{LeeLib.Ch02.blockswap_eq_rotate}
\uses{lem:finrotate-pow-val}
\leanok



For $n=k+l$, the block swap is cyclic rotation by $k$.
\end{lemma}
\begin{proof}
\leanok
\sketch The swap adds $k$ modulo $n$.
\end{proof}

\begin{proposition}[The Classical Hodge Involution]
\label{prop:hodge-star-classical-involution}
\lean{LeeLib.Ch02.hodgeStar_hodgeStar}
\uses{prop:hodge-star-graded, prop:hodge-star-involution, lem:wedge-covectors-sign, lem:hodge-star-congr-smul, lem:blockswap-eq-rotate}
\leanok



For an oriented $n$-dimensional inner product space and a $k$-covector,
\[
{*}{*}\omega=(-1)^{k(n-k)}\omega.
\]
\end{proposition}
\begin{proof}
\leanok
\sketch The second classical block order differs from the mirrored order by the block-swap rotation, whose sign is $(-1)^{k(n-k)}$. Combine this comparison with Proposition \ref{prop:hodge-star-involution}.
\end{proof}

\begin{proposition}\label{prop:hodge-star-operator}\dcref{ch2:2-18}


\uses{prop:hodge-star-pointwise,prop:hodge-star-graded,prop:hodge-star-degree-zero,prop:hodge-star-involution,prop:hodge-star-classical-involution,prop:fiber-metric-forms}
For an oriented Riemannian $n$-manifold, there is a unique smooth Hodge star
$*:\Lambda^kT^*M\to\Lambda^{n-k}T^*M$ satisfying
\[
\omega\wedge {*}\eta=\langle\omega,\eta\rangle_gdV_g.
\]
It satisfies $*f=fdV_g$ for $k=0$ and
$*^2\omega=(-1)^{k(n-k)}\omega$.
\end{proposition}

\begin{lemma}[A Wedge with a Repeated Covector Vanishes]
\label{lem:wedge-covectors-repeat-zero}
\lean{LeeLib.Ch02.wedgeCovectors_eq_zero_of_repeat}
\uses{def:wedge-covectors}
\leanok



If two covectors in a finite family coincide, their wedge is zero.
\end{lemma}
\begin{proof}
\leanok
\sketch The determinant defining the wedge has equal rows.
\end{proof}

\begin{proposition}[The Hodge Star of a Coordinate Covector]
\label{prop:hodge-star-euclidean-single}
\lean{LeeLib.Ch02.hodgeStar_flatL_single, LeeLib.Ch02.hodgeStar_basisFun_single}
\uses{prop:hodge-star-graded, prop:hodge-star-pointwise, lem:wedge-covectors-sign, lem:wedge-covectors-repeat-zero, def:inner-k-covectors, def:wedge-covectors}
\leanok


\dcref{ch2:2-19}

For an oriented orthonormal coframe $(e^0,\ldots,e^{n-1})$,
\[
{*} e^i=(-1)^i e^{i_1}\wedge\cdots\wedge e^{i_{n-1}},
\]
where the remaining indices are increasing; in $1$-based indexing the sign is $(-1)^{i-1}$.
\end{proposition}
\begin{proof}
\leanok
\uses{prop:hodge-star-graded,prop:hodge-star-pointwise,lem:wedge-covectors-sign,lem:wedge-covectors-repeat-zero,def:inner-k-covectors,def:wedge-covectors}
\sketch Verify the star characterization on coframe covectors. The cyclic reordering contributes $(-1)^i$ when the index agrees; otherwise the wedge repeats a covector and vanishes.
\end{proof}

\begin{lemma}[The Sign Law for Wedges of Covectors]
\label{lem:wedge-covectors-sign}
\lean{LeeLib.Ch02.wedgeCovectors_comp_perm}
\uses{def:wedge-covectors}
\leanok



Permuting a covector family multiplies its wedge by the permutation sign.
\end{lemma}
\begin{proof}
\leanok
\sketch The permutation permutes determinant rows.
\end{proof}

\begin{lemma}[Frame Independence up to a Scalar]
\label{lem:hodge-star-congr-smul}
\lean{LeeLib.Ch02.hodgeStarSum_congr_smul}
\uses{prop:hodge-star-pointwise, def:inner-k-covectors}
\leanok



If two reference volume forms satisfy $E^\varepsilon=sF^\zeta$, then $*_{(e,\varepsilon)}\eta=s\,*_{(f,\zeta)}\eta$.
\end{lemma}
\begin{proof}
\leanok
\sketch Compare the two star characterizations and use uniqueness.
\end{proof}

\begin{proposition}[The Hodge Star of an Ordered Basis Wedge]
\label{prop:hodge-star-basis-wedge}
\lean{LeeLib.Ch02.hodgeStar_wedgeCovectors_perm}
\uses{prop:hodge-star-pointwise, prop:hodge-star-graded, lem:wedge-covectors-sign, lem:wedge-covectors-repeat-zero, def:wedge-covectors, def:inner-k-covectors}
\leanok



For a permutation $\sigma$ and $k+l=n$, $*(e^{\sigma(0)}\wedge\cdots\wedge e^{\sigma(k-1)})=(\operatorname{sgn}\sigma)e^{\sigma(k)}\wedge\cdots\wedge e^{\sigma(n-1)}$.
\end{proposition}
\begin{proof}
\leanok
\uses{prop:hodge-star-pointwise,lem:wedge-covectors-sign,lem:wedge-covectors-repeat-zero,def:wedge-covectors,def:inner-k-covectors}
\sketch Verify on coframe wedges; repeated indices make both sides vanish.
\end{proof}

\begin{proposition}[The Hodge Star is an Involution on 2-Forms of a 4-Space]
\label{prop:hodge-star-two-involution}
\lean{LeeLib.Ch02.hodgeStar_hodgeStar_two}
\uses{prop:hodge-star-classical-involution}
\leanok



On an oriented $4$-space, $*^2$ is the identity on $2$-forms.
\end{proposition}
\begin{proof}
\leanok
\sketch Apply the classical involution with $n=k=2$.
\end{proof}

\begin{proposition}[The Hodge Star of a Coordinate 2-Covector on a 4-Space]
\label{prop:hodge-star-euclidean-pair}
\lean{LeeLib.Ch02.hodgeStar_wedgeCovectors_pair, LeeLib.Ch02.hodgeStar_basisFun_pair, LeeLib.Ch02.hodgeStar_e01, LeeLib.Ch02.hodgeStar_e02, LeeLib.Ch02.hodgeStar_e03, LeeLib.Ch02.hodgeStar_e23, LeeLib.Ch02.hodgeStar_e13, LeeLib.Ch02.hodgeStar_e12}
\uses{prop:hodge-star-basis-wedge, prop:hodge-star-two-involution}
\leanok


\dcref{ch2:2-19}

For an oriented orthonormal coframe in dimension $4$,
\[
{*}(e^0\wedge e^1)=e^2\wedge e^3,\quad {*}(e^0\wedge e^2)=-e^1\wedge e^3,\quad {*}(e^0\wedge e^3)=e^1\wedge e^2,
\]
\[
{*}(e^2\wedge e^3)=e^0\wedge e^1,\quad {*}(e^1\wedge e^3)=-e^0\wedge e^2,\quad {*}(e^1\wedge e^2)=e^0\wedge e^3.
\]
\end{proposition}
\begin{proof}
\leanok
\uses{prop:hodge-star-basis-wedge,prop:hodge-star-two-involution}
\sketch Use the ordered-basis formula for three values and $*^2=1$ for the rest.
\end{proof}

\begin{proposition}[Self-Dual / Anti-Self-Dual Decomposition]
\label{prop:self-dual-decomposition}
\lean{LeeLib.Ch02.selfDualForms, LeeLib.Ch02.antiSelfDualForms, LeeLib.Ch02.isCompl_selfDual_antiSelfDual, LeeLib.Ch02.existsUnique_selfDual_add_antiSelfDual}
\uses{prop:hodge-star-two-involution}
\leanok



The $+1$ and $-1$ eigenspaces of $*$ on $\Lambda^2(V^*)$ are complementary.
\end{proposition}
\begin{proof}
\leanok
\sketch Decompose $\omega=\frac12(\omega+*\omega)+\frac12(\omega-*\omega)$ and use $*^2=1$.
\end{proof}

\begin{proposition}[Self-Dual and Anti-Self-Dual 2-Forms on $\mathbb R^4$]
\label{prop:self-dual-euclidean}
\lean{LeeLib.Ch02.selfDual_add_e01_e23, LeeLib.Ch02.selfDual_sub_e02_e13, LeeLib.Ch02.selfDual_add_e03_e12, LeeLib.Ch02.antiSelfDual_sub_e01_e23, LeeLib.Ch02.antiSelfDual_add_e02_e13, LeeLib.Ch02.antiSelfDual_sub_e03_e12}
\uses{prop:hodge-star-euclidean-pair, prop:self-dual-decomposition}


\dcref{ch2:2-20}

The self-dual forms are $e^0\wedge e^1+e^2\wedge e^3$, $e^0\wedge e^2-e^1\wedge e^3$, and $e^0\wedge e^3+e^1\wedge e^2$; the anti-self-dual forms have the corresponding opposite signs.
\end{proposition}
\begin{proof}
\leanok
\uses{prop:hodge-star-euclidean-pair,prop:self-dual-decomposition}
\sketch Apply the six coordinate star identities termwise.
\end{proof}

\begin{proposition}\label{prop:local-isometry-distance}\dcref{ch2:2-31}


For connected Riemannian manifolds and a local isometry
$\varphi:M\to\widetilde M$,
$d_{\widetilde g}(\varphi(x),\varphi(y))\le d_g(x,y)$. Equality need not hold.
\end{proposition}

\begin{proposition}\label{prop:arc-length-integral}\dcref{ch2:2-32}
For a smooth curve $\gamma:[a,b]\to M$ and continuous $f$, define
\[
\int_\gamma f\,ds=\int_a^bf(t)|\gamma'(t)|_g\,dt.
\]
This integral is invariant under diffeomorphic reparametrization. For an
embedded curve it agrees with integration over the image using the induced
volume element.
\end{proposition}

\begin{exercise}\label{prob:arc-length-integral}\dcref{ch2:2-32}
Let $(M,g)$ be a Riemannian manifold and $\gamma: [a,b] \to M$ a smooth curve segment. For each continuous function $f: [a,b] \to \mathbb{R}$, we define the \textit{integral of $f$ with respect to arc length}, denoted by $\int_\gamma f \, ds$, by
\[
\int_\gamma f \, ds = \int_a^b f(t) |\gamma'(t)|_g \, dt.
\]

\begin{enumerate}
\item[(a)] Show that $\int_\gamma f \, ds$ is independent of parametrization in the following sense: if $\varphi: [c,d] \to [a,b]$ is a diffeomorphism, then
\[
\int_a^b f(t) |\gamma'(t)|_g \, dt = \int_c^d \tilde{f}(u) |\tilde{\gamma}'(u)|_g \, du,
\]
where $\tilde{f} = f \circ \varphi$ and $\tilde{\gamma} = \gamma \circ \varphi$.
\item[(b)] Suppose now that $\gamma$ is a smooth embedding, so that $C = \gamma([a,b])$ is an embedded submanifold with boundary in $M$. Show that
\[
\int_\gamma f \, ds = \int_C (f \circ \gamma^{-1}) \, d\tilde{V},
\]
where $d\tilde{V}$ is the Riemannian volume element on $C$ associated with the induced metric and the orientation determined by $\gamma$.
\end{enumerate}

(Used on p. 273.)
\end{exercise}


\begin{exercise}
\label{exer:characterize-local-isometry}
\dcref{ch2:2.7}
\lean{LeeLib.Ch02.isLocalIsometry_iff_isMetricPreserving}
\leanok
\uses{def:local-isometry,def:metric-preserving,def:isometry,lem:metric-preserving-equidim-local-diffeomorphism}
Prove that if $(M,g)$ and $(\widetilde{M}, \widetilde{g})$ are Riemannian manifolds of the same dimension, a smooth map $\varphi: M \to \widetilde{M}$ is a local isometry if and only if $\varphi^* \widetilde{g} = g$.
\end{exercise}
\begin{proof}
\leanok
\uses{def:local-isometry,lem:metric-preserving-equidim-local-diffeomorphism}
Suppose first that $\varphi$ is a local isometry, and let $p \in M$. Then $\varphi$ restricts, on some neighborhood $U$ of $p$, to an isometry onto an open subset of $\widetilde{M}$, so in particular $\widetilde{g}_{\varphi(p)}(d\varphi_p v, d\varphi_p w) = g_p(v,w)$ for all $v,w \in T_pM$. Since every point of $M$ lies in such a $U$, this identity holds at every point, which is $\varphi^*\widetilde{g} = g$. This half uses neither the smoothness hypothesis nor the equality of dimensions.

Conversely, suppose $\varphi^* \widetilde{g} = g$. By Lemma \ref{lem:metric-preserving-equidim-local-diffeomorphism}, $\varphi$ is a local diffeomorphism, so each $p \in M$ has a neighborhood $U$ on which $\varphi$ restricts to a diffeomorphism onto an open subset of $\widetilde{M}$. That restriction preserves the metric, because $\varphi$ does at every point of $M$ and metric preservation is a pointwise condition on the differential. Hence $\varphi|_U$ is an isometry onto an open subset, and $\varphi$ is a local isometry.
\end{proof}

\begin{exercise}
\label{exer:gram-schmidt-adapted-frame}
\dcref{ch2:2.15}
\uses{prop:adapted-orthonormal-frames,prop:slice-coordinates,prop:existence-orthonormal-frames}
Prove the preceding proposition. [Hint: Apply the Gram--Schmidt algorithm to a coordinate frame in slice coordinates (see Prop. A.22).]
\end{exercise}

\begin{exercise}
\label{exer:horizontal-vector-field-not-lift}
\lean{LeeLib.Ch02.horizontalNonLiftField_mem_horizontalSpace, LeeLib.Ch02.not_isHorizontalLift_horizontalNonLiftField}
\leanok
\dcref{ch2:2.26}
\uses{def:horizontal-space,def:horizontal-lift-tangent,ex:riemannian-submersions}
Let $\pi: \mathbb{R}^2 \to \mathbb{R}$ be the projection map $\pi(x,y) = x$, and let $\tilde{W}$ be the smooth vector field $y\partial_x$ on $\mathbb{R}^2$. Show that $\tilde{W}$ is horizontal, but there is no vector field on $\mathbb{R}$ whose horizontal lift is equal to $\tilde{W}$.
\end{exercise}

\begin{exercise}
\label{exer:inverse-metric-ambiguity}
\lean{LeeLib.Ch02.RiemannianMetric.innerDual_dualCoframe}
\leanok
\dcref{ch2:2.38}
\uses{def:dual-coframe,def:frame-matrix,def:sharp-map}
If $g$ is a Riemannian metric on $M$ and $(E_i)$ is a local frame on $M$, there is a potential ambiguity about what the expression $(g^{ij})$ represents: we have defined it to mean the inverse matrix of $(g_{ij})$, but one could also interpret it as the components of the contravariant 2-tensor field $g^{\#\#}$ obtained by raising both of the indices of $g$. Show that these two interpretations lead to the same result.
\end{exercise}

\begin{exercise}\label{exer:isometry-preserves-volume-form}\dcref{ch2:2.42}
\lean{LeeLib.Ch02.IsMetricPreserving.pullbackAlternating_volumeForm}
\leanok
\uses{def:pullback-alternating,def:volume-form-pointwise,prop:volume-form-uniqueness,prop:volume-form-value-one,lem:orthonormal-pushforward,lem:push-basis}
Suppose $(M,g)$ and $(\tilde{M},\tilde{g})$ are oriented Riemannian manifolds, and $\varphi: M \to \tilde{M}$ is an orientation-preserving isometry. Prove that $\varphi^* dV_{\tilde{g}} = dV_g$.
\end{exercise}
\begin{proof}
\leanok
\uses{def:pullback-alternating,prop:volume-form-uniqueness,prop:volume-form-value-one,lem:orthonormal-pushforward,lem:push-basis}
Fix $x \in M$ and let $(E_1,\ldots,E_n)$ be a local oriented orthonormal frame near $x$. By
Definition \ref{def:pullback-alternating},
\[
(\varphi^* dV_{\tilde{g}})|_x(E_1|_x,\ldots,E_n|_x)
  = dV_{\tilde{g}}|_{\varphi(x)}\bigl(d\varphi_x E_1|_x, \ldots, d\varphi_x E_n|_x\bigr).
\]
By Lemmas \ref{lem:orthonormal-pushforward} and \ref{lem:push-basis} the family
$\bigl(d\varphi_x E_i|_x\bigr)$ is an orthonormal basis of $T_{\varphi(x)}\tilde{M}$, and it is
positively oriented because $\varphi$ is orientation-preserving. So by Proposition
\ref{prop:volume-form-value-one} the right-hand side equals $1$.

Thus $(\varphi^* dV_{\tilde{g}})|_x$ is an $n$-form taking the value $1$ on the oriented
orthonormal frame $(E_i|_x)$. By Proposition \ref{prop:volume-form-uniqueness} the only such
$n$-form is $dV_g|_x$, so $(\varphi^* dV_{\tilde{g}})|_x = dV_g|_x$. As $x$ was arbitrary,
$\varphi^* dV_{\tilde{g}} = dV_g$.
\end{proof}

\begin{exercise}
\label{exer:orthonormal-frames-and-coframes}
\lean{LeeLib.Ch02.RiemannianMetric.tfae_orthonormal_frame_coframe}
\leanok
\dcref{ch2:2.39}
\uses{exer:inverse-metric-ambiguity,def:dual-coframe,def:sharp-map,def:frame-matrix}
Let $(M, g)$ be a Riemannian manifold with or without boundary, let $(E_i)$ be a local frame for $M$, and let $(e^i)$ be its dual coframe. Show that the following are equivalent:
\begin{enumerate}
\item[(a)] $(E_i)$ is orthonormal.
\item[(b)] $(e^i)$ is orthonormal.
\item[(c)] $(e^i)^{\sharp} = E_i$ for each $i$.
\end{enumerate}
\end{exercise}

\begin{exercise}
\label{exer:orthonormal-frames-proposition}
\dcref{ch2:2.10}
Use local orthonormal frames to prove the preceding proposition.
\end{exercise}

\begin{exercise}[Use a partition of unity to prove the preceding proposition]
\label{exer:partition-unity-proof}
\lean{LeeLib.Ch02.exists_riemannianMetric}
\leanok
\uses{prop:existence-riemannian-metric}
\dcref{ch2:2.5}
Use a partition of unity to prove the preceding proposition.
\end{exercise}

\begin{exercise}
\label{exer:polarization-proof}
\lean{LeeLib.Ch02.real_inner_eq_inner_add_add_self_sub_inner_sub_sub_self_div_four}
\leanok
\uses{lem:polarization-identity}
\dcref{ch2:2.2}
Prove the preceding lemma.
\end{exercise}

\begin{exercise}
\label{exer:polarization-scalar-product}
\lean{LeeLib.Ch02.BilinForm.IsSymm.inner_eq_apply_add_sub_apply_sub_div_four}
\leanok
\dcref{ch2:2.58}
\uses{exer:polarization-proof}
Prove that the polarization identity (2.2) holds for every scalar product.
\end{exercise}
\begin{proof}
\leanok
The proof of Lemma \ref{lem:polarization-identity} used only bilinearity and symmetry of $\langle \cdot, \cdot \rangle$, never positive definiteness: expanding $\langle v \pm w, v \pm w \rangle = \langle v,v \rangle \pm 2\langle v,w \rangle + \langle w,w \rangle$ and subtracting gives $4\langle v,w \rangle$. Since a scalar product is by definition a symmetric bilinear form, the identity holds verbatim.
\end{proof}

\begin{exercise}\label{exer:preceding-lemma-nondegenerate}
\leanok
\dcref{ch2:2.61}
\uses{lem:nondegenerate-subspace}
Prove the preceding lemma.
\end{exercise}
\begin{proof}
\leanok
This is the content of Lemma \ref{lem:nondegenerate-subspace}, proved above.
\end{proof}

\begin{exercise}[Induced Metric on $T^n$]
\label{exer:product-metric-torus}
\lean{LeeLib.Ch02.circleMetric, LeeLib.Ch02.torusMetric, LeeLib.Ch02.torusEmbedding_preservesMetric}
\leanok
\dcref{ch2:2.23}
\uses{def:product-metric,def:induced-metric}
Show that the induced metric on $T^n$ described in Exercise 2.21 is equal to the product metric obtained from the usual induced metric on $S^1 \subseteq \mathbb{R}^2$.

The formalization treats the case $n = 2$, where $T^2 = S^1 \times S^1$ needs no quotient-manifold construction: \lean{LeeLib.Ch02.circleMetric} is the round metric on $S^1 \subseteq \mathbb{C} = \mathbb{R}^2$ (the pullback of the Euclidean metric along the inclusion), \lean{LeeLib.Ch02.torusMetric} is the product metric $\mathtt{prodMetric}$ of two copies of it, and \lean{LeeLib.Ch02.torusEmbedding_preservesMetric} shows that the metric induced by the $L^2$-embedding $S^1 \times S^1 \hookrightarrow \mathbb{R}^4$ equals exactly that product metric --- so the induced flat metric on $T^2$ is the product of the induced circle metrics.
\end{exercise}

\begin{exercise}
\label{exer:product-pseudo-riemannian}
\dcref{ch2:2.68}
\uses{def:pseudo-riemannian-metric,def:pseudo-metric-has-signature,def:product-metric}

Suppose $(M_1, g_1)$ and $(M_2, g_2)$ are pseudo-Riemannian manifolds of signatures $(r_1, s_1)$ and $(r_2, s_2)$, respectively. Show that $(M_1 \times M_2, g_1 \oplus g_2)$ is a pseudo-Riemannian manifold of signature $(r_1 + r_2, s_1 + s_2)$.
\end{exercise}

\begin{exercise}
\label{exer:prove-isometry-invariance}
\leanok
\dcref{ch2:2.52}
\uses{prop:isometry-invariance-riemannian-distance}

Prove the preceding proposition.
\end{exercise}

\begin{exercise}
\label{exer:prove-lemma-2-11}
\lean{LeeLib.Ch02.pullbackForm_posDef_iff_immersion}
\leanok
\dcref{ch2:2.12}
\uses{lem:pullback-immersion-metric}
Prove Lemma 2.11.
\end{exercise}

\begin{exercise}
\label{exer:prove-nondegenerate-equivalent}
\leanok
\dcref{ch2:2.57}
\uses{lem:nondegenerate-symmetric-bilinear}
Prove the preceding lemma.
\end{exercise}
\begin{proof}
\leanok
This is the content of Lemma \ref{lem:nondegenerate-symmetric-bilinear}, proved above.
\end{proof}

\begin{exercise}
\label{exer:prove-orthonormal-frames}
\leanok
\dcref{ch2:2.67}
\uses{prop:orthonormal-frames-pseudo-riemannian}
Prove the preceding proposition.
\end{exercise}
\begin{proof}
\leanok
\uses{prop:orthonormal-frames-pseudo-riemannian}
This is the content of Proposition \ref{prop:orthonormal-frames-pseudo-riemannian}, proved above.
\end{proof}

\begin{exercise}
\label{exer:prove-outward-normal}
\dcref{ch2:2.18}
\uses{prop:induced-orientation-boundary}
Prove this proposition. [Hint: Use the paragraph preceding Prop. B.17 as a starting point.]
\end{exercise}

\begin{exercise}[Prove Proposition 2.47]
\label{exer:prove-prop-247}
\leanok
\dcref{ch2:2.48}
\uses{prop:length-additive,prop:length-reparam-invariant,prop:length-isometry-invariant}
\end{exercise}

\begin{exercise}\label{exer:riemannian-density-formula}\dcref{ch2:2.45}
\lean{LeeLib.Ch02.densityL_eq_abs_wedgeCovectors}
\leanok
\uses{prop:riemannian-density,def:density-vector-space}
Prove this proposition by showing that $\mu$ can be defined in terms of any local orthonormal frame by
\[
\mu = |\varepsilon^1 \wedge \cdots \wedge \varepsilon^n|.
\]
The formalization is the pointwise (fibrewise) identity $\mu_x = |\varepsilon^1 \wedge \cdots \wedge \varepsilon^n|$ on each tangent space $T_xM$, from which Lee's statement is the specialization $V := T_xM$; both sides equal $|\det[\langle E_i, v_j\rangle]|$.
\end{exercise}

\begin{exercise}
\label{exer:torus-covering}
\dcref{ch2:2.36}
\uses{ex:n-torus-riemannian-submanifold}
Let $T^n \subseteq \mathbb{R}^{2n}$ be the $n$-torus with its induced metric. Show that the map $X: \mathbb{R}^n \to T^n$ of Example 2.21 is a Riemannian covering.
\end{exercise}

\begin{exercise}
\label{exer:verify-surface-revolution-examples}
\dcref{ch2:2.22}
\uses{ex:metrics-in-graph-coordinates,ex:surfaces-of-revolution,ex:n-torus-riemannian-submanifold}

Verify the claims in Examples 2.19–2.21.
\end{exercise}

\begin{exercise}
\label{prob:codifferential-operator}
\dcref{ch2:2-29}
Let $(M,g)$ be a compact oriented Riemannian $n$-manifold. For $1 \leq k \leq n$, define a map $d^* : \Omega^k(M) \to \Omega^{k-1}(M)$ by $d^* \omega = (-1)^{n(k+1)+1} * d * \omega$, where $*$ is the Hodge star operator defined in Problem 2-18. Extend this definition to 0-forms by defining $d^*\omega = 0$ for $\omega \in \Omega^0(M)$.

\begin{enumerate}
\item[(a)] Show that $d^* \circ d^* = 0$.
\item[(b)] Show that the formula
\[
(\omega, \eta) = \int_M \langle \omega, \eta \rangle_g \, dV_g
\]
defines an inner product on $\Omega^k(M)$ for each $k$, where $(\cdot, \cdot)_g$ is the inner product on forms defined in Problem 2-16.
\item[(c)] Show that $(d^*\omega, \eta) = (\omega, d\eta)$ for all $\omega \in \Omega^k(M)$ and $\eta \in \Omega^{k-1}(M)$.
\end{enumerate}
\end{exercise}

\begin{exercise}
\label{prob:covering-volume-form}
\dcref{ch2:2-14}
Suppose $(\tilde{M}, \tilde{g})$ and $(M, g)$ are compact connected Riemannian manifolds, and $\pi : \tilde{M} \to M$ is a $k$-sheeted Riemannian covering. Prove that $\text{Vol}(\tilde{M}) = k \cdot \text{Vol}(M)$. (Used on p. 363.)
\end{exercise}

\begin{exercise}
\label{prob:curl-vector-field}
\dcref{ch2:2-27}
Suppose $(M,g)$ is an oriented Riemannian 3-manifold.

(a) Define $\beta : TM \to \Lambda^2 T^* M$ by $\beta(X) = X \mathbin{\lrcorner} dV_g$. Show that $\beta$ is a smooth bundle isomorphism, and thus we can define the \emph{curl} of a vector field $X \in \mathfrak{X}(M)$ by
\[
\operatorname{curl} X = \beta^{-1} d(X^b).
\]

(b) Show that the following diagram commutes:
\[
\begin{array}{ccccccc}
C^{\infty}(M) & \xrightarrow{\operatorname{grad}} & \mathfrak{X}(M) & \xrightarrow{\operatorname{curl}} & \mathfrak{X}(M) & \xrightarrow{\operatorname{div}} & C^{\infty}(M) \\
\downarrow \text{Id} & & \downarrow b & & \downarrow \beta & & \downarrow * \\
\Omega^0(M) & \xrightarrow{d} & \Omega^1(M) & \xrightarrow{d} & \Omega^2(M) & \xrightarrow{d} & \Omega^3(M)
\end{array}
\]
where $*(f) = f \, dV_g$, and use this to prove that $\operatorname{curl}(\operatorname{grad} f) = 0$ for every $f \in C^{\infty}(M)$, and $\operatorname{div}(\operatorname{curl} X) = 0$ for every $X \in \mathfrak{X}(M)$.

(c) Compute the formula for the curl in standard coordinates on $\mathbb{R}^3$ with the Euclidean metric.
\end{exercise}

\begin{exercise}\label{prob:dirichlet-neumann-eigenvalues}\dcref{ch2:2-25}
Let $(M,g)$ be a compact connected Riemannian $n$-manifold with nonempty
boundary. A number $\lambda \in \mathbb{R}$ is called a \emph{Dirichlet eigenvalue for $M$} if there
exists a smooth real-valued function $u$ on $M$, not identically zero, such that
$-\Delta u = \lambda u$ and $u|_{\partial M} = 0$. Similarly, $\lambda$ is called a \emph{Neumann eigenvalue} if there exists a smooth real-valued function $u$ on $M$, not identically zero, such that $-\Delta u = \lambda u$ and $Nu|_{\partial M} = 0$, where $N$ is the outward unit normal vector field along $\partial M$.

(a) Show that every Dirichlet eigenvalue is strictly positive.
(b) Show that 0 is a Neumann eigenvalue, and all other Neumann eigenvalues are strictly positive.
\end{exercise}

\begin{exercise}[Dirichlet's Principle]
\label{prob:dirichlet-principle}
\dcref{ch2:2-26}
Suppose $(M,g)$ is a compact connected Riemannian $n$-manifold with nonempty boundary. Prove that a function $u \in C^{\infty}(M)$ is harmonic if and only if it minimizes $\int_M |\operatorname{grad} u|^2 \, dV_g$ among all smooth functions with the same boundary values. [Hint: For any function $f \in C^{\infty}(M)$ that vanishes on $\partial M$, expand $\int_M |\operatorname{grad}(u + \varepsilon f)|^2 \, dV_g$ and use Problem 2-22.]
\uses{prob:divergence-theorem}
\end{exercise}

\begin{exercise}\label{prob:distance-disconnected-riemannian}\dcref{ch2:2-30}
Suppose $(M,g)$ is a (not necessarily connected) Riemannian manifold. Show that there is a distance function $d$ on $M$ that induces the given topology and restricts to the Riemannian distance on each component of $M$. (Used on p. 187.)
\end{exercise}

\begin{exercise}\label{prob:distance-preserving-map-linear-isometry}\dcref{ch2:2-2}
\lean{LeeLib.Ch01.exists_linearIsometryEquiv_of_dist_eq}
\leanok
Suppose $V$ and $W$ are finite-dimensional real inner product spaces of the same dimension, and $F: V \to W$ is any map (not assumed to be linear or even continuous) that preserves the origin and all distances: $F(0) = 0$ and $|F(x) - F(y)| = |x - y|$ for all $x, y \in V$. Prove that $F$ is a linear isometry. [Hint: First show that $F$ preserves inner products, and then show that it is linear.] (Used on p. 187.)
\end{exercise}

\begin{exercise}\label{prob:divergence-laplacian-formulas-ch2}\dcref{ch2:2-21}
Prove Proposition~2.46 (the coordinate formulas for the divergence and the Laplacian).
\uses{prop:coordinate-representations-divergence-laplacian}
\end{exercise}

\begin{exercise}\label{prob:divergence-theorem}\dcref{ch2:2-22}
Suppose $(M,g)$ is a compact Riemannian manifold with boundary.

(a) Prove the following \textbf{divergence theorem} for $X \in \mathfrak{X}(M)$:
\[
\int_M (\operatorname{div} X) dV_g = \int_{\partial M} \langle X, N\rangle_g d\bar{g},
\]
where $N$ is the outward unit normal to $\partial M$ and $\bar{g}$ is the induced metric on $\partial M$. [Hint: Prove it first in the case that $M$ is orientable, and then apply that case to the orientation covering of $M$ (Prop.~B.18).]
\uses{prop:orientation-covering-theorem}

(b) Show that the divergence operator satisfies the following product rule for
$u \in C^\infty(M)$ and $X \in \mathfrak{X}(M)$:
\[
\operatorname{div}(uX) = u \operatorname{div} X + \langle \operatorname{grad} u, X \rangle_g,
\]
and deduce the following ``integration by parts'' formula:
\[
\int_M \langle \operatorname{grad} u, X \rangle_g \, dV_g = \int_{\partial M} u(X,N)_g \, d\tilde{V}_g - \int_M u \operatorname{div} X \, dV_g.
\]
What does this say when $M$ is a compact interval in $\mathbb{R}$?

(Used on p.\ 149.)
\end{exercise}

\begin{exercise}\label{prob:eigenvalues}\dcref{ch2:2-24}
Let $(M,g)$ be a compact Riemannian manifold (without boundary). A real
number $\lambda$ is called an \emph{eigenvalue of $M$} if there exists a smooth function $u$
on $M$, not identically zero, such that $-\Delta u = \lambda u$. In this case, $u$ is called an
\emph{eigenfunction} corresponding to $\lambda$.

(a) Prove that $0$ is an eigenvalue of $M$, and that all other eigenvalues are
strictly positive.

(b) Show that if $u$ and $v$ are eigenfunctions corresponding to distinct eigen-
values, then $\int_M uv \, dV_g = 0$.

(Used on p.\ 149.)
\end{exercise}

\begin{exercise}
\label{prob:fiber-metric-exterior-powers}
\lean{LeeLib.Ch02.RiemannianMetric.riemannian_fiberMetric_forms, LeeLib.Ch02.RiemannianMetric.innerFormsAt_self_pos, LeeLib.Ch02.RiemannianMetric.innerFormsAt_comm, LeeLib.Ch02.RiemannianMetric.contMDiff_innerFormsAt}
\leanok
\uses{prop:fiber-metric-forms,prop:inner-k-covectors-unique,prop:inner-k-covectors-frame-independent,lem:inner-k-covectors-posdef}
\dcref{ch2:2-16}
Let $(M, g)$ be a Riemannian $n$-manifold. Show that for each $k = 1, \ldots, n$, there is a unique fiber metric $\langle \cdot, \cdot \rangle_g$ on the bundle $\Lambda^k T^* M$ that satisfies
\[
\langle \omega^1 \wedge \cdots \wedge \omega^k, \eta^1 \wedge \cdots \wedge \eta^k \rangle_g = \det\left(\langle \omega^i, \eta^j \rangle_g\right) \quad (2.26)
\]
whenever $\omega^1, \ldots, \omega^k, \eta^1, \ldots, \eta^k$ are covectors at a point $p \in M$. [Hint: Define the inner product locally by declaring the set of $k$-covectors
\[
\{e^{i_1} \wedge \cdots \wedge e^{i_k} |_p : i_1 < \cdots < i_k\}
\]
to be an orthonormal basis for $\Lambda^k(T_p^* M)$ whenever $(e^i)$ is a local orthonormal coframe for $T^* M$, and then prove that the resulting inner product satisfies (2.26) and is independent of the choice of frame.]
\end{exercise}

\begin{exercise}\label{prob:gradient-characterization}
\dcref{ch2:2-10}
\lean{LeeLib.Ch02.eq_grad_iff_of_forall_ne_zero}
\leanok
\uses{def:gradient,def:regular-set,lem:gradient-characterization,lem:transversal-representative-eq-gradient}
Suppose $(M,g)$ is a Riemannian manifold, $f \in C^\infty(M)$, and $X \in \mathfrak{X}(M)$ is a nowhere-vanishing vector field. Prove that $X = \operatorname{grad} f$ if and only if $Xf = |X|_g^2$ and $X$ is orthogonal to the level sets of $f$ at all regular points of $f$. (Used on pp.~161, 180.)
\end{exercise}
\begin{proof}
\leanok
Throughout, the tangent space to the level set of $f$ through a regular point $p$ is $\ker df_p$, so ``$X$ is orthogonal to the level sets at regular points'' means $\langle X_p, w\rangle_g = 0$ for every regular $p$ and every $w \in \ker df_p$; and $Xf = df(X)$, $|X|_g^2 = \langle X, X\rangle_g$.

Suppose first that $X = \operatorname{grad} f$. Then $Xf = df(X) = \langle \operatorname{grad} f, X\rangle_g = \langle X, X\rangle_g = |X|_g^2$ by Lemma \ref{lem:gradient-characterization}, and for $w \in \ker df_p$ the same lemma gives $\langle X_p, w\rangle_g = df_p(w) = 0$. Neither implication used the hypothesis that $X$ vanishes nowhere.

Conversely, suppose $Xf = |X|_g^2$ and $X$ is orthogonal to the level sets at regular points. Fix $p \in M$. Since $X_p \neq 0$ and $g$ is positive definite, $\langle X_p, X_p\rangle_g > 0$, so
\[
df_p(X_p) = \langle X_p, X_p \rangle_g > 0 .
\]
In particular $df_p \neq 0$, so \emph{every} point of $M$ is a regular point of $f$ and the orthogonality hypothesis is available at $p$. The claim $X_p = \operatorname{grad} f|_p$ is then Lemma \ref{lem:transversal-representative-eq-gradient}.
\end{proof}

\begin{exercise}\label{prob:gradient-orthogonal-level-sets}
\lean{LeeLib.Ch02.innerAt_grad_mfderiv_regularLevelSetIncl_eq_zero}
\leanok
\dcref{ch2:2-9}
\uses{prop:regular-level-set-hypersurface,prop:gradient-normal-regular-part,lem:gradient-nonvanishing-regular-part}
Prove Proposition 2.37 (the gradient is orthogonal to regular level sets).
\end{exercise}

\begin{exercise}\label{prob:harmonic-functions}\dcref{ch2:2-23}
Let $(M,g)$ be a compact Riemannian manifold with or without boundary.
A function $u \in C^\infty(M)$ is said to be \emph{harmonic} if $\Delta u = 0$, where $\Delta$ is the
Laplacian defined on page~32.

(a) Prove \emph{Green's identities}:
\[
\int_M u\Delta v \, dV_g = \int_{\partial M} u N v \, d\tilde{V}_g - \int_M \langle \operatorname{grad} u, \operatorname{grad} v \rangle_g \, dV_g,
\]
\[
\int_M (u\Delta v - v\Delta u) \, dV_g = \int_{\partial M} (u N v - v N u) \, d\tilde{V}_g,
\]
where $N$ is the outward unit normal vector field on $\partial M$ and $\tilde{g}$ is the
induced metric on $\partial M$.

(b) Show that if $M$ is connected, $\partial M \ne \emptyset$, and $u, v$ are harmonic functions
on $M$ whose restrictions to $\partial M$ agree, then $u \equiv v$.

(c) Show that if $M$ is connected and $\partial M = \emptyset$, then the only harmonic
functions on $M$ are the constants, and every smooth function $u$ satisfies
$\int_M \Delta u \, dV_g = 0$.

(Used on pp.\ 149, 223.)
\end{exercise}

\begin{exercise}\label{prob:hodge-star-euclidean}\dcref{ch2:2-19}
\lean{LeeLib.Ch02.hodgeStar_basisFun_single, LeeLib.Ch02.hodgeStar_basisFun_pair}
\leanok
\uses{prop:hodge-star-euclidean-single,prop:hodge-star-euclidean-pair}
Regard $\mathbb{R}^n$ as a Riemannian manifold with the Euclidean metric and the standard orientation, and let $*$ denote the Hodge star operator defined in Problem~2-18.

(a) Calculate $* dx^i$ for $i = 1, \ldots, n$.

(b) Calculate $*(dx^i \wedge dx^j)$ in the case $n = 4$.
\end{exercise}

\begin{exercise}
\label{prob:hodge-star-exterior-derivative}
\dcref{ch2:2-28}
Let $(M,g)$ be an oriented Riemannian manifold and let $*$ denote its Hodge star operator (Problem 2-18). Show that for every $X \in \mathfrak{X}(M)$,
\[
X \mathbin{\lrcorner} dV_g = *(X^b),
\]
\[
\operatorname{div} X = *d*(X^b),
\]
and, when $\dim M = 3$,
\[
\operatorname{curl} X = (*d(X^b))^{\#},
\]
where the curl of a vector field is defined as in Problem 2-27.
\uses{prob:hodge-star-operator,prob:curl-vector-field}
\end{exercise}

\begin{exercise}\label{prob:hodge-star-operator}\dcref{ch2:2-18}
\lean{LeeLib.Ch02.RiemannianMetric.riemannian_hodgeStar, LeeLib.Ch02.RiemannianMetric.contMDiff_hodgeStarAt, LeeLib.Ch02.RiemannianMetric.hodgeStarAt_add, LeeLib.Ch02.RiemannianMetric.hodgeStarAt_smul}
\leanok
\uses{prop:hodge-star-pointwise,prop:hodge-star-graded,prop:hodge-star-degree-zero,prop:hodge-star-involution,prop:hodge-star-classical-involution,prob:fiber-metric-exterior-powers}
Let $(M,g)$ be an oriented Riemannian $n$-manifold.

(a) For each $k = 0, \ldots, n$, show that there is a unique smooth bundle homomorphism $*: \Lambda^k T^* M \to \Lambda^{n-k}T^*M$, called the \textbf{Hodge star operator}, satisfying
\[
\omega \wedge *\eta = \langle\omega,\eta\rangle_g dV_g
\]
for all smooth $k$-forms $\omega, \eta$, where $\langle\cdot,\cdot\rangle_g$ is the inner product on $k$-forms defined in Problem~2-16. (For $k = 0$, interpret the inner product as ordinary multiplication.) [Hint: First prove uniqueness, and then define $*$ locally by setting
\[
*(e^{i_1} \wedge \cdots \wedge e^{i_k}) = \pm e^{j_1} \wedge \cdots \wedge e^{j_{n-k}}
\]
in terms of an orthonormal coframe $(e^i)$, where the indices $j_1,\ldots,j_{n-k}$ are chosen such that $(i_1,\ldots,i_k,j_1,\ldots,j_{n-k})$ is some permutation of $(1,\ldots,n)$.]

(b) Show that $*: \Lambda^0 T^*M \to \Lambda^n T^*M$ is given by $*f = f dV_g$.

(c) Show that $* * \omega = (-1)^{k(n-k)}\omega$ if $\omega$ is a $k$-form.
\end{exercise}

\begin{exercise}\label{prob:horizontal-vector-fields}
\dcref{ch2:2-5}
\lean{LeeLib.Ch02.exists_horizontalLift_eq}
\leanok
\uses{prop:horizontal-lift-of-vector-field,prop:horizontal-lift-realizes-horizontal-vector}
Prove parts (b) and (c) of Proposition 2.25 (properties of horizontal vector fields). (Used on p.~146.)
\end{exercise}

\begin{exercise}
\label{prob:hypersurface-volume-form}
\lean{LeeLib.Ch02.exists_isSmoothOrientation_iff_exists_isSmoothNormalField}
\leanok
\dcref{ch2:2-13}
\uses{prop:hypersurface-volume-form,prop:hypersurface-volume-form-formula,prop:hypersurface-volume-form-if,prop:hypersurface-volume-form-only-if}
Prove Proposition 2.43 (characterizing the volume form of a Riemannian hypersurface). [Hint: To prove (2.17), use an adapted orthonormal frame.]
\end{exercise}

\begin{exercise}
\label{prob:inner-products-k-forms}
\dcref{ch2:2-17}
\uses{prob:fiber-metric-exterior-powers,prop:inner-products-of-tensors}
Because we regard the bundle $\Lambda^k T^* M$ of $k$-forms as a subbundle of the bundle $T^k T^* M$ of covariant $k$-tensors, we have two inner products to choose from on $k$-forms: the one defined in Problem 2-16, and the restriction of the tensor inner product defined in Proposition 2.40. For this problem, we use the notation $\langle \cdot, \cdot \rangle$ to denote the inner product of Problem 2-16, and $\langle\langle \cdot, \cdot \rangle\rangle$ to denote the restriction of the tensor inner product.

(a) Using the convention for the wedge product that we use in this book (see p. 400), prove that
\[
\langle \cdot, \cdot \rangle_g = \frac{1}{k!} \langle\langle \cdot, \cdot \rangle\rangle_g.
\]

(b) Show that if the Alt convention is used for the wedge product (p. 401), the formula becomes
\[
\langle\langle \cdot, \cdot \rangle\rangle_g = k! \langle \cdot, \cdot \rangle_g.
\]
\end{exercise}

\begin{exercise}\label{prob:inner-products-tensor-bundles}
\dcref{ch2:2-11}
\uses{prop:inner-products-of-tensors}
Prove Proposition 2.40 (inner products on tensor bundles).
\end{exercise}

\begin{exercise}\label{prob:level-sets-pseudo-riemannian}\dcref{ch2:2-35}
\lean{LeeLib.Ch02.hasSignature_pseudoPullbackMetric_levelSet_of_pseudoGrad_pos, LeeLib.Ch02.hasSignature_pseudoPullbackMetric_levelSet_of_pseudoGrad_neg, LeeLib.Ch02.form_pseudoGrad_mfderiv_levelSet_val_eq_zero}
\leanok
\uses{cor:pseudo-riemannian-embedding-criterion}
Prove Corollary 2.71 (about level sets in pseudo-Riemannian manifolds). (Used on p.\ 63.)
\end{exercise}

\begin{exercise}\label{prob:local-isometry-distance}\dcref{ch2:2-31}
\lean{LeeLib.Ch02.IsLocalIsometry.riemannianEDist_le}
\leanok
Suppose $(M,g)$ and $(\tilde{M}, \tilde{g})$ are connected Riemannian manifolds, and $\varphi: M \to \tilde{M}$ is a local isometry. Show that $d_{\tilde{g}}(\varphi(x), \varphi(y)) \leq d_g(x,y)$ for all $x, y \in M$. Give an example to show that equality need not hold. (Used on p. 37.)
\end{exercise}

\begin{exercise}\label{prob:lorentz-metric-existence}\dcref{ch2:2-34}
\lean{LeeLib.Ch02.exists_isLorentzMetric_iff_exists_isRankOneDistribution}
\leanok
\uses{thm:lorentz-metrics}
Prove Theorem 2.69 (existence of Lorentz metrics), as follows.

\begin{enumerate}
\item[(a)] For sufficiency, assume that $D \subseteq TM$ is a 1-dimensional distribution, and choose any Riemannian metric $g$ on $M$. Show that locally it is possible to choose a $g$-orthonormal frame $(E_i)$ and dual coframe $(\varepsilon^i)$ such that $E_1$ spans $D$; and then show that the Lorentz metric $-(\varepsilon^1)^2 + (\varepsilon^2)^2 + \cdots + (\varepsilon^n)^2$ is independent of the choice of frame.

\item[(b)] To prove necessity, suppose that $g$ is a Lorentz metric on $M$, and let $g_0$ be any Riemannian metric. Show that for each $p \in M$, there are exactly two $g_0$-unit vectors $v_0, -v_0$ on which the function $v \mapsto g(v, v)$ takes its minimum among all unit vectors in $T_p M$, and use Lagrange multipliers to conclude that there exists a number $\lambda(p) < 0$ such that $g(v_0, w) = \lambda(p)g_0(v_0, w)$ for all $w \in T_p M$. You may use the following standard result from perturbation theory: if $U$ is an open subset of $\mathbb{R}^n$ and $A: U \to GL(n, \mathbb{R})$ is a smooth matrix-valued function such that $A(x)$ is symmetric and has exactly one negative eigenvalue for each $x \in U$, then there exist smooth functions $\lambda: U \to (-\infty, 0)$ and $X: U \to \mathbb{R}^n \setminus \{0\}$ such that $A(x)X(x) = \lambda(x)X(x)$ for all $x \in U$.
\end{enumerate}
\end{exercise}

\begin{exercise}\label{prob:mobius-band-metric}
\dcref{ch2:2-8}
\uses{ex:open-mobius-band}
Prove that the action of $\mathbb{Z}$ on $\mathbb{R}^2$ defined in Example 2.35 is smooth, free, proper, and isometric, and therefore the open Möbius band inherits a flat Riemannian metric such that the quotient map is a Riemannian covering.
\end{exercise}

\begin{exercise}\label{prob:riemannian-1-manifold-flat}\dcref{ch2:2-1}
Show that every Riemannian $1$-manifold is flat. (Used on pp. 13, 193.)
\end{exercise}

\begin{exercise}\label{prob:riemannian-metric-submersion}
\lean{LeeLib.Ch02.existsUnique_isRiemannianSubmersion_metric}
\leanok
\dcref{ch2:2-6}
\uses{thm:group-action-riemannian-submersion}
Prove Theorem 2.28 (if $\pi: \tilde{M} \to M$ is a surjective smooth submersion, and a group acts on $\tilde{M}$ isometrically, vertically, and transitively on fibers, then $M$ inherits a unique Riemannian metric such that $\pi$ is a Riemannian submersion).
\end{exercise}

\begin{exercise}
\label{prob:riemannian-volume-form-existence}
\lean{LeeLib.Ch02.RiemannianMetric.riemannian_volumeForm}
\leanok
\dcref{ch2:2-12}
\uses{prop:riemannian-volume-form}
Prove Proposition 2.41 (existence and uniqueness of the Riemannian volume form).
\end{exercise}

\begin{exercise}\label{prob:self-dual-anti-self-dual}\dcref{ch2:2-20}
\lean{LeeLib.Ch02.existsUnique_selfDual_add_antiSelfDual, LeeLib.Ch02.selfDual_add_e01_e23, LeeLib.Ch02.selfDual_sub_e02_e13, LeeLib.Ch02.selfDual_add_e03_e12, LeeLib.Ch02.antiSelfDual_sub_e01_e23, LeeLib.Ch02.antiSelfDual_add_e02_e13, LeeLib.Ch02.antiSelfDual_sub_e03_e12}
\leanok
\uses{prop:self-dual-decomposition,prop:self-dual-euclidean}
Let $M$ be an oriented Riemannian 4-manifold. A 2-form $\omega$ on $M$ is said to be \textbf{self-dual} if $*\omega = \omega$, and \textbf{anti-self-dual} if $*\omega = -\omega$.

(a) Show that every 2-form $\omega$ on $M$ can be written uniquely as a sum of a self-dual form and an anti-self-dual form.

(b) On $M = \mathbb{R}^4$ with the Euclidean metric, determine the self-dual and anti-self-dual 2-forms in standard coordinates.
\end{exercise}

\begin{exercise}\label{prob:surface-of-revolution-warped-product}\dcref{ch2:2-3}
\uses{ex:warped-products}
Given a smooth embedded $1$-dimensional submanifold $C \subseteq H$ as in Example 2.24(b), show that the surface of revolution $S_C \subseteq \mathbb{R}^3$ with its induced metric is isometric to the warped product $C \times_a S^1$, where $a: C \to \mathbb{R}$ is the distance to the $z$-axis.
\end{exercise}

\begin{exercise}\label{prob:sylvesters-law-inertia}\dcref{ch2:2-33}
\leanok
\uses{prop:sylvesters-law-inertia}

Prove Proposition 2.65 (Sylvester's law of inertia).
\end{exercise}

\begin{exercise}\label{prob:warped-product-isometry}
\dcref{ch2:2-4}
\lean{LeeLib.Ch02.isMetricPreserving_polarMap}
\leanok
\uses{ex:warped-products,def:polar-map,lem:mfderiv-polar-map,lem:sphere-tangent-orthogonal-radius,def:metric-preserving,def:open-submanifold-metric,prop:warped-product-constant-one}
Let $\rho: \mathbb{R}^+ \to \mathbb{R}$ be the restriction of the standard coordinate function, and let $\mathbb{R}^+ \times_\rho S^{n-1}$ denote the resulting warped product (see Example 2.24(c)). Define $\Phi: \mathbb{R}^+ \times_\rho S^{n-1} \to \mathbb{R}^n \setminus \{0\}$ by $\Phi(\rho,\omega) = \rho\omega$. Show that $\Phi$ is an isometry between the warped product metric and the Euclidean metric on $\mathbb{R}^n \setminus \{0\}$. (Used on p.~293.)
\end{exercise}
\begin{proof}
\leanok
\uses{def:polar-map,lem:mfderiv-polar-map,lem:sphere-tangent-orthogonal-radius,def:metric-preserving}
We verify $\Phi^*\bar{g} = g_{\mathbb{R}^+} \oplus \rho^2 \mathring{g}$ at each point; that $\Phi$ is a bijection onto $\mathbb{R}^n \setminus \{0\}$ is checked separately below.

Fix $(r,\omega)$ and let $(a,w), (a',w')$ be tangent vectors there. By Lemma \ref{lem:mfderiv-polar-map},
\[
\langle d\Phi(a,w), d\Phi(a',w')\rangle
 = \langle a\omega + r\,d\iota(w),\ a'\omega + r\,d\iota(w')\rangle,
\]
which expands into four terms:
\[
a a' \langle\omega,\omega\rangle + a r \langle \omega, d\iota(w')\rangle + a' r \langle d\iota(w), \omega\rangle + r^2 \langle d\iota(w), d\iota(w')\rangle.
\]
Now $\langle\omega,\omega\rangle = |\omega|^2 = 1$ because $\omega \in S^{n-1}$, and the two middle terms vanish by Lemma \ref{lem:sphere-tangent-orthogonal-radius}. Since $\mathring{g}_\omega(w,w') = \langle d\iota(w), d\iota(w')\rangle$ by the definition of the round metric as $\iota^*\bar{g}$, what remains is
\[
a a' + r^2\,\mathring{g}_\omega(w,w'),
\]
which is exactly the warped product metric $g_{\mathbb{R}^+} \oplus \rho^2\mathring{g}$ evaluated at $(r,\omega)$, the radial factor $g_{\mathbb{R}^+}$ being $\iota^*\bar{g} = d\rho^2$ by Lemma \ref{lem:mfderiv-open-inclusion} and the warping function being $\rho(r) = r$.

For the bijection: $|\Phi(r,\omega)| = r|\omega| = r > 0$, so $\Phi$ never takes the value $0$; if $\Phi(r,\omega) = \Phi(r',\omega')$ then taking norms gives $r = r'$ and then cancelling the nonzero scalar $r$ gives $\omega = \omega'$, so $\Phi$ is injective; and any $x \neq 0$ equals $\Phi(|x|, x/|x|)$, so $\Phi$ is onto $\mathbb{R}^n \setminus \{0\}$.
\end{proof}

\begin{exercise}
\label{prob:warped-product-volume-form}
\dcref{ch2:2-15}
Suppose $(M_1, g_1)$ and $(M_2, g_2)$ are oriented Riemannian manifolds of dimensions $k_1$ and $k_2$, respectively. Let $f : M_1 \to \mathbb{R}^+$ be a smooth function, and let $g = g_1 \oplus f^2 g_2$ be the corresponding warped product metric on $M_1 \times_f M_2$. Prove that the Riemannian volume form of $g$ is given by
\[
dV_g = f^{k_2} dV_{g_1} \wedge dV_{g_2},
\]
where $f$, $dV_{g_1}$, and $dV_{g_2}$ are understood to be pulled back to $M_1 \times M_2$ by the projection maps. (Used on p. 295.)
\end{exercise}
